\chapter{Discharge of a witness and restricted equality rules}
\label{prov:ch:disch}
\label{disch:chap}

A derivation in the native scalar core $\Nat$ (\cref{base:sec:kernel})
supplies data at three places. The rule \rulename{Tr} takes an equation
$s=t$ together with a term template. The rule \rulename{DE} opens a paired
formula at a fresh witness and is separate from every equality rule. The
rule \rulename{QE}, and the use of a clause that replaces it in a clause
calculus (\cref{base:sec:clauses}), takes an arbitrary term. The calculi of
this chapter fuse or restrict these interfaces.

In \cref{disch:sec:name} one rule, \rulename{Name}, discharges an opened
witness $b$ together with a proved equation $b=t$; it replaces both the
equality rule and the elimination of $\D{x}{y}$. In
\cref{disch:sec:inferred} neither the equation nor the term $t$ is
supplied: both are read off two rows of the matrix. In
\cref{disch:sec:subst} the fusion is carried out on the canonical recursor
syntax, where no normalization is available and reflexivity is derived. In
\cref{disch:sec:restricted} the equality rule substitutes for a variable
only, for the newest variable only, or, on raw terms, for a free parameter;
an obstruction shows that unary congruence is not enough. In
\cref{disch:sec:spec} the use of a clause takes no term argument. In
\cref{disch:sec:alias} a primitive fresh alias, whose discharge substitutes
into the conclusion, and substitution for a parameter in the backward
direction replace the equality rule; the source of equality and the target
of every use of a clause are then the newest parameter, and both new rules
can be restricted to literal equations. \Cref{disch:sec:const} records an induction interface that one of the
packages treats separately. \Cref{disch:sec:impl,disch:sec:sizes} summarize
the checkers, and \cref{disch:sec:further} lists what is left open.

Throughout, $\equiv$ is literal identity of syntax up to the names of bound
variables, a \emph{parameter} is a free numerical variable, and a
\emph{literal equation} is a formula whose root is an equation. The numerals
$\Suc 0$ and $\Suc\Suc 0$ are written $1$ and $2$, and $H$ is the equation
$0=0$.

\section{Normal forms, clause contexts, and the reference calculus}
\label{disch:sec:setting}

This section fixes what is common to \pkg{witness\_discharge},
\pkg{inferred\_witness}, \pkg{deterministic\_opening},
\pkg{first\_instance}, \pkg{clause\_discharge} and
\pkg{variable\_substitution}. Full normalization is the subject of
\cref{comp:sec:full} and the clause calculi on raw terms that of
\cref{claus:sec:five}; they are cited, and only what is specific to the
calculi of this chapter is stated here. The remark on the strict traversal
is from \pkg{first\_instance}.

\subsection{Conventions on full normal forms}

In \cref{disch:sec:setting,disch:sec:name,disch:sec:inferred} and in
\cref{disch:sec:vsub,disch:sec:newest} every term and every formula is in
full normal form: $\nf$ is the normal form for the reduction of
\cref{comp:def:reduction}, which exists and is unique by
\cref{comp:thm:sn,comp:thm:confluence}, and every displayed instance
$B(t)$, $B(t,r)$ or $F\subst{t}{b}$ is formed with the normalizing
substitution of \cref{comp:def:nusubst}. Substitution is simultaneous: an
inserted term is not traversed again, so that $F\subst{t}{b}$ is defined
also when $t$ contains $b$. By \cref{comp:prop:nfadm} the map $\nf$ is an
admissible normalizer; in particular it commutes with injective renaming,
it leaves the root of a formula unchanged, and it is coherent with
substitution,
\begin{equation}
 \nf(T\sigma)=\nf\bigl(\nf(T)[x\mapsto\nf(\sigma x)]\bigr),
 \label{disch:eq:coherence}
\end{equation}
so that instances compose (\cref{comp:lem:laws}) and it is immaterial
whether a compound expression is normalized before or after one of its
variables is instantiated. The counter of an iterator in a normal form is
neutral (\cref{comp:lem:fullgrammar}).

Let $P$ be the parity term of \cref{base:def:sel}, written $N$ there,
$P(y)=\It{y}{0}{u}{\sel(u,1,0)}$. On normal forms the computations of
\cref{base:eq:par012} and of the selector are literal identities:
\begin{equation}
 \sel(0,a,c)\equiv a,\qquad \sel(\Suc v,a,c)\equiv c,\qquad
 P(1)\equiv 1,\qquad P(2)\equiv 0.
 \label{disch:eq:endpoints}
\end{equation}
No induction, and no statement about $P$ at other arguments, is used
anywhere in this chapter.

\begin{remark}[Termination of a strict traversal]
\label{disch:rem:strict}
The checkers of this chapter normalize the counter, the seed and the body
first and, at a successor counter, normalize the predecessor iterate
$\It{n}{a}{u}{h}$ before inserting it into the body. If the body does not
mention its state variable, that value is erased by the root contraction,
so the strategy is not a reduction sequence of \cref{comp:def:reduction}
as it stands. \pkg{first\_instance} justifies it as follows. Add an
administrative term $\mathsf{let}\ u=a\ \mathsf{in}\ h$ with the reduction
to $h\subst{a}{u}$, and replace the successor contraction by the
\emph{strict expansion}
\[
 \It{\Suc n}{a}{u}{h}\longrightarrow
 \mathsf{let}\ v=\It{n}{a}{u}{h}\ \mathsf{in}\ h\subst{v}{u},
 \qquad v\text{ fresh},
\]
reducing under all constructors and under both components of a
$\mathsf{let}$. No new normal forms arise: every $\mathsf{let}$ has a root
contraction, and that contraction composes with a strict expansion to the
original successor contraction. A $\mathsf{let}$ is strongly normalizing
whenever its value is strongly normalizing and its body sends every
strongly normalizing term substituted for $u$ to a strongly normalizing
term; this is proved by induction on the sum of the reduction heights of
the two components, an internal contraction decreasing the sum and the
root contraction being covered by the hypothesis on the body. The proof of
\cref{comp:lem:snclosure} uses this fact at the successor root, after
obtaining the predecessor iterate by the same lexicographic induction, and
the structural induction of \cref{comp:thm:sn}, extended by the new
constructor, gives strong normalization of the auxiliary syntax. The
strict traversal is an evaluation strategy of the auxiliary calculus:
after a strict expansion it normalizes the value and then contracts the
$\mathsf{let}$. Hence it terminates, also when the value is erased, and
since forgetting the administrative terms preserves conversion it returns
the unique normal form $\nf(t)$.
\end{remark}

\begin{remark}[Bounded checkers]
\label{disch:rem:fuel}
By \cref{comp:prop:notpr} the normalizer is total recursive and not
primitive recursive; it is part of the specification of every calculus on
normal forms and is counted in every checker of such a calculus.
\pkg{witness\_discharge}, \pkg{inferred\_witness}, \pkg{first\_instance}
and \pkg{clause\_discharge} each supply a second checker with a fuel
field, as in \cref{comp:prop:fuel}. The metering differs from the one
there: one unit is consumed at every entry into the recursive traversal of
syntax and into the recursive proof check, not at every successor
contraction. The conclusions are the same, by the same argument. An
accepting unbounded run with $k$ metered entries is accepted with fuel $k$
and rejected with fuel $k-1$, and erasing the fuel turns an accepting
bounded run into an accepting unbounded run, so that
\[
 \exists p\ \operatorname{Accept}(A,p)
 \quad\Longleftrightarrow\quad
 \exists p\,\exists k\ \operatorname{Accept}_{\mathrm{fuel}}(A,p,k);
\]
and since every cycle of the call graph passes through a metered entry,
the bounded acceptance relation is primitive recursive. Nothing is
asserted about a primitive recursive bound on the fuel that a derivable
formula needs.
\end{remark}

\subsection{Formulas, clauses, and the two conventions}

The formulas of
\cref{disch:sec:setting,disch:sec:name,disch:sec:inferred} are the
$Q$-free formulas of \cref{claus:def:tau} over normal forms, with the
closure abbreviation $q_yA(y)=\D{w}{y}A(y)$, $w$ unused, of that
definition. Contexts are lists of clause occurrences
(\cref{base:sec:clauses,claus:sec:origins}). For a matrix $B(x,y)$ and a
term $r$ the \emph{opened clause} is
\[
 \mathcal C_B(r)\defeq[B(r,0);\,y.\,B(r,y)].
\]
Two conventions for assumptions and reflexivity occur in the packages.
\begin{itemize}
\item \emph{The convention with \rulename{Refl}.} Every clause has exactly
 one antecedent, an assumption $F$ is the clause
 $\operatorname{ord}(F)=[H;\,v.\,F]$ of \cref{claus:sec:origins}, and
 reflexivity is a rule.
\item \emph{The convention with the reflexivity clause.} Clauses are the
 uniform clauses of \cref{claus:def:v}, with at most one antecedent; an
 assumption $F$ is the unconditional clause $[\ ;\,v.\,F]$; every context
 contains one distinguished occurrence of the clause
 $\mathcal R=[\ ;\,y.\,y=y]$ of that definition; and there is no
 reflexivity rule: $t=t$ is a use of $\mathcal R$ at $t$.
\end{itemize}
In both, $\Gamma,F$ denotes $\Gamma$ extended by the assumption clause of
$F$. \Cref{disch:lem:conventions} shows that the two conventions are
interchangeable. All proofs below are written in the convention with
\rulename{Refl}.

\subsection{The shared rules and the reference calculus}

The rules \rulename{Use}, \rulename{Elim} and \rulename{Intro} are those of
\cref{claus:def:c3}, read on full normal forms with normalizing
substitution; \rulename{Intro} is the witness-selected introduction of
\cref{comp:def:wi} with a clause as its assumption. The equality rule of
the reference calculus is \rulename{Tr}, as in $\Nat$.

\begin{definition}[Shared rules]
\label{disch:def:rules}
\begin{align}
 &\infrule{}{\Gamma\proves t=t}\ \rulename{Refl}
 \label{disch:eq:refl}\\[1ex]
 &\infrule{\Gamma\proves A}{\Gamma\proves B(t)}\ \rulename{Use},
 \qquad [A;\,y.\,B(y)]\text{ an occurrence in }\Gamma
 \label{disch:eq:use}\\[1ex]
 &\infrule{\Gamma,\ B(t,\eta)\proves B(t,\Suc z)}
          {\Gamma\proves\D{x}{y}B(x,y)}\ \rulename{Intro},
 \qquad
 \eta=\begin{cases}
 z,&B\text{ is a literal equation and }t\equiv 0,\\
 0,&\text{otherwise}
 \end{cases}
 \label{disch:eq:intro}\\[1ex]
 &\infrule{\Gamma\proves\D{x}{y}B(x,y)\qquad
           \Gamma,\ \mathcal C_B(b)\proves A}
          {\Gamma\proves A}\ \rulename{Elim}
 \label{disch:eq:elim}\\[1ex]
 &\infrule{\Gamma\proves s=t\qquad\Gamma\proves h(s)=r}
          {\Gamma\proves r=h(t)}\ \rulename{Tr}
 \label{disch:eq:tr}
\end{align}
In \rulename{Use} the term $t$ is arbitrary; for an unconditional clause
the rule has no premise. In \rulename{Intro} the variable $z$ is fresh for
$\Gamma$, for the witness $t$ and for $\D{x}{y}B$; the instance with
$\eta=0$ is the \emph{case mode} and the instance with $\eta=z$ the
\emph{progress mode}, and the mode is determined by the literal syntax of
the normal forms, not by provable equality. In \rulename{Elim} the
variable $b$ is fresh for $\Gamma$, $A$ and $\D{x}{y}B$. In \rulename{Tr},
$h(w)$ is an arbitrary term template.
\end{definition}

\begin{definition}[The reference calculus]
\label{disch:def:reference}
$\mathcal T$ is the calculus with the five rules \rulename{Refl},
\rulename{Use}, \rulename{Intro}, \rulename{Elim} and \rulename{Tr}, in the
convention with \rulename{Refl}.
\end{definition}

$\mathcal T$ is the fully normalized clause calculus with which the
packages of this section compare themselves. It is the clause form of the
six-rule calculus $\mathsf F$ of \cref{comp:def:fsix}.

\begin{proposition}[The reference calculus and the kernel]
\label{disch:prop:kernel}
For a native formula $A$ let $A^{\flat}\defeq\tau(\nf(A))$, with $\tau$
the $Q$-free form of \cref{claus:def:tau}. Then $\Nat$ derives
$\Gamma\proves A$ if and only if $\mathcal T$ derives $A^{\flat}$ from the
assumption clauses of the formulas $G^{\flat}$, $G\in\Gamma$; and $\Nat$
derives $A$ from $A^\flat$ and conversely. The translation is the identity
on normalized $Q$-free formulas and fixes $0=\Suc 0$.
\end{proposition}

\begin{proof}
$\Nat$ derives $A$ from $\nf(A)$ and conversely by
\cref{comp:lem:conversion}, and $\nf(A)$ from $A^\flat$ and conversely by
\cref{claus:lem:unq}(iii). By \cref{comp:thm:full}, $\Nat$ and the kernel
$\Nat^{\mathrm f}$ of \cref{comp:def:nfull} on full normal forms have the
same consequences, and by \cref{comp:thm:fsix} so does $\mathsf F$; it
remains to show, for
$\Gamma$ and $A$ in full normal form, that $\Gamma\proves_{\mathsf F}A$ if
and only if $\operatorname{ord}(\tau\Gamma)\proves_{\mathcal T}\tau A$.
This is the argument of \cref{claus:thm:five} with $\mathsf F$ in the
place of $\Nat$ and $\mathcal T$ in the place of $\mathsf C_3$; it is
shorter, because $\mathsf F$ and $\mathcal T$ have the same equality rule
and the same introduction.

From $\mathcal T$ to $\mathsf F$, for an annotated context
(\cref{claus:def:origin}) and its native image: \rulename{Refl} and
\rulename{Tr} are rules of $\mathsf F$; a \rulename{Use} of an assumption
occurrence is \rulename{Hyp}, and a \rulename{Use} of an opened occurrence
is \rulename{QE}; \rulename{Elim} is \rulename{DE}; and \rulename{Intro}
is \rulename{WI} in the same mode, its new occurrence being an assumption.

From $\mathsf F$ to $\mathcal T$: \rulename{Hyp} is \rulename{Refl} at $0$
followed by \rulename{Use}; \rulename{Refl} and \rulename{Tr} are rules;
\rulename{QE} is the second rule of \cref{disch:lem:closure} below, since
$\tau(\Q{y}B)\equiv q_y\tau B$; and \rulename{WI} is \rulename{Intro} in
the same mode, because $\tau$ sends literal equations to literal equations
and all other formulas to formulas with root $D$. For \rulename{DE} apply
\rulename{Elim} to the translated major premise. In the extended context,
under the assumption $\tau B(b,0)$, the opened clause at $\Suc z$ gives
$\tau B(b,\Suc z)$, so the first rule of \cref{disch:lem:closure} gives
$q_y\tau B(b,y)$, which is $\tau(\Q{y}B(b,y))$; weaken the translated
branch and cut this formula into it (\cref{disch:lem:structural}).

Finally, from $\operatorname{ord}(\tau\Gamma)\proves_{\mathcal T}\tau A$
the first translation gives $\tau\Gamma\proves_{\mathsf F}\tau A$, and
\cref{claus:lem:unq}(iii), which holds in $\mathsf F$ by
\cref{comp:thm:fsix}, with cut replaces $\tau\Gamma$ and $\tau A$ by
$\Gamma$ and $A$.
\end{proof}

\subsection{Structural operations}

\begin{lemma}[Weakening, renaming, cut]
\label{disch:lem:structural}
Let $\mathcal S$ be a calculus whose rules are among those of
\cref{disch:def:rules} and the rules \rulename{Name}
\eqref{disch:eq:name}, \rulename{Open} \eqref{disch:eq:open},
\rulename{Export} \eqref{disch:eq:export} and \rulename{VSub}
\eqref{disch:eq:vsub} introduced below, and which contains \rulename{Refl}
and \rulename{Use}. The following are admissible in $\mathcal S$, by
direct induction on derivations.
\begin{enumerate}[label=\textup{(\alph*)}]
\item Weakening: if $\Gamma\proves A$ and every occurrence of $\Gamma$ is
 an occurrence of $\Gamma'$, then $\Gamma'\proves A$.
\item Injective renaming: if $\Gamma\proves A$ and $\iota$ is an injective
 map of parameters, then $\Gamma\iota\proves A\iota$.
\item Cut of an assumption: if $\Gamma\proves F$ and $\Gamma,F\proves A$,
 then $\Gamma\proves A$.
\end{enumerate}
None of the three substitutes a term for a parameter in a derivation.
\end{lemma}

\begin{proof}
The inductions are those of \cref{claus:lem:struct}; for \rulename{Refl},
\rulename{Use}, \rulename{Elim}, \rulename{Intro} and \rulename{Tr} nothing
changes, since normalization commutes with injective renaming. Only the
side conditions of the four new rules have to be checked under an
injective renaming. Such a renaming preserves and reflects the property of
being a variable, the literal identities tested by the selector of
\cref{disch:def:selector}, and the absence of a parameter from a formula.
Hence the selected witness of \rulename{Open} and \rulename{Export}, the
source of \rulename{VSub} and the export conditions of \rulename{Name} and
\rulename{Export} are preserved. For weakening and cut, no side condition
of the new rules refers to the clauses of the context other than through
freshness, and the selector depends on the matrix of the major premise
alone.
\end{proof}

Substitution of a term for a parameter does not preserve rule instances in
general. It can turn a witness into the literal $0$ and so change the mode
of \rulename{Intro} (\cref{comp:ex:modechange,claus:prop:subst}), it can
make the test of the selector succeed, and it can turn the variable source
of \rulename{VSub} into a compound term. Admissibility of substitution is
proved for the calculus $\mathcal O$ in \cref{disch:thm:subst} and follows
for the others from the equivalences.

\subsection{Closures, numerical cases, and the two introductions}

The lemmas of this subsection and the next are the counterparts on full
normal forms of \cref{claus:lem:closure,claus:lem:cases,claus:lem:eq} and
\cref{claus:lem:ftransport}. They are restated because the calculi of
\cref{disch:sec:name,disch:sec:inferred} do not have \rulename{Elim} as a
rule, and one of them has it at first for two kinds of matrices only; the
hypotheses record what each argument uses.

\begin{lemma}[Closure and numerical cases]
\label{disch:lem:closure}
Let $\mathcal S$ contain \rulename{Use} and \rulename{Intro} and admit
weakening, and suppose that \rulename{Elim} is a primitive or derived rule
of $\mathcal S$ for every matrix in which the witness coordinate does not
occur. Then the following rules are derived, for every formula $A(y)$,
every term $t$ and $z$ fresh:
\[
 \infrule{\Gamma,\ A(0)\proves A(\Suc z)}{\Gamma\proves q_yA(y)}
 \qquad
 \infrule{\Gamma\proves q_yA(y)\qquad\Gamma\proves A(0)}
         {\Gamma\proves A(t)}
 \qquad
 \infrule{\Gamma\proves A(0)\qquad\Gamma\proves A(\Suc z)}
         {\Gamma\proves A(t)}
\]
No equality rule is used.
\end{lemma}

\begin{proof}
The second rule is \cref{claus:lem:closure}(ii), whose proof eliminates a
matrix with an unused witness coordinate, and the third follows from the
first two as in \cref{claus:cor:induction}. The first rule is obtained
differently from \cref{claus:lem:closure}(i), with no appeal to an
unrestricted case introduction: apply \rulename{Intro} to the matrix
$A(y)$, with the unused witness coordinate, at the witness $1$; since
$1\not\equiv 0$ the rule is in case mode at every root, and its premise is
the given one.
\end{proof}

The third rule is a distinction of numerical cases, not an induction: its
successor premise has no hypothesis.

For a literal equation $E$ let $G(v,E)$ be the gate of
\cref{comp:lem:envelope}, with $G(0,E)\equiv H$ and
$G(\Suc w,E)\equiv E$, and let the \emph{envelope} in the present
formulas be $L(E)\defeq q_vG(v,E)$.

\begin{lemma}[The two unrestricted introductions]
\label{disch:lem:repair}
Let $\mathcal S$ contain \rulename{Refl}, \rulename{Use} and
\rulename{Intro}, admit weakening and cut, and have \rulename{Elim} as a
primitive or derived rule for every matrix in which the witness coordinate
does not occur and for every matrix with root $D$. Then for every matrix
$B$, every literal equation $E(x,y)$ and every witness $t$ the rules
\[
 \infrule{\Gamma,\ B(t,0)\proves B(t,\Suc z)}
         {\Gamma\proves\D{x}{y}B(x,y)}
 \qquad
 \infrule{\Gamma,\ E(t,z)\proves E(t,\Suc z)}
         {\Gamma\proves\D{x}{y}E(x,y)}
\]
are derived, with $z$ fresh as in \rulename{Intro}.
\end{lemma}

\begin{proof}
\emph{Case introduction.} \rulename{Intro} is the first rule unless $B$ is
a literal equation $E(x,y)$ and $t\equiv 0$. For that case the derivation
is the one given after \cref{comp:lem:envelope} for the missing case of
\cref{comp:thm:fsix}, read in the present formulas: $\Q{v}$ becomes $q_v$;
the two derivations $\Gamma,E\proves L(E)$ and $\Gamma,L(E)\proves E$ of
\cref{comp:lem:envelope} use the first and the second rule of
\cref{disch:lem:closure}, with \rulename{Refl} for the base $H$; the
matrix $L(E(x,y))$ has root $D$, so it is introduced at the witness $0$ in
case mode and opened by \rulename{Elim} at a fresh $c$; and the retained
closure is the opened clause, used by \rulename{Use} at $\Suc z$. The
derivation eliminates matrices with an unused witness coordinate and one
matrix with root $D$, and no others.

\emph{Progress introduction.} This is the third statement of
\cref{claus:cor:induction}, with the same proof: \rulename{Intro} for the
matrix $E(t,y)$ at the literal witness $0$ of an unused coordinate is in
progress mode and yields $q_yE(t,y)$; in the context $\Gamma,E(t,0)$ the
second rule of \cref{disch:lem:closure} gives $E(t,\Suc z')$; and the case
introduction at the witness $t$ concludes.
\end{proof}

\subsection{Equality from one term transport}

\emph{Uncrossed} transport is the rule
\begin{equation}
 \infrule{\Gamma\proves s=t\qquad\Gamma\proves h(s)=r}
         {\Gamma\proves h(t)=r}\ \rulename{UTr}
 \label{disch:eq:utr}
\end{equation}
for an arbitrary term template $h(w)$.

\begin{lemma}[Equality algebra]
\label{disch:lem:equality}
Let $\mathcal S$ have \rulename{Refl}, as a rule or derived. If
\rulename{UTr} is derivable in $\mathcal S$, or if \rulename{Tr} is, then
so are both, and so are symmetry, transitivity, the congruence
$s=t\proves h(s)=h(t)$, and transport in every equation template:
$s=t,\ p(s)=q(s)\proves p(t)=q(t)$.
\end{lemma}

On raw terms the same consequences are \cref{claus:lem:eq};
what is specific here is that one term transport, crossed or uncrossed,
suffices once reflexivity is available.

\begin{proof}
From \rulename{UTr}: with $h(w)=w$ and $r\equiv s$ the premises are $s=t$
and the reflexive $s=s$, and the conclusion is $t=s$; this is symmetry, and
\rulename{UTr} followed by symmetry is \rulename{Tr}. From \rulename{Tr}
and reflexivity, \cref{comp:lem:eqalg}(b) gives symmetry, transitivity,
congruence and transport in an equation template, and \rulename{Tr}
followed by symmetry is \rulename{UTr}.
\end{proof}

\begin{lemma}[Transport in a formula]
\label{disch:lem:formula}
Let $\mathcal S$ contain \rulename{Use} and \rulename{Intro}, admit
weakening, have \rulename{Elim} as a primitive or derived rule for all
matrices, and derive symmetry and transport in every equation template.
Then for every formula template $F(w)$
\[
 \infrule{\Gamma\proves s=t\qquad\Gamma\proves F(s)}{\Gamma\proves F(t)}
\]
is derived in $\mathcal S$.
\end{lemma}

\begin{proof}
The proof of \cref{claus:lem:ftransport} applies, by induction on the
syntax of the template. Its one appeal to \cref{claus:lem:cases} is an
\rulename{Intro} at the fresh parameter opened by \rulename{Elim}; a
parameter is not the literal $0$, so that instance is in case mode
whatever the root of the matrix is, and the case introduction at the
witness $0$ of \cref{disch:lem:repair} is not needed.
\end{proof}

\subsection{Two-row packets}

The device by which a calculus without a general equality rule recovers
\rulename{Tr} is a paired formula whose matrix has prescribed rows at
$y=0,1,2$. It first occurs in \pkg{variable\_substitution}.

\begin{definition}[Packet]
\label{disch:def:packet}
Let $E_1(x)=(a_1=c_1)$ and $E_2(x)=(a_2=c_2)$ be equations in which $x$
may occur and $y$ does not. The \emph{packet} with rows $E_1,E_2$ is the
matrix
\[
 \Pi(x,y)\defeq\Bigl(\sel\bigl(y,0,\sel(P(y),a_2,a_1)\bigr)
 =\sel\bigl(y,0,\sel(P(y),c_2,c_1)\bigr)\Bigr).
\]
\end{definition}

By \eqref{disch:eq:endpoints} its rows are literally
\begin{equation}
 \Pi(x,0)\equiv H,\qquad \Pi(x,1)\equiv E_1(x),\qquad
 \Pi(x,2)\equiv E_2(x).
 \label{disch:eq:rows}
\end{equation}

\begin{lemma}[Introduction of a packet]
\label{disch:lem:packet}
Let $\mathcal S$ contain \rulename{Use} and \rulename{Intro}, admit
weakening, and have the numerical-case rule of \cref{disch:lem:closure}. If
$\Gamma\proves E_1(s)$ and $\Gamma\proves E_2(s)$, then
$\Gamma\proves\D{x}{y}\Pi(x,y)$. In any context that contains the opened
clause $\mathcal C_\Pi(b)$, \rulename{Refl} and \rulename{Use} derive
$E_1(b)$ and $E_2(b)$.
\end{lemma}

\begin{proof}
Apply \rulename{Intro} at the witness $s$. Its goal is
\[
 \Pi(s,\Suc z)\equiv
 \bigl(\sel(P(\Suc z),a_2(s),a_1(s))=\sel(P(\Suc z),c_2(s),c_1(s))\bigr).
\]
Let $A(v)\defeq\bigl(\sel(v,a_2(s),a_1(s))=\sel(v,c_2(s),c_1(s))\bigr)$.
Then $A(0)\equiv E_2(s)$, $A(\Suc z')\equiv E_1(s)$ and, by
\eqref{disch:eq:coherence}, $A(P(\Suc z))\equiv\Pi(s,\Suc z)$. The
numerical-case rule at the target $P(\Suc z)$, with the two given
derivations weakened to the local context and an eigenparameter $z'$
chosen fresh independently of $z$, proves the goal. The assumption of
\rulename{Intro} is not used, so the derivation is valid in either mode, in
particular when $s\equiv 0$. The second statement is \eqref{disch:eq:rows}:
the antecedent of $\mathcal C_\Pi(b)$ is $H$, proved by \rulename{Refl},
and \rulename{Use} at $1$ and at $2$ gives the two rows at $b$.
\end{proof}

\subsection{Interchangeability of the two conventions}

\begin{lemma}[Reflexivity as a fixed clause]
\label{disch:lem:conventions}
Let $\mathcal S$ be a calculus in the convention with \rulename{Refl}
whose other rules are among \rulename{Use}, \rulename{Intro},
\rulename{Elim}, \rulename{Tr}, \rulename{Name}, \rulename{Open},
\rulename{Export} and \rulename{VSub}, and let $\mathcal S^{\mathcal R}$
have the same rules without \rulename{Refl}, in the convention with the
reflexivity clause. For a context $\Delta$ of uniform clauses let
$L(\Delta)$ be the context $\Delta^{\circ}$ of \cref{claus:def:v}, which
replaces every unconditional occurrence $[\ ;\,y.\,B]$ by $[H;\,y.\,B]$.
Then for every formula $F$
\[
 \mathcal R,\Delta\proves_{\mathcal S^{\mathcal R}}F
 \quad\Longleftrightarrow\quad
 L(\Delta)\proves_{\mathcal S}F .
\]
The statement covers arbitrary contexts, not only those that arise in a
derivation from $\mathcal R$ alone.
\end{lemma}

\begin{proof}
The two inductions of the proof of \cref{claus:thm:uniform} apply. They
translate only the uses of $\mathcal R$, the uses of unconditional
clauses, \rulename{Refl}, and the assumption added by \rulename{Intro},
and they copy every other rule; occurrences correspond by position, also
when two have the same value. The rules that are copied here are
\rulename{Tr}, \rulename{VSub}, and the rules \rulename{Elim},
\rulename{Name}, \rulename{Open} and \rulename{Export}, each of which adds
the clause $\mathcal C_B(r)$ with one antecedent on both sides and has the
same premises and side conditions in the two conventions.
\end{proof}

\section{Naming packets: equality at the discharge of a witness}
\label{disch:sec:name}

This section draws on \pkg{witness\_discharge}. The packet is that of
\pkg{variable\_substitution} (\cref{disch:sec:vsub}); its final
substitution and the elimination of its witness become one inference here.

\subsection{The rule \texorpdfstring{\rulename{Name}}{Name}}

Equality substitution is performed only while the fresh witness of a
paired formula is discharged:
\begin{equation}
 \infrule{\Gamma\proves\D{x}{y}B(x,y)\qquad
          \Gamma,\ \mathcal C_B(b)\proves b=t\qquad
          \Gamma,\ \mathcal C_B(b)\proves F}
         {\Gamma\proves F\subst{t}{b}}\ \rulename{Name}
 \label{disch:eq:name}
\end{equation}
The variable $b$ is fresh for $\Gamma$ and $\D{x}{y}B$, and the exported
formula $F\subst{t}{b}$ must not contain $b$. The two branch premises have
the same extended context. The term $t$ and the formula $F$ may contain
$b$; in particular $t\equiv b$ is allowed. The side condition concerns the
exported normal form and is not a prohibition of $b$ in $t$.

\begin{definition}[The calculus of witness discharge]
\label{disch:def:name}
$\mathcal W$ is the calculus with the four rules \rulename{Refl},
\rulename{Use}, \rulename{Intro} and \rulename{Name}, in the convention
with \rulename{Refl}. $\mathcal W^{\mathcal R}$ has the three rule families
\rulename{Use}, \rulename{Intro} and \rulename{Name} and the fixed clause
$\mathcal R$. There is no further logical, equality, computation, case,
induction, weakening or cut rule.
\end{definition}

\begin{lemma}[Soundness of \rulename{Name}]
\label{disch:lem:name-sound}
\rulename{Name} is sound for the ordinary reading in $\Teq$
(\cref{base:sec:teq}), for every formula $F$ and also when $t$ contains
$b$.
\end{lemma}

\begin{proof}
Eliminate the existential quantifier in the reading of the major premise
at a fresh $b$; the local assumption is the reading of $\mathcal C_B(b)$.
The two branch derivations give $b=t$ and $F$. Equality substitution in
the template $F\subst{w}{b}$ gives the raw instance at $t$, and the
computation equations identify it with its normal form $F\subst{t}{b}$.
This step does not need $b\notin\FV(t)$: an equation $b=t(b)$ supports the
simultaneous substitution of $t(b)$ for $b$ in a fixed formula. What the
elimination of the quantifier needs is the absence of $b$ from the result,
which is the side condition.
\end{proof}

\rulename{Use} has the reading of its clause, and \rulename{Intro} is sound
by coverage in case mode and by induction for literal equations in
progress mode.

\begin{lemma}[Ordinary elimination]
\label{disch:lem:name-elim}
\rulename{Elim} is a derived rule of $\mathcal W$.
\end{lemma}

\begin{proof}
Take $t\equiv b$ in \rulename{Name}. The equality premise is
\rulename{Refl}, the other premise is the given derivation of $A$, and
$A\subst{b}{b}\equiv A$ does not contain $b$.
\end{proof}

By \cref{disch:lem:structural,disch:lem:name-elim} the closure rules, the
numerical cases, the two unrestricted introductions and the introduction
of packets
(\cref{disch:lem:closure,disch:lem:repair,disch:lem:packet}) are available
in $\mathcal W$. None of them uses an instance of \rulename{Name} other
than with $t\equiv b$.

\subsection{Transport from one discharge}

\begin{theorem}[Transport in $\mathcal W$]
\label{disch:thm:name-transport}
For every term template $h(w)$ the rule \rulename{UTr}
\eqref{disch:eq:utr} is derived in $\mathcal W$. Apart from ordinary
eliminations, its derivation contains one instance of \rulename{Name}.
\end{theorem}

\begin{proof}
Let $\Gamma\proves s=t$ and $\Gamma\proves h(s)=r$, let $x$ be fresh for
$h$, $r$ and $t$, and let $\Pi$ be the packet with the rows
\[
 E_1(x)=(h(x)=r),\qquad E_2(x)=(x=t).
\]
The two given equations are $E_1(s)$ and $E_2(s)$, so
$\Gamma\proves\D{x}{y}\Pi$ by \cref{disch:lem:packet}. Apply
\rulename{Name} with a fresh $b$. In the context $\Gamma,\mathcal C_\Pi(b)$
the same lemma gives the row $b=t$, which is the equality premise, and the
row $h(b)=r$, which is taken as $F$. The exported formula is
\[
 (h(b)=r)\subst{t}{b}\equiv(h(t)=r)
\]
by \eqref{disch:eq:coherence}, because $b$ occurs neither in $h$ nor in
$r$; it does not contain $b$. The identity is one of normal forms and not
an application of an equality rule.
\end{proof}

\begin{corollary}[Equality in $\mathcal W$]
\label{disch:cor:name-equality}
$\mathcal W$ derives symmetry, transitivity, congruence, \rulename{Tr}, and
transport in every formula template.
\end{corollary}

\begin{proof}
\Cref{disch:lem:equality} applies by \cref{disch:thm:name-transport}; in
particular symmetry is obtained from the packet with the identity template
and a reflexive minor premise, and is not assumed by it.
\Cref{disch:lem:formula} applies by \cref{disch:lem:name-elim}.
\end{proof}

\subsection{Equivalence}

\begin{theorem}[$\mathcal W$ and the reference calculus]
\label{disch:thm:name}
For every context $\Gamma$ and formula $F$,
\[
 \Gamma\proves_{\mathcal W}F\quad\Longleftrightarrow\quad
 \Gamma\proves_{\mathcal T}F .
\]
Both directions are effective transformations of derivations that leave
the end sequent unchanged.
\end{theorem}

\begin{proof}
From $\mathcal T$ to $\mathcal W$, translate a derivation recursively:
\rulename{Elim} is replaced by \cref{disch:lem:name-elim} and
\rulename{Tr} by \cref{disch:cor:name-equality}; the other three rules are
shared. The constructions that replace a \rulename{Tr} node were derived
without \rulename{Tr} and without a translated node, so the recursion is
on the given derivation alone.

From $\mathcal W$ to $\mathcal T$: $\mathcal T$ derives symmetry and
\rulename{UTr} from \rulename{Tr} by \cref{disch:lem:equality}, and hence
transport in every formula template by \cref{disch:lem:formula}. At a
\rulename{Name} node the translated branch premises give $b=t$ and $F$ in
the context $\Gamma,\mathcal C_B(b)$. Transport in the template
$F\subst{w}{b}$, $w$ fresh, gives $F\subst{t}{b}$ in that context; the
substitution is simultaneous, so this is correct also when $t$ contains
$b$. By the side condition the result does not contain $b$, and
\rulename{Elim} discharges the clause. The finitely many eigenparameters
are renamed before each replacement.
\end{proof}

\begin{corollary}[Three rule families]
\label{disch:cor:name-three}
For every context $\Delta$ of clauses with at most one antecedent and every
formula $F$,
$\mathcal R,\Delta\proves_{\mathcal W^{\mathcal R}}F$ if and only if
$L(\Delta)\proves_{\mathcal T}F$. Consequently $\Nat$ derives
$\Gamma\proves A$ if and only if $\mathcal W^{\mathcal R}$ derives
$A^{\flat}$ from $\mathcal R$ and the unconditional assumption clauses of
the formulas of $\Gamma^{\flat}$, and the ordinary reading of
$\mathcal W^{\mathcal R}$ is bi-interpretable with $\iPRA$ in the sense of
\cref{base:thm:bridge}: whole domains, actual equality, both numerical
composites with graph $x=y$, and both formula round trips.
\end{corollary}

\begin{proof}
The first statement is \cref{disch:lem:conventions,disch:thm:name}; the
second adds \cref{disch:prop:kernel}. The comparisons are the identity on
normalized terms and formulas and fix $0=\Suc 0$, so the interpretations
and composites of \cref{base:thm:bridge} are unchanged.
\end{proof}

\begin{remark}[What the count of three means; Lean]
\label{disch:rem:name-lean}
The number three counts inference patterns after reflexivity has been moved
into one explicitly listed axiom clause; it does not say that reflexivity
is proved from nothing. The reduction specific to \pkg{witness\_discharge}
is the localization of all equality elimination at the discharge of a
witness; the clauses with at most one antecedent and the fixed clause are
choices of presentation taken from \pkg{uniform\_clauses}.

In Lean the kernel form is \Lean{NMPrf}: on full normal forms, with
\rulename{Name} in place of \rulename{Tr} and \rulename{DE} and with the
hypothesis $\Q{y}B(b,y)$ in the opened context, the kernel has the six
rules \rulename{Hyp}, \rulename{Refl}, \rulename{Name}, \rulename{QE},
\rulename{DI} and \rulename{QI}, and \Lean{nmprf\_iff\_nprf}
(\Lean{NativeWitnessDischarge}) proves it equivalent to $\Nat$; the packet
is \Lean{NMPrf.packet}. The clause form $\mathcal W^{\mathcal R}$ is
\Lean{WDPrf}, with \Lean{wdprf\_iff\_nprf} (\Lean{NativeNameClauses}).
\end{remark}

\begin{remark}[Substitution in $\mathcal W$]
The substitution inside \rulename{Name} acts on the proved formula and is
justified by the proved equation; it is not a substitution into the branch
derivation. Substitution of a term for a parameter in a derivation of
$\mathcal W$ is admissible through \cref{disch:thm:name} and the
admissibility in $\mathcal O$ (\cref{disch:thm:subst,disch:thm:inferred}),
not by an induction that preserves the rule at each node.
\end{remark}

\section{The inferred witness}
\label{disch:sec:inferred}

This section draws on \pkg{inferred\_witness},
\pkg{deterministic\_opening}, \pkg{first\_instance} and
\pkg{clause\_discharge}. The four packages contain one idea in two forms;
\cref{disch:rem:four} records how they differ.

\subsection{The selector and the two forms of the rule}

\rulename{Name} has three premises, a supplied term $t$ and a supplied
formula $F$. The equation $b=t$ need not be proved: it can be read off the
opened clause.

\begin{definition}[Selector]
\label{disch:def:selector}
For a matrix $B(x,y)$ and a variable $b$ that does not occur in
$\D{x}{y}B$, the \emph{selected witness} is
\[
 \rho_B(b)\defeq
 \begin{cases}
 r,&B(b,0)\equiv H\ \text{ and }\ B(b,1)\equiv(b=r),\\
 b,&\text{otherwise.}
 \end{cases}
\]
In the first case the test \emph{succeeds}; in the second the selector
\emph{falls back}.
\end{definition}

The test inspects two rows of the matrix, at $y=0$ and at $y=1$, and
compares normal forms literally. The left side of the row at $1$ must be
the variable $b$ itself; an equation $r=b$, an equation provably of that
shape, and an equation whose left side merely contains $b$ are not
recognized. The right side $r$ may contain $b$. There is no search, no
decision of provable equality, and no choice left to the certificate.

\begin{lemma}[Properties of the selector]
\label{disch:lem:selector}
\begin{enumerate}[label=\textup{(\alph*)}]
\item For an injective renaming $\iota$,
 $\rho_{B\iota}(\iota b)\equiv\rho_B(b)\iota$, and the test succeeds for
 $B\iota$ at $\iota b$ if and only if it succeeds for $B$ at $b$.
\item If $x$ does not occur in $B$, or if $B$ has root $D$, the selector
 falls back.
\item If the test succeeds for $B$ with $\rho_B(b)\equiv r$, and $\sigma$
 is a substitution with $\sigma b\equiv b$ under which $b$ remains fresh
 for $\D{x}{y}B\sigma$, then it succeeds for $B\sigma$ with
 $\rho_{B\sigma}(b)\equiv r\sigma$.
\item A fallback is not stable under substitution.
\end{enumerate}
\end{lemma}

\begin{proof}
(a) Normalization commutes with $\iota$, and $\iota$ preserves and
reflects literal identity. (b) If $x$ does not occur in $B$, the row
$B(b,1)$ does not contain $b$ and is not an equation with left side $b$.
If $B$ has root $D$, so has $B(b,0)$, because normalization does not remove
a formula constructor, and it is not $H$. (c) By
\eqref{disch:eq:coherence}, $(B\sigma)(b,0)\equiv B(b,0)\sigma\equiv
H\sigma\equiv H$ and
$(B\sigma)(b,1)\equiv(b=r)\sigma\equiv(b=r\sigma)$. (d) For a parameter
$a$ and a term $r$ let $B(x,y)\defeq\bigl(\sel(y,a,x)=\sel(y,a,r)\bigr)$.
Then $B(b,0)\equiv(a=a)\not\equiv H$ and the selector falls back, while
under $a\mapsto 0$ the rows are $H$ and $b=r$ and the test succeeds.
\end{proof}

\begin{lemma}[The selected equation]
\label{disch:lem:selected}
In every context that contains $\mathcal C_B(b)$, the equation
$b=\rho_B(b)$ is derived by \rulename{Refl} and \rulename{Use}.
\end{lemma}

\begin{proof}
If the selector falls back, the equation is $b=b$. If the test succeeds,
the antecedent of $\mathcal C_B(b)$ is $H$, proved by \rulename{Refl}, and
\rulename{Use} at $1$ gives $B(b,1)\equiv(b=\rho_B(b))$.
\end{proof}

The witness can be used in two ways. Either the clause is opened at the
selected witness and the goal is unchanged,
\begin{equation}
 \infrule{\Gamma\proves\D{x}{y}B(x,y)\qquad
          \Gamma,\ \mathcal C_B(\rho_B(b))\proves A}
         {\Gamma\proves A}\ \rulename{Open}
 \label{disch:eq:open}
\end{equation}
with $b$ fresh for $\Gamma$, $A$ and $\D{x}{y}B$; or the clause is opened
at the fresh variable and the selected witness is substituted in the
exported formula,
\begin{equation}
 \infrule{\Gamma\proves\D{x}{y}B(x,y)\qquad
          \Gamma,\ \mathcal C_B(b)\proves F}
         {\Gamma\proves F\subst{\rho_B(b)}{b}}\ \rulename{Export}
 \label{disch:eq:export}
\end{equation}
with $b$ fresh for $\Gamma$ and $\D{x}{y}B$ and $b$ not in the exported
formula. In \rulename{Open} the parameter $b$ stays available in the
branch whether or not the selector returns it, and the clause
$\mathcal C_B(\rho_B(b))=[B(\rho_B(b),0);\,y.\,B(\rho_B(b),y)]$ may contain
it. Both rules have two premises and no term argument. Against
\rulename{Name}, the rule \rulename{Export} drops the equality premise and
the term $t$; the rule \rulename{Open} drops in addition the formula $F$.

\begin{definition}[The two calculi of the inferred witness]
\label{disch:def:inferred}
$\mathcal O$ has the rules \rulename{Refl}, \rulename{Use},
\rulename{Intro} and \rulename{Open}; $\mathcal X$ has the rules
\rulename{Refl}, \rulename{Use}, \rulename{Intro} and \rulename{Export}.
Their forms $\mathcal O^{\mathcal R}$ and $\mathcal X^{\mathcal R}$ with
the reflexivity clause have three rule families each.
\end{definition}

\begin{proposition}[Soundness]
\label{disch:prop:inferred-sound}
\rulename{Open} and \rulename{Export} are sound for the ordinary reading in
$\Teq$, also when $\rho_B(b)$ contains $b$. Moreover every instance of
\rulename{Export} is derivable in $\mathcal W$ with the same end sequent.
\end{proposition}

\begin{proof}
Eliminate the existential quantifier of the major premise at $b$; the
local assumption is the reading of $\mathcal C_B(b)$, and
\cref{disch:lem:selected} gives $b=\rho_B(b)$. For \rulename{Export},
equality substitution gives $F\subst{\rho_B(b)}{b}$, which does not contain
$b$. For \rulename{Open}, equality substitution in the template
$\mathcal C_B(v)^\circ$ gives the reading of $\mathcal C_B(\rho_B(b))$, the
branch gives $A$, and $A$ does not contain $b$. No occurs check on the
selected witness is needed: the rules do not assert that the equation
$b=r(b)$ has a solution independently of the major premise. The last
statement: keep the two premises of \rulename{Export}, insert the
derivation of \cref{disch:lem:selected}, and apply \rulename{Name} with
$t\equiv\rho_B(b)$.
\end{proof}

\subsection{The export form}

\begin{lemma}[Ordinary elimination in $\mathcal X$]
\label{disch:lem:export-elim}
\rulename{Elim} is a derived rule of $\mathcal X$, by one
\rulename{Export}.
\end{lemma}

\begin{proof}
Apply \rulename{Export} with $F\equiv A$. Since $A$ does not contain $b$,
$A\subst{\rho_B(b)}{b}\equiv A$ whichever term the selector returns. The
argument does not need the selector to fall back.
\end{proof}

Hence \cref{disch:lem:closure,disch:lem:repair,disch:lem:packet} hold in
$\mathcal X$.

\begin{theorem}[Transport in $\mathcal X$]
\label{disch:thm:export-transport}
\rulename{UTr} is derived in $\mathcal X$ for every term template $h(w)$.
\end{theorem}

\begin{proof}
Let $\Gamma\proves s=t$ and $\Gamma\proves h(s)=r$, and let $\Pi$ be the
packet with the rows, exchanged against \cref{disch:thm:name-transport},
\[
 E_1(x)=(x=t),\qquad E_2(x)=(h(x)=r).
\]
\Cref{disch:lem:packet} gives $\Gamma\proves\D{x}{y}\Pi$; the equation
$s=t$ is used here, in the proof of the major premise, and not as a
premise of the discharge. Apply \rulename{Export} with a fresh $b$. By
\eqref{disch:eq:rows}, $\Pi(b,0)\equiv H$ and $\Pi(b,1)\equiv(b=t)$, so
the test succeeds and $\rho_\Pi(b)\equiv t$. In the branch, \rulename{Use}
of $\mathcal C_\Pi(b)$ at $2$ with the premise $H$ gives
$F\equiv(h(b)=r)$. The exported formula is
$(h(b)=r)\subst{t}{b}\equiv(h(t)=r)$ by \eqref{disch:eq:coherence}, and it
does not contain $b$.
\end{proof}

The row that the rule must read is at the fixed position $1$; the row that
the branch supplies is read by \rulename{Use} at $2$.

\begin{theorem}[$\mathcal X$ and $\mathcal W$]
\label{disch:thm:export}
$\Gamma\proves_{\mathcal X}F$ if and only if $\Gamma\proves_{\mathcal W}F$,
for every context $\Gamma$ and formula $F$, by effective transformations
that leave the end sequent unchanged.
\end{theorem}

\begin{proof}
From $\mathcal X$ to $\mathcal W$ use the last statement of
\cref{disch:prop:inferred-sound} at every \rulename{Export} node. From
$\mathcal W$ to $\mathcal X$: by
\cref{disch:thm:export-transport,disch:lem:equality,disch:lem:formula,disch:lem:export-elim},
$\mathcal X$ derives transport in every formula template. At a
\rulename{Name} node the translated branches give $b=t$ and $F$ under the
opened clause; transport in the template $F\subst{w}{b}$ gives
$F\subst{t}{b}$; by the side condition of \rulename{Name} it does not
contain $b$, and \cref{disch:lem:export-elim} discharges the witness. This
covers a term $t$ that contains $b$. The order of the constructions is:
ordinary elimination, numerical cases, the packet, transport in a term,
transport in a formula; none uses a later one, and no substitution into a
derivation is used.
\end{proof}

\subsection{The open form}

In $\mathcal O$ ordinary elimination is not immediate: the branch of
\rulename{Open} assumes $\mathcal C_B(\rho_B(b))$ where \rulename{Elim}
assumes $\mathcal C_B(b)$, and it does not help that the goal is free of
$b$. The assumption is changed by substituting in the branch derivation.
The order of the argument avoids a circle: the instances of
\rulename{Open} on which the selector necessarily falls back give
numerical cases and the unrestricted case introduction; these give
substitution; substitution gives ordinary elimination.

\begin{lemma}[Fallback openings]
\label{disch:lem:fallback}
If the witness coordinate does not occur in $B$, or if $B$ has root $D$,
then \rulename{Open} for $B$ is literally \rulename{Elim} for $B$.
Consequently
\cref{disch:lem:closure,disch:lem:repair,disch:lem:packet} hold in
$\mathcal O$.
\end{lemma}

\begin{proof}
\Cref{disch:lem:selector}(b). The three lemmas require \rulename{Elim}
only for these two kinds of matrices, and weakening and cut hold by
\cref{disch:lem:structural}.
\end{proof}

\begin{theorem}[Substitution in $\mathcal O$]
\label{disch:thm:subst}
If $\Gamma\proves_{\mathcal O}A$ and $\sigma$ is a simultaneous
substitution of terms for parameters, then
$\Gamma\sigma\proves_{\mathcal O}A\sigma$. The transformation of
derivations is effective. It does not preserve the rule instance at each
node.
\end{theorem}

\begin{proof}
By induction on the given derivation, with the statement quantified over
all substitutions. The induction hypothesis is applied to proper
subderivations of the given derivation only, never to a derivation that
the construction has produced. A substitution may send a parameter to a
term that contains that parameter; images are inserted once.

\rulename{Refl} is immediate. For \rulename{Use} of an occurrence
$[A';\,y.\,B]$ at $t$, the context $\Gamma\sigma$ contains
$[A'\sigma;\,y.\,B\sigma]$, the induction hypothesis gives $A'\sigma$, and
\rulename{Use} at $t\sigma$ gives $(B\sigma)(t\sigma)\equiv B(t)\sigma$ by
\eqref{disch:eq:coherence}.

For \rulename{Intro} at the witness $t$ with the matrix $B$, rename $z$
away from the images of $\sigma$ and let $B'\equiv B\sigma$,
$t'\equiv t\sigma$. The induction hypothesis gives
$\Gamma\sigma,B'(t',\eta)\proves B'(t',\Suc z)$ with the old $\eta$. A
progress instance remains one: the root is still an equation and the
literal $0$ is still $0$. A case instance remains one unless $B$ is a
literal equation and $t\not\equiv 0\equiv t'$. In that event the
substituted premise is the premise of the unrestricted case introduction
of \cref{disch:lem:repair}, which holds in $\mathcal O$ by
\cref{disch:lem:fallback} and was proved before the present theorem.

For \rulename{Open} with the major premise $\D{x}{y}B$ and the parameter
$b$, rename $b$ away from the images of $\sigma$, let $B'\equiv B\sigma$
and $\rho'\equiv\rho_{B'}(b)$, and translate the major premise by the
induction hypothesis. There are two cases, according to the test for the
\emph{old} matrix.

If the old test succeeds with $\rho_B(b)\equiv r$, then by
\cref{disch:lem:selector}(c) the new test succeeds and
$\rho'\equiv r\sigma$. Apply the induction hypothesis to the branch with
$\sigma$, which fixes $b$. Its clause becomes
\[
 \mathcal C_B(r)\sigma\equiv[B'(r\sigma,0);\,y.\,B'(r\sigma,y)]
 \equiv\mathcal C_{B'}(\rho')
\]
by \eqref{disch:eq:coherence}; this does not need $r$ to be free of $b$.

If the old selector falls back, the old branch has the clause
$\mathcal C_B(b)$. The new test may fail or succeed. In both events apply
the induction hypothesis to the old branch with the substitution that
agrees with $\sigma$ on the other parameters and sends $b$ to $\rho'$.
The clause becomes $\mathcal C_{B'}(\rho')$, and the context and the goal
become $\Gamma\sigma$ and $A\sigma$, because $\Gamma$ and $A$ do not
contain $b$ and no image of $\sigma$ does.

In either case \rulename{Open} applies to the translated premises and
proves $A\sigma$.
\end{proof}

The second case is necessary: by \cref{disch:lem:selector}(d) the selector
is not an invariant test, and an argument that treats it as one leaves a
gap. The checkers of \pkg{inferred\_witness} and
\pkg{deterministic\_opening} are tested on a derivation whose opening
falls back and for which a substitution of a passive parameter makes the
test succeed: the old branch is then rejected, and the old branch
instantiated at the newly selected witness is accepted.

\begin{corollary}[Ordinary elimination in $\mathcal O$]
\label{disch:cor:open-elim}
\rulename{Elim} is admissible in $\mathcal O$ for every matrix.
\end{corollary}

\begin{proof}
Given $\Gamma,\mathcal C_B(b)\proves A$ with $b$ fresh for $\Gamma$ and
$A$, apply \cref{disch:thm:subst} with the substitution
$b\mapsto\rho_B(b)$ that fixes all other parameters. The result is
$\Gamma,\mathcal C_B(\rho_B(b))\proves A$, the branch premise of
\rulename{Open}.
\end{proof}

\begin{theorem}[Transport in $\mathcal O$]
\label{disch:thm:open-transport}
\rulename{UTr} is derived in $\mathcal O$ for every term template $h(w)$.
\end{theorem}

\begin{proof}
Take the packet of \cref{disch:thm:export-transport}, with the rows
$x=t$ and $h(x)=r$; by \cref{disch:lem:packet,disch:lem:fallback},
$\Gamma\proves\D{x}{y}\Pi$. Apply \rulename{Open} with a fresh $b$. The
test succeeds with $\rho_\Pi(b)\equiv t$, so the branch has the clause
$\mathcal C_\Pi(t)=[H;\,y.\,\Pi(t,y)]$. \rulename{Use} at $2$ with the
premise $H$ gives $\Pi(t,2)\equiv(h(t)=r)$, the unchanged outer goal,
which does not contain $b$. There is no equality premise and no
substitution in an exported formula; the derivation does not use
\cref{disch:thm:subst}.
\end{proof}

\begin{theorem}[$\mathcal O$ and the reference calculus]
\label{disch:thm:open}
$\Gamma\proves_{\mathcal O}A$ if and only if $\Gamma\proves_{\mathcal T}A$,
for every context $\Gamma$ and formula $A$, by effective transformations
that leave the end sequent unchanged.
\end{theorem}

\begin{proof}
From $\mathcal T$ to $\mathcal O$: \rulename{Elim} is replaced by
\cref{disch:cor:open-elim}, and \rulename{Tr} by
\cref{disch:thm:open-transport} followed by symmetry
(\cref{disch:lem:equality}).

From $\mathcal O$ to $\mathcal T$: $\mathcal T$ derives transport in every
formula template. At an \rulename{Open} node translate the major premise
and open it by \rulename{Elim} at $b$, keeping the clause
$\mathcal C_B(b)$. Write $\rho\equiv\rho_B(b)$;
\cref{disch:lem:selected} gives $b=\rho$. Derive $q_yB(\rho,y)$ by the
first closure rule of \cref{disch:lem:closure}: from its assumption
$B(\rho,0)$ transport along $\rho=b$ in the template $B(w,0)$ to $B(b,0)$,
use $\mathcal C_B(b)$ at $\Suc z$, and transport along $b=\rho$ in the
template $B(w,\Suc z)$ to $B(\rho,\Suc z)$. \rulename{Elim} for
$q_yB(\rho,y)$, whose witness coordinate is unused, opens exactly the
clause $[B(\rho,0);\,y.\,B(\rho,y)]=\mathcal C_B(\rho)$. The translated
branch, weakened by the clause $\mathcal C_B(b)$ that is still present,
proves $A$. Both eliminations close, since $A$ contains neither the unused
witness nor $b$. The construction is valid when $\rho$ contains $b$:
throughout the inner part $b$ is an available parameter, and only the
outer elimination discharges it.
\end{proof}

\subsection{Summary and comparison of the four packages}

\begin{theorem}[The inferred witness]
\label{disch:thm:inferred}
The calculi $\mathcal T$, $\mathcal W$, $\mathcal X$ and $\mathcal O$
derive the same sequents. For every context $\Delta$ of clauses with at
most one antecedent and every formula $F$ the following are equivalent:
$L(\Delta)\proves_{\mathcal T}F$;
$\mathcal R,\Delta\proves_{\mathcal W^{\mathcal R}}F$;
$\mathcal R,\Delta\proves_{\mathcal X^{\mathcal R}}F$;
$\mathcal R,\Delta\proves_{\mathcal O^{\mathcal R}}F$. Substitution of
terms for parameters is admissible in each of them. $\Nat$ derives
$\Gamma\proves A$ if and only if $\mathcal O^{\mathcal R}$, equivalently
$\mathcal X^{\mathcal R}$, derives $A^{\flat}$ from $\mathcal R$ and the
unconditional assumption clauses of $\Gamma^{\flat}$; and the ordinary
readings of $\mathcal O^{\mathcal R}$ and $\mathcal X^{\mathcal R}$ are
bi-interpretable with $\iPRA$ as in \cref{disch:cor:name-three}.
\end{theorem}

\begin{proof}
\Cref{disch:thm:name,disch:thm:export,disch:thm:open} give the first
statement, \cref{disch:lem:conventions} the second.
\Cref{disch:thm:subst} gives substitution in $\mathcal O$, and the other
calculi inherit it through the equivalences, which leave end sequents
unchanged. The rest is \cref{disch:prop:kernel} and
\cref{base:thm:bridge}.
\end{proof}

In Lean the rule on the kernel, with the selected witness read off the
normalized matrix in place of \rulename{DE}, is \Lean{FOPrf} with
\Lean{foprf\_iff\_nprf} (\Lean{NativeInferredWitness}); its completeness
proof translates a kernel derivation under an arbitrary substitution and
at \rulename{DE} sends the fresh variable to the selected witness whatever
it is, as in the second case of the proof of \cref{disch:thm:subst}; no
distinction of cases on the test is needed there. The clause form
$\mathcal O^{\mathcal R}$ is \Lean{OCPrf} with \Lean{ocprf\_iff\_nprf}
(\Lean{NativeOpenClauses}), and $\mathcal X^{\mathcal R}$ is \Lean{XCPrf}
with \Lean{xcprf\_iff\_nprf} (\Lean{NativeExportClauses}).

\begin{remark}[How the four packages differ]
\label{disch:rem:four}
\emph{Rows inspected.} There is no difference: all four inspect the rows
$B(b,0)$ and $B(b,1)$ of the normalized matrix, require the first to be
literally $0=0$ and the second to be literally an equation whose left side
is the fresh variable, and take the right side of the second as the
witness.

\emph{The reflexivity clause.} There is no difference either: all four are
stated in the convention with the fixed unconditional clause
$[\ ;\,y.\,y=y]$ and clauses with at most one antecedent, and have three
rule families.

\emph{Use of the witness.} \pkg{inferred\_witness} and
\pkg{deterministic\_opening} state \rulename{Open}: the clause is opened
at the selected witness and the branch proves the unchanged goal. Their
proofs are those of
\cref{disch:lem:fallback,disch:thm:subst,disch:cor:open-elim,disch:thm:open-transport,disch:thm:open}.
\pkg{first\_instance} and \pkg{clause\_discharge} state \rulename{Export}:
the clause is opened at the fresh variable and the witness is substituted
in the exported formula. Their proofs are those of
\cref{disch:lem:export-elim,disch:thm:export-transport,disch:thm:export},
and they need no substitution theorem. \pkg{inferred\_witness} keeps the
export form as a second write-up and as a control checker, and
\pkg{deterministic\_opening} keeps it as an intermediate variant with its
own checker.

\emph{Reference calculus.} \pkg{inferred\_witness} proves the equivalence
with $\mathcal T$ in the convention with the reflexivity clause; the other
three prove it with $\mathcal W^{\mathcal R}$. In the translation of
\rulename{Open} back to the reference calculus, \pkg{inferred\_witness}
transports the base of the closure in both directions, as in the proof of
\cref{disch:thm:open}; \pkg{deterministic\_opening} distinguishes the two
cases of the test and, when it succeeds, uses that the base is $H$ and
ignores the assumption of the closure.

\emph{Additional material.} \pkg{first\_instance} compares its calculus
with the raw clause calculus directly instead of through $\mathcal W$, by
the argument cited in \cref{disch:prop:kernel}, and contributes
\cref{disch:rem:strict}. \pkg{clause\_discharge} delivers, besides the
checker of $\mathcal X^{\mathcal R}$, a smaller checker of the unchanged
$\mathcal W^{\mathcal R}$. \pkg{deterministic\_opening} delivers a
reordering of an inherited compact checker that is independent of the new
rule (\cref{disch:sec:impl}).

\emph{Names.} The rule is called $\mathrm O^*$ in \pkg{inferred\_witness},
$\mathrm O$ in \pkg{deterministic\_opening}, Name in
\pkg{first\_instance} and Row in \pkg{clause\_discharge}; the Lean
development calls \rulename{Export} Discharge.
\end{remark}

\begin{remark}[The conditions of the test]
\label{disch:rem:conditions}
Three conditions cannot be removed. \pkg{deterministic\_opening} exhibits,
for each, a certificate that the correspondingly weakened checker accepts:
without the requirement that the left side of the row at $1$ is the fresh
witness, and without the requirement that the row at $0$ is $0=0$, the
checker accepts $1=0$; without the requirement that the outer goal is in
the outer scope, it accepts a formula in which the fresh parameter
escapes. In \pkg{clause\_discharge} the source of the replacement must be
the witness opened by the same inference: an equation whose left side is
an older parameter, a constant or a compound term triggers no replacement.
With an older parameter as source, the argument of
\cref{disch:lem:export-elim} fails, since the substitution then acts on a
goal that mentions that parameter; the generalization is not part of any
of the calculi (\cref{disch:sec:further}).
\end{remark}

\section{Substitutive discharge on the canonical recursor syntax}
\label{disch:sec:subst}

This section draws on \pkg{substitutive\_discharge}. Its term syntax, and
the passage from that syntax to iterator terms, are those of
\pkg{canonical\_recursors}.

The calculi of \cref{disch:sec:name,disch:sec:inferred} depend on full
normalization: the rows of a packet must be literal before any equality
rule is available. If the discharge rule is given a formula template, so
that the discharged equation may have an arbitrary source, no packet is
needed, the only normalization is a structural erasure, and reflexivity is
derived from one computation schema.

\subsection{Syntax}

In this section the convention on normal forms is not in force. The terms
are those of the seeded recursor $R(n,a;u.h)$ with one bound variable of
\cref{ctr:def:languages}, with the defined successor
$\Suc a\defeq R(0,a;u.0)$ and the computation equations
\eqref{ctr:eq:r1zero} and \eqref{ctr:eq:r1succ}. The zero computation is
made part of the syntax: terms and formulas of the calculus are zero-body
canonical, and substitution is the normalized substitution
$T\langle\sigma\rangle$, both in the sense of \cref{ctr:def:nu}. The
canonicalization $\nu$ does not unfold a successor iteration, does not
execute a body and does not test whether a term denotes zero; it preserves
the root constructor \eqref{ctr:eq:nuroot}, so that $\nu(t)\equiv 0$ if
and only if $t\equiv 0$; and it is coherent with substitution, which is
therefore associative, commutes with renaming and introduces no free
variables beyond those of its arguments (\cref{ctr:lem:coherence}). The
syntax is that of the calculi $\mathcal C_\nu$ and $\mathcal K_6$ of
\cref{ctr:def:free,ctr:def:six}.

Formulas are equations and $\D{x}{y}B$. A context is a finite list of
clauses $[A;\,y.\,B(y)]$ with one antecedent, and an assumption $A$ is the
clause $[0=0;\,v.\,A]$. The truth of $0=0$ is derived, not assumed, and no
clause is fixed in advance. The contradiction is $0=\Suc 0$.

\subsection{Four rule families}

\begin{definition}[The calculus of substitutive discharge]
\label{disch:def:subdis}
$\mathcal K$ has the following four rules. All displayed expressions are
canonical.
\begin{align}
 &\infrule{}{\Gamma\proves R(\Suc n,a;\,u.\,h)=h\subst{R(n,a;\,u.\,h)}{u}}\
 \rulename{Calc}
 \label{disch:eq:calc}\\[1ex]
 &\infrule{\Gamma\proves A}{\Gamma\proves B(t)}\ \rulename{Use},
 \qquad [A;\,y.\,B(y)]\text{ an occurrence in }\Gamma
 \notag\\[1ex]
 &\infrule{\Gamma,\ B(t,\eta)\proves B(t,\Suc z)}
          {\Gamma\proves\D{x}{y}B(x,y)}\ \rulename{Intro},
 \qquad
 \eta=\begin{cases}
 z,&B\text{ is a literal equation and }t\equiv 0,\\
 0,&\text{otherwise}
 \end{cases}
 \notag\\[1ex]
 &\infrule{\Gamma\proves\D{x}{y}B\qquad
           \Gamma,\ \mathcal C_B(b)\proves s=t\qquad
           \Gamma,\ \mathcal C_B(b)\proves F(s)}
          {\Gamma\proves F(t)}\ \rulename{SDis}
 \label{disch:eq:sdis}
\end{align}
\rulename{Calc} is the schema \eqref{ctr:eq:Rplus}: the predecessor iterate
on the right is canonicalized, and at $n\equiv 0$ it is $\Suc a$. In \rulename{Intro}, $z$ is fresh for
$\Gamma$, $t$ and $\D{x}{y}B$, and the test is literal. In \rulename{SDis},
$b$ is fresh for $\Gamma$ and $\D{x}{y}B$, the exported formula $F(t)$
does not contain $b$, and $F(v)$ is an arbitrary formula template whose
distinguished variable $v$ is different from $b$; a constant template is
allowed, and $s$, $t$ and the template may contain $b$. The calculus
$\mathcal K_{\mathrm{sep}}$ has \rulename{Calc}, \rulename{Use},
\rulename{Intro}, \rulename{Elim} and the transport rule
\begin{equation}
 \infrule{\Gamma\proves s=t\qquad\Gamma\proves F(s)}{\Gamma\proves F(t)}
 \ \rulename{TrF}
 \label{disch:eq:trf}
\end{equation}
for every formula template $F$.
\end{definition}

$\mathcal K$ has no hypothesis rule, no reflexivity axiom, no separate
equality rule, no separate elimination of $\D{x}{y}$, no coverage or
separation axiom and no structural rule. The one computation schema is
\rulename{Calc}; the progress mode of \rulename{Intro} carries the
induction. Weakening, exchange, contraction, injective renaming and cut of
an assumption are admissible by induction on derivations, as in
\cref{disch:lem:structural,ctr:lem:struct}; injective renaming preserves the literal test
of \rulename{Intro}. Substitution of a term for a parameter can turn a
witness into the literal $0$; it is admissible through
\cref{disch:prop:sd-basis}, not by an induction that preserves the rules.

\subsection{The bootstrap of equality}

Let, for a term $t$ and variables $u,x,y$ that do not occur in it,
\[
 W_t\defeq R(1,0;\,u.\,t),\qquad J\defeq(W_0=0),\qquad
 \Theta\defeq\D{x}{y}J .
\]
\rulename{Calc} proves $W_t=t$: the body is constant and ignores the
predecessor iterate.

\begin{lemma}[A closed paired theorem]
\label{disch:lem:theta}
$\Theta$ is derived by \rulename{Calc} and \rulename{Intro} alone.
\end{lemma}

\begin{proof}
Apply \rulename{Intro} at the witness $1$. The matrix does not depend on
its bound variables, so the goal of the premise is $J$, which
\rulename{Calc} proves without the assumption. Neither the use of an
assumption nor reflexivity is involved.
\end{proof}

\begin{lemma}[Transport in a formula]
\label{disch:lem:sd-transport}
\rulename{TrF} is a derived rule of $\mathcal K$.
\end{lemma}

\begin{proof}
Use $\Theta$ as the major premise of \rulename{SDis}, and weaken the two
given derivations into the opened context; they do not mention the fresh
witness, and $F(t)$ does not contain it. No law of equality is used.
\end{proof}

\begin{lemma}[Reflexivity and assumptions]
\label{disch:lem:sd-refl}
$\Gamma\proves t=t$ for every term $t$, and an assumption can be used:
$\Gamma,A\proves A$.
\end{lemma}

\begin{proof}
\rulename{Calc} gives $W_t=t$. Use it as both premises of \rulename{TrF}
with the template $F(v)=(v=t)$: the instance at the source is $W_t=t$ and
the instance at the target is $t=t$. For an assumption $A$, take $t\equiv
0$ and apply \rulename{Use} to $[0=0;\,v.\,A]$.
\end{proof}

\begin{lemma}[Equality and ordinary elimination]
\label{disch:lem:sd-eq}
Symmetry, transitivity, the congruence $s=t\proves h(s)=h(t)$ and
\rulename{Elim} are derived rules of $\mathcal K$.
\end{lemma}

\begin{proof}
From $s=t$, transport the reflexive $s=s$ in the template $(v=s)$ to
$t=s$. From $s=t$ and $t=r$, transport $t=r$ along $t=s$ in the template
$(v=r)$ to $s=r$. From $s=t$, transport the reflexive $h(s)=h(s)$ in the
template $(h(s)=h(v))$ to $h(s)=h(t)$. For \rulename{Elim}, let the opened
branch prove a formula $A$ without $b$. Apply \rulename{SDis} with the
constant template $F(v)=A$ and with the equation $J$, proved by
\rulename{Calc}, as the equality premise; the third premise is the given
branch, and the exported formula is $A$. No packet and no normalization of
a selector is needed.
\end{proof}

\begin{theorem}[$\mathcal K$ and the calculus with separate rules]
\label{disch:thm:sd}
For every context $\Gamma$ of canonical clauses and every formula $A$,
$\Gamma\proves_{\mathcal K}A$ if and only if
$\Gamma\proves_{\mathcal K_{\mathrm{sep}}}A$. The translations are
effective and preserve the end sequent, including its parameters.
\end{theorem}

\begin{proof}
An instance of \rulename{SDis} is \rulename{TrF} inside the opened context
followed by \rulename{Elim}, whose side condition is that of
\rulename{SDis}. Conversely \rulename{TrF} is
\cref{disch:lem:sd-transport} and \rulename{Elim} is
\cref{disch:lem:sd-eq}. The other rules are shared. Induct on derivations,
renaming eigenparameters when the fixed auxiliary derivations are
inserted.
\end{proof}

\subsection{The canonical basis and the bridge}

\begin{proposition}[Recovered rules]
\label{disch:prop:sd-basis}
The following are derived in $\mathcal K$, with $q_yB\defeq\D{w}{y}B$ for
an unused $w$.
\begin{enumerate}[label=\textup{(\alph*)}]
\item Case introduction for every matrix and witness: from
 $\Gamma,B(t,0)\proves B(t,\Suc z)$ infer $\Gamma\proves\D{x}{y}B$.
\item From $\Gamma,B(0)\proves B(\Suc z)$ infer $\Gamma\proves q_yB$; from
 $\Gamma\proves q_yB$ and $\Gamma\proves B(0)$ infer $\Gamma\proves B(t)$;
 and, for a literal equation $B$, from $\Gamma,B(z)\proves B(\Suc z)$
 infer $\Gamma\proves q_yB$. Hence induction for literal equations under
 an arbitrary context and numerical cases for arbitrary formulas.
\item A context that contains the clause $[B(0);\,y.\,B(y)]$ proves
 $q_yB$.
\end{enumerate}
\end{proposition}

\begin{proof}
(a) Only a literal equation at the witness $0$ is missing from
\rulename{Intro}. The term $c\defeq W_0$ is canonical, is not the literal
$0$, and satisfies $c=0$ by \rulename{Calc}. Given the premise at the
witness $0$, start from the assumption $B(c,0)$, transport it to $B(0,0)$,
insert the given derivation by cut, and transport its conclusion
$B(0,\Suc z)$ back to $B(c,\Suc z)$. \rulename{Intro} at $c$ is in case
mode and concludes $\D{x}{y}B$. No induction is used.

(b) The first rule is (a) for the unused coordinate. For the second,
\rulename{Elim} opens $[B(0);\,y.\,B(y)]$, and \rulename{Use} with the
proof of $B(0)$ gives $B(t)$. The third is \rulename{Intro} at the literal
witness $0$ of the unused coordinate, in progress mode.

(c) Use (a) at the unused witness $1$: under its assumption $B(0)$, use
the clause at $\Suc z$.
\end{proof}

The term $W_0$ is the term $c=\operatorname{wr}(0)$ of \cref{ctr:lem:sixeq},
and (a) is the repair of \cref{ctr:lem:sixcase}, with \rulename{TrF} in
place of the equality of $\mathcal K_6$; (b) and (c) are the clause forms
of \cref{ctr:lem:resources} and of \cref{claus:lem:closure}(i)--(iv).

The Eq/$Q$/$D$ basis of \pkg{canonical\_recursors} is the six-rule calculus
$\mathcal K_6$ of \cref{ctr:def:six}. By (c) a derivation of $\mathcal K_6$
whose elimination of $\D{x}{y}$ keeps a formula $\Q{y}B(b,y)$ is
translated by deriving $q_yB(b,y)$ from the opened clause and cutting it
into the translated branch; the two representations of a context are
related by derivations and are not identified by convention. Reflexivity,
symmetry and \rulename{TrF} give the crossed equation transport (X) of
$\mathcal K_6$, including its restriction on the source, and the crossed
elimination (QX) is the second rule of (b) followed by symmetry when the
conclusion is an equation. Conversely $\mathcal K_6$ derives formula
transport (\cref{ctr:lem:sixtransport}) and the elimination of $\D{x}{y}$,
which validate \rulename{SDis}, a clause $[A;\,y.\,B]$ being represented by
its ordinary reading. So $\mathcal K$ and $\mathcal K_6$ have the same
consequences on formulas, and by \cref{ctr:thm:six,ctr:thm:canonical} the
same as $\mathcal C_\nu$ and $\mathcal R_{\mathrm{raw}}$.

\begin{theorem}[Bridge for $\mathcal K$]
\label{disch:thm:sd-ipra}
The ordinary reading of $\mathcal K$ and $\iPRA$ are bi-interpretable with
whole domains, actual equality, both numerical composites with graph
$x=y$, and both formula round trips. In particular they are equiconsistent.
\end{theorem}

\begin{proof}
By the comparison above, $\mathcal K$ and $\mathcal K_6$ derive the same
formulas from the same canonical contexts, so the statement is
\cref{ctr:thm:sixbridge}, whose bi-interpretation is that of
\cref{ctr:thm:transfer} for $\mathsf T_R^1$. The direct interpretations, in
the form described in \cref{ctr:rem:direct}, read as follows on the rules
of $\mathcal K$.

From $\mathcal K$ to $\iPRA$: interpret each fixed term $R(n,a;u.h)$ by
primitive recursion with the seed $\Suc a$ and the displayed step. The
erasure of a body at the counter $0$ preserves the value. \rulename{Calc}
is a defining equation, \rulename{Use} is implication and universal
elimination, \rulename{Intro} is coverage in case mode and induction for
an equation under its assumed base in progress mode, and \rulename{SDis}
is equality substitution inside the scope of an existential elimination.

From $\iPRA$ to $\mathcal K$: the predecessor $p$ of \eqref{ctr:eq:seedpred}
(\cref{seed:thm:succ-pred}, written $P$ there) is constructed before any
numerical selector, cancellation of successor or reflection of equality is
available, and the iterator of \cref{base:sec:terms} is the initialized
iteration $\widehat I(n,a;u.h)=p(R(n,a;v.\Suc(h\subst{p(v)}{u})))$ of
\cref{seed:thm:succ-iter} and \cref{ctr:def:translations}: at the counter
$0$ its value is $p(\Suc a)=a$, and at a successor \rulename{Calc} exposes
an outer successor that $p$ removes. By
\cref{disch:prop:sd-basis,disch:lem:sd-refl,disch:lem:sd-eq} the rules
needed by the logical macros and by the arithmetic of
\cref{base:thm:pred,base:thm:reflect,base:thm:primrec} are available in
$\mathcal K$, with $\sel(v,a,b)$ read as $\widehat I(v,a;u.b)$. This
gives the defining equations and the quantifier-free induction of $\iPRA$.

Both numerical composites are the identity by the argument of
\cref{ctr:thm:transfer}(2): recurrence uniqueness (\cref{ctr:lem:unique}),
applied to each fixed description in the order of its dependencies, with
pointwise equality of bodies lifted to the iteration by body
extensionality (\cref{ctr:lem:bodyext}) and not assumed as a congruence;
the interpretation of successor returns $\Suc a$. The logical round trips
of \cref{base:thm:bridge} identify the reading of $\D{x}{y}$ with its
returned macro, and induction on formulas extends the numerical
identities to both formula composites. The contradiction $0=\Suc 0$
corresponds to falsum in both directions.
\end{proof}

\begin{remark}[Checking, the raw companion, Lean]
\label{disch:rem:sd}
The syntax recursion and the proof recursion of $\mathcal K$ are
structural, and the only canonicalization erases a body at a literal
counter $0$. Under a standard coding of finite trees the certificate
relation is therefore primitive recursive, with no fuel field, in contrast
to \cref{disch:rem:fuel}.

The package contains a companion on raw iterator terms without any
normalization: the two computation equations of \cref{base:sec:terms}, the
use of a clause, \rulename{Intro}, and the same discharge rule. An
assumption $F$ is $[H_0;\,v.\,F]$ with
$H_0\defeq(\It{0}{0}{u}{0}=0)$, one instance of the computation schema;
$\D{x}{y}H_0$ is introduced at the witness $\Suc 0$ without its
assumption; transport is the discharge of this formula; reflexivity is the
transport of $\It{0}{t}{u}{0}=t$ along itself in the template $(v=t)$; and
\rulename{Elim} is the discharge with a constant template and the equation
$H_0$. The proofs are those of
\cref{disch:lem:theta,disch:lem:sd-transport,disch:lem:sd-refl,disch:lem:sd-eq}.

In Lean the discharge rule on the kernel, in place of \rulename{Tr} and
\rulename{DE}, gives the seven-rule calculus \Lean{SDPrf}, with
\Lean{sdprf\_iff\_nprf} (\Lean{NativeSubstitutiveDischarge}): the closed
theorem is $\D{x}{y}(\It{0}{0}{u}{u}=0)$, proved by cases from a
computation; a discharge of it that does not use the opened context is
transport in an arbitrary formula (\Lean{SDPrf.tr}), and a constant
template gives \rulename{DE} (\Lean{SDPrf.dE}). The four families on the
canonical recursor syntax are not formalized.
\end{remark}

\section{Restricted substitution}
\label{disch:sec:restricted}

This section draws on \pkg{variable\_substitution}
(\cref{disch:sec:vsub,disch:sec:newest}), \pkg{parameter\_substitution}
(\cref{disch:sec:param}) and \pkg{substitutive\_discharge}
(\cref{disch:sec:unary}).

\subsection{Substitution for a variable}
\label{disch:sec:vsub}

The convention on normal forms of \cref{disch:sec:setting} is in force.
The equality rule substitutes for a parameter and does nothing else:
\begin{equation}
 \infrule{\Gamma\proves x=t\qquad\Gamma\proves u=v}
         {\Gamma\proves v=u\subst{t}{x}}\ \rulename{VSub},
 \qquad x\text{ a parameter.}
 \label{disch:eq:vsub}
\end{equation}
The rule substitutes simultaneously for every occurrence of $x$ in $u$. It
does not remove $x$ from the parameters and does not substitute in
$\Gamma$ or in $v$. The term $t$ may contain $x$. No template, no selected
occurrence, no compound source and no matching of a pattern is supplied.
A variable bound by an iterator is not a source.

\begin{definition}
\label{disch:def:vsub}
$\mathcal V$ is the calculus with the five rules \rulename{Refl},
\rulename{VSub}, \rulename{Use}, \rulename{Elim} and \rulename{Intro}, in
the convention with \rulename{Refl}.
\end{definition}

\Cref{disch:lem:structural,disch:lem:closure,disch:lem:packet} apply to
$\mathcal V$, in which \rulename{Elim} is primitive. In particular the
numerical cases are derived with no equality rule, and neither the case
introduction at the witness $0$ nor an envelope is needed for them.

\begin{theorem}[Transport in $\mathcal V$]
\label{disch:thm:vsub-transport}
For every term template $h(w)$ the rule \rulename{Tr} is derived in
$\mathcal V$. The derivation contains exactly one instance of
\rulename{VSub}.
\end{theorem}

\begin{proof}
Let $\Gamma\proves s=t$ and $\Gamma\proves h(s)=r$, and take the packet
$\Pi$ of \cref{disch:thm:name-transport}, with the rows $h(x)=r$ and
$x=t$. By \cref{disch:lem:packet}, $\Gamma\proves\D{x}{y}\Pi$. Open it by
\rulename{Elim} at a fresh $b$; in the branch the same lemma gives
$b=t$ and $h(b)=r$. One \rulename{VSub} with the major premise $b=t$ and
the minor premise $h(b)=r$ yields
\[
 r=(h(b))\subst{t}{b}\equiv h(t)
\]
by \eqref{disch:eq:coherence}. The conclusion does not contain $b$, and
\rulename{Elim} closes.
\end{proof}

Everything before the one \rulename{VSub} uses the logical rules,
reflexivity and literal computation only: no symmetry, transitivity or
congruence, no substitution into a derivation, no reflection of equality
and no induction. Occurrences of $s$ in the passive part of $h$ are not
touched, because $b$ is fresh for $h$; normalization may duplicate or
erase the distinguished occurrences, and \eqref{disch:eq:coherence}
accounts for this without counting them.

\begin{theorem}[$\mathcal V$ and the reference calculus]
\label{disch:thm:vsub}
$\Gamma\proves_{\mathcal V}A$ if and only if $\Gamma\proves_{\mathcal T}A$,
for every context $\Gamma$ and formula $A$, by effective transformations
that are the identity on the end sequent.
\end{theorem}

\begin{proof}
A \rulename{VSub} node is one \rulename{Tr}: with $w$ fresh let
$K(w)\defeq u\subst{w}{x}$; then $K(x)\equiv u$ and
$K(t)\equiv u\subst{t}{x}$, also when $t$ contains $x$, and \rulename{Tr}
from $x=t$ and $K(x)=v$ concludes $v=K(t)$. Conversely a \rulename{Tr}
node is replaced by the derivation of \cref{disch:thm:vsub-transport}
after its premises are translated. The fresh names of that derivation are
chosen away from the context, the conclusion and the terms involved, so
all conditions of \rulename{Elim} and \rulename{Intro} hold.
\end{proof}

Together with \cref{disch:prop:kernel} and \cref{base:thm:bridge} this
gives the comparison with $\Nat$ and the bridge to $\iPRA$, as in
\cref{disch:cor:name-three}.

\begin{proposition}[The native form with six rules]
\label{disch:prop:vsub-native}
On normal forms with the formulas $s=t$, $\Q{y}A$, $\D{x}{y}A$ and
contexts of formulas, the calculus with \rulename{Hyp}, \rulename{Refl},
\rulename{VSub}, \rulename{QE}, \rulename{DE} and the introduction
\rulename{Intro} of \cref{disch:def:rules} derives the same sequents as
the calculus with \rulename{Tr} in place of \rulename{VSub}, under the
identity translation.
\end{proposition}

\begin{proof}
\rulename{VSub} is one \rulename{Tr} as before. For the converse the
numerical cases use no equality rule: introduce $\D{x}{y}F(y)$, with $x$
unused, at the witness $1$, ignoring the assumption $F(0)$; the branch of
\rulename{DE} has the hypothesis $\Q{y}F(y)$, and \rulename{QE} with the
given proof of $F(0)$ gives $F(t)$. The packet is introduced and
eliminated as in \cref{disch:thm:vsub-transport}; the two uses of the
clause $[H;\,y.\,\Pi(b,y)]$ become two applications of \rulename{QE} to
the hypothesis $\Q{y}\Pi(b,y)$ with the reflexive premise $H$, and the
final \rulename{VSub} is unchanged.
\end{proof}

The calculus with \rulename{Tr} in this proposition is the six-rule
calculus of \pkg{full\_normalization}. In Lean the kernel on full normal
forms with \rulename{VSub} in place of \rulename{Tr} is \Lean{VSPrf}, with
\Lean{vsprf\_iff\_nprf} (\Lean{NativeVariableSubstitution}) and the packet
\Lean{VSPrf.packet}.

\begin{remark}[What is reduced, and what the proof needs]
The source of the equality rule is a variable, and the action of the rule
is the substitution that the binders already require. The freedom of an
ordinary equality rule is recovered by introducing, using and eliminating
a fresh name. Full normalization is what makes the argument short: the
rows of the packet at $0$, $1$ and $2$ are literal before any congruence
is available. The proof uses the logical rules; no completeness statement
is made for the equality rules in isolation.
\end{remark}

\subsection{The newest variable}
\label{disch:sec:newest}

In this subsection the parameters of a sequent form a finite ordered list,
and the eigenparameter of \rulename{Elim} or \rulename{Intro} is put at
the front of the list of the premise it belongs to. $\mathcal V_0$ is
$\mathcal V$ with the source of \rulename{VSub} restricted to the front
parameter of the list at its node; with de Bruijn indices the source is
the variable of index $0$.

Weakening is not admissible in $\mathcal V_0$ by an induction that
preserves the rule at each node: a parameter inserted in front turns the
former front parameter into an older one. \Cref{disch:lem:structural} is
not available, and the completeness proof uses an invariant instead.

\begin{definition}[Internally sourced derivations]
\label{disch:def:internal}
A derivation in $\mathcal V_0$ with the initial parameter list $\Xi$ is
\emph{internally sourced} if at every \rulename{VSub} node the source is
the eigenparameter of an \rulename{Elim} node of the derivation, the
\rulename{VSub} node lies in the branch of that node, and no other
\rulename{Elim} or \rulename{Intro} node lies between the two.
\end{definition}

So the source is never a member of $\Xi$, and it is the front parameter at
its node.

\begin{lemma}[Safe structural changes]
\label{disch:lem:internal}
Let an internally sourced derivation of $\Gamma\proves A$ over $\Xi$ be
given, let $\iota$ be an injection of $\Xi$ into a list $\Xi'$, not
necessarily monotone, and let $\Gamma'$ be a context over $\Xi'$ that
contains every occurrence of $\Gamma\iota$. Then there is an internally
sourced derivation of $\Gamma'\proves A\iota$ over $\Xi'$.
\end{lemma}

\begin{proof}
Rename the eigenparameters of the derivation away from $\Xi'$ and
$\Gamma'$. A node of the derivation has a sequent over a list
$(e_k,\dots,e_1,\Xi)$, where $e_1,\dots,e_k$ are the eigenparameters of
the \rulename{Elim} and \rulename{Intro} nodes on the path from the root,
and a context $\Gamma,\Lambda$, where $\Lambda$ consists of the clauses
those nodes add. Replace it by the sequent over $(e_k,\dots,e_1,\Xi')$
with the context $\Gamma',\Lambda\iota$ and the formula renamed by
$\iota$, which is extended by the identity on the $e_i$. The initial list
is not reset when the transformation descends into a premise. Every
instance of \rulename{Refl}, \rulename{Use}, \rulename{Elim} and
\rulename{Intro} is carried to an instance, as in
\cref{disch:lem:structural}: normalization and substitution commute with
injective renaming, the mode of \rulename{Intro} is preserved, and the
added clauses do not affect any side condition. At a \rulename{VSub} node
the source is $e_k$ with $k\geq 1$, and it is still the front parameter of
the new list and still bound by the same \rulename{Elim} node; inserting,
permuting or renaming members of the initial list changes neither. The
result is internally sourced over $\Xi'$.
\end{proof}

The lemma does not give weakening for arbitrary derivations of
$\mathcal V_0$.

\begin{theorem}[The newest variable suffices]
\label{disch:thm:newest}
Every derivation in $\mathcal T$ over a parameter list $\Xi$ has an
internally sourced translation in $\mathcal V_0$ with the same end sequent
and the same list. Consequently $\mathcal V_0$, $\mathcal V$ and
$\mathcal T$ derive the same sequents, and weakening, substitution of
terms for parameters and \rulename{Tr} are admissible in $\mathcal V_0$.
\end{theorem}

\begin{proof}
By induction on the derivation in $\mathcal T$. \rulename{Refl},
\rulename{Use}, \rulename{Elim} and \rulename{Intro} are copied. The
translated premises are internally sourced over their own lists; a source
that is bound by an \rulename{Elim} node inside a premise is bound by a
node inside the combined derivation, so the invariant is preserved.

At a \rulename{Tr} node apply the construction of
\cref{disch:thm:vsub-transport} to the two translated premises. The
construction places copies of them under the eigenparameters and the
assumption clauses of the \rulename{Intro} of the packet and of the
numerical-case rule; each copy is obtained by \cref{disch:lem:internal},
which applies because the premises are internally sourced by the induction
hypothesis. The remaining nodes of the construction are \rulename{Refl},
\rulename{Use}, \rulename{Elim}, \rulename{Intro} and the one
\rulename{VSub}. That node stands at the end of the branch of the
\rulename{Elim} which opens the packet at $b$; its premises $b=t$ and
$h(b)=r$ are derived from the opened clause by \rulename{Refl} and
\rulename{Use}, with no eigenparameter in between. So its source $b$ is
the front parameter and is bound by an \rulename{Elim} node of the
constructed derivation, with no node of \rulename{Elim} or
\rulename{Intro} between them.

Conversely a restricted \rulename{VSub} is an instance of \rulename{VSub}
and hence one \rulename{Tr} (\cref{disch:thm:vsub}). So the three calculi
agree. A derivation of $\mathcal V_0$ is weakened, or instantiated, or
used as a premise of \rulename{Tr}, by translating it to $\mathcal T$,
operating there, and translating back. The completeness proof itself uses
none of these operations on arbitrary derivations.
\end{proof}

The same invariant and translation apply to the six-rule form of
\cref{disch:prop:vsub-native}, with \rulename{QE} for \rulename{Use} and
\rulename{DE} for \rulename{Elim}: the final \rulename{VSub} stands under
its own newest eigenparameter of \rulename{DE}. In Lean the kernel form
with the newest variable is \Lean{V0Prf}, with \Lean{v0prf\_iff\_nprf}
(\Lean{NativeNewestVariable}); the same proof applies, because the packet
substitutes only for the witness it has just opened and no derivation is
renamed. The clause form $\mathcal V_0$ is \Lean{VCPrf}, with
\Lean{vcprf\_iff\_nprf} (\Lean{NativeVariableClauses}); there the premises
needed under further variables are obtained from the induction hypothesis
at a shifted renaming, which is the protected translation of
\cref{disch:lem:internal}.

\subsection{Substitution for a parameter on raw terms}
\label{disch:sec:param}

In this subsection terms are the raw terms of \cref{base:sec:terms} and
substitution is ordinary capture-avoiding substitution; $\It{0}{a}{u}{h}$
and $a$ are different terms, and nothing is normalized. The rows of a
packet are then not literal, and substitution for a variable has to be
combined with the joining of two equations at the root. Formulas are
equations and $\D{x}{y}B$; contexts consist of clauses $[A;\,y.\,B(y)]$
with one antecedent; the assumption clause of $F$ is
$\operatorname{ord}(F)\defeq[H;\,v.\,F]$.

For terms $s,t,u$ define
\[
 \psi(s,t,u)\defeq
 \begin{cases}
 u\subst{t}{s},&s\text{ is a parameter},\\
 t,&s\text{ is not a parameter and }u\equiv s,\\
 u,&s\text{ is not a parameter and }u\not\equiv s.
 \end{cases}
\]
The first case substitutes for a variable in one pass and is defined also
when $t$ contains $s$. For every $s$, $\psi(s,t,s)\equiv t$. For a source
that is not a variable no proper subterm of $u$ is inspected. Define also
the partial one-step computation
\[
 \kappa(\It{0}{a}{u}{h})\defeq a,\qquad
 \kappa(\It{\Suc n}{a}{u}{h})\defeq h\subst{\It{n}{a}{u}{h}}{u},
\]
undefined on all other terms; the iterator in the second result is left as
it is.

\begin{definition}[The calculus of parameter substitution]
\label{disch:def:param}
$\mathcal P$ has the following five rule families.
\begin{align}
 &\infrule{}{\Gamma\proves b=\kappa(b)}\ \rulename{Compute},
 \qquad\kappa(b)\text{ defined}
 \label{disch:eq:compute}\\[1ex]
 &\infrule{\Gamma\proves s=t\qquad\Gamma\proves u=r}
          {\Gamma\proves r=\psi(s,t,u)}\ \rulename{Param}
 \label{disch:eq:param}\\[1ex]
 &\infrule{\Gamma\proves A}{\Gamma\proves B(t)}\ \rulename{Use},
 \qquad [A;\,y.\,B(y)]\text{ an occurrence in }\Gamma
 \notag\\[1ex]
 &\infrule{\Gamma\proves\D{x}{y}B(x,y)\qquad
           \Gamma,\ \mathcal C_B(b)\proves F}
          {\Gamma\proves F}\ \rulename{Open}
 \notag\\[1ex]
 &\infrule{\Gamma,\ B(t,\eta)\proves B(t,\Suc z)}
          {\Gamma\proves\D{x}{y}B(x,y)}\ \rulename{Close},
 \qquad
 \eta=\begin{cases}
 z,&B\text{ is a literal equation and }t\equiv 0,\\
 0,&\text{otherwise}
 \end{cases}
 \notag
\end{align}
\rulename{Open} is \rulename{Elim} and \rulename{Close} is
\rulename{Intro}, on raw syntax and with the same freshness conditions;
the tests of \rulename{Close} are tests on raw syntax. In this subsection
\rulename{Open} denotes this rule and not the rule of
\cref{disch:eq:open}.
\end{definition}

There is no reflexivity axiom, no hypothesis rule, no template, no
constructor $\Q{}$, no case rule and no induction schema.
\rulename{Compute} contains the two computation equations; counting it as
one family is a grouping of the implementation and does not turn two
equations into one.

\begin{lemma}[Equality at the root]
\label{disch:lem:root}
In $\mathcal P$ the following are derived under every context:
\[
 s=t,\ s=r\proves r=t;\qquad
 \proves a=a;\qquad
 s=t\proves t=s;\qquad
 s=t,\ t=r\proves s=r .
\]
\end{lemma}

\begin{proof}
The first, \emph{joining at a common left side}, is \rulename{Param} with
$u\equiv s$, since $\psi(s,t,s)\equiv t$. For reflexivity,
\rulename{Compute} proves $\It{0}{a}{u}{0}=a$, and joining this equation
with itself gives $a=a$. For symmetry join the reflexive $s=s$, as major
premise, with $s=t$. For transitivity take the symmetric $t=s$ and join
$t=r$, as major premise, with $t=s$. Only \rulename{Compute} and joining
at the root are used.
\end{proof}

\begin{lemma}[Elementary operations]
\label{disch:lem:param-struct}
The use of an assumption, weakening, injective renaming and cut of an
assumption are admissible in $\mathcal P$.
\end{lemma}

\begin{proof}
Reflexivity proves $H$, and \rulename{Use} of $[H;\,v.\,F]$ proves $F$.
Weakening and injective renaming are inductions on derivations; an
injective renaming preserves being a parameter, literal identity of terms,
the root of a formula and the literal $0$, hence the case of $\psi$ and
the mode of \rulename{Close}. For cut, replace every \rulename{Use} of the
occurrence $[H;\,v.\,F]$ by the given derivation of $F$, weakened to the
local context, after renaming eigenparameters. No substitution of a term
for a parameter in a derivation is used.
\end{proof}

\begin{lemma}[Fresh aliases]
\label{disch:lem:alias}
For a term $a$ and a variable $b$ fresh for $\Gamma$, $F$ and $a$,
\[
 \infrule{\Gamma,\ b=a\proves F}{\Gamma\proves F}
\]
is admissible in $\mathcal P$. Its derivation uses the two computation
equations, joining at the root, the logical rules and cut; it uses no
substitution in a proper subterm.
\end{lemma}

\begin{proof}
Let $B(x,y)\defeq\bigl(\sel(y,a,x)=a\bigr)$, where
$\sel(y,a,x)=\It{y}{a}{u}{x}$. \rulename{Close} at the witness $a$ proves
$\D{x}{y}B$: its goal $\sel(\Suc z,a,a)=a$ is the successor case of
\rulename{Compute} for the constant body $a$. The assumption is not used,
so the derivation is valid in either mode, also when $a\equiv 0$.

Open at a fresh $b$. The clause is
$[\sel(0,a,b)=a;\,y.\,\sel(y,a,b)=a]$. Its antecedent is the zero case of
\rulename{Compute}, and \rulename{Use} at $1$ gives $\sel(1,a,b)=a$. The
successor case of \rulename{Compute} gives $\sel(1,a,b)=b$. Joining, with
the former as major premise, gives $b=a$. Cut this into the given
derivation, weakened by the clause; \rulename{Open} discharges $b$.
\end{proof}

A second derivation uses the case mode of \rulename{Close} only: take the
witness $\It{0}{a}{v}{0}$, which is not the literal $0$; the successor
case of \rulename{Compute} gives
$\sel(\Suc z,a,\It{0}{a}{v}{0})=\It{0}{a}{v}{0}$, the zero case gives
$\It{0}{a}{v}{0}=a$, and transitivity joins them. So the lemma does not
depend on the progress mode. The alias is the device of the optional
argument in \pkg{template\_free}; the point here is that it needs neither
substitution in a proper subterm nor transport in a template nor
induction, and therefore cannot hide the principle it is used to derive.

\begin{theorem}[Substitution for a parameter gives every template]
\label{disch:thm:param-congruence}
For every term template $h(v)$, if $\Gamma\proves_{\mathcal P}s=t$ then
$\Gamma\proves_{\mathcal P}h(s)=h(t)$. Repeated occurrences of $v$,
occurrences under the binder of an iterator and the absence of $v$ are all
permitted.
\end{theorem}

\begin{proof}
Work under the assumption $b=s$, with $b$ fresh for $\Gamma$, $h$, $s$ and
$t$. Transitivity gives $b=t$. Reflexivity gives $h(b)=h(b)$.
\rulename{Param} with the major premise $b=s$ and this minor premise gives
$h(b)=h(s)$: the source is the parameter $b$, and substitution for $b$
marks exactly the distinguished occurrences of the template. In the same
way $b=t$ gives $h(b)=h(t)$. Joining $h(b)=h(t)$, as major premise, with
$h(b)=h(s)$ gives $h(s)=h(t)$, which does not contain $b$, and
\cref{disch:lem:alias} discharges the alias.
\end{proof}

The argument is not a pass that replaces every occurrence of the possibly
compound term $s$ in $h(s)$; such a pass would also replace occurrences of
$s$ in the passive part of the template. Only the parameter case of
$\psi$ acts on a proper subterm.

\begin{corollary}[The equality interface of $\mathcal P$]
\label{disch:cor:param-eq}
$\mathcal P$ derives \rulename{Tr}, $s=t,\ h(s)=r\proves r=h(t)$, and
transport in every equation template,
$s=t,\ p(s)=q(s)\proves p(t)=q(t)$. It derives the replacement rule of
\pkg{template\_free}: from $s=t$ and $u=v$ infer $v=u'$, where $u'$
results from $u$ by replacing, in one pass from the root, the occurrences
of $s$ by $t$.
\end{corollary}

\begin{proof}
Join the congruence $h(s)=h(t)$, as major premise, with $h(s)=r$. Apply
this to $p$ and then to $q$, with symmetry, for the equation template. For
the replacement rule, a pass from the root selects finitely many pairwise
disjoint occurrences of $s$ in $u$; replacing exactly these by a fresh
variable gives a template $h$ with $h(s)\equiv u$ and $h(t)\equiv u'$.
The extraction of the template belongs to the producer of the proof and
not to the checker.
\end{proof}

\begin{lemma}[Every case introduction]
\label{disch:lem:param-case}
For every matrix $B$ and witness $t$, from
$\Gamma,B(t,0)\proves B(t,\Suc z)$ one derives
$\Gamma\proves\D{x}{y}B(x,y)$ in $\mathcal P$.
\end{lemma}

\begin{proof}
Only a literal equation $E$ at the witness $0$ is not already the case
mode of \rulename{Close}. Work under a fresh alias $b=0$. From the
assumption $E(b,0)$, transport in an equation template gives $E(0,0)$; cut
in the given derivation to obtain $E(0,\Suc z)$; transport along the
symmetric equation gives $E(b,\Suc z)$. \rulename{Close} at the variable
$b$ is in case mode and concludes $\D{x}{y}E$, which does not contain $b$;
\cref{disch:lem:alias} discharges the alias. Only transport in equations
is used.
\end{proof}

\begin{lemma}[Closures, induction, and transport in formulas]
\label{disch:lem:param-closure}
With $q_yB\defeq\D{w}{y}B$ for an unused $w$, the following are derived in
$\mathcal P$.
\begin{enumerate}[label=\textup{(\alph*)}]
\item From $\Gamma,B(0)\proves B(\Suc z)$ infer $\Gamma\proves q_yB$.
\item From $\Gamma\proves q_yB$ and $\Gamma\proves B(0)$ infer
 $\Gamma\proves B(t)$.
\item A context that contains $[B(0);\,y.\,B(y)]$ proves $q_yB$.
\item For a literal equation $E$, from $\Gamma,E(z)\proves E(\Suc z)$
 infer $\Gamma\proves q_yE$; with (b), induction for literal equations
 under an arbitrary context.
\item Transport in every formula template.
\end{enumerate}
\end{lemma}

\begin{proof}
(a) \rulename{Close} at the unused witness $1$ is in case mode. (b)
\rulename{Open} for the unused witness adds $[B(0);\,y.\,B(y)]$, and
\rulename{Use} gives $B(t)$. (c) In (a), under the assumption $B(0)$, use
the clause at $\Suc z$. (d) \rulename{Close} at the unused literal witness
$0$ is in progress mode. (e) By induction on the template, as in
\cref{disch:lem:formula}: the equation case is
\cref{disch:cor:param-eq}; for $\D{x}{y}K(v,x,y)$ open the source at $b$
and introduce the target at the same $b$ in case mode, converting the base
backwards and the successor instance forwards.
\end{proof}

\begin{theorem}[$\mathcal P$ and the kernel]
\label{disch:thm:param}
Let $\tau$ be the translation of native formulas that leaves terms
unchanged and satisfies $\tau(s=t)=(s=t)$,
$\tau(\D{x}{y}B)=\D{x}{y}\tau B$ and $\tau(\Q{y}B)=\D{w}{y}\tau B$ with
$w$ fresh, and let $j$ be the inclusion of the formulas of $\mathcal P$
into native formulas. For every context $\Gamma$ of native formulas and
every native formula $A$,
\[
 \Gamma\proves_{\Nat}A\quad\Longleftrightarrow\quad
 \operatorname{ord}(\tau\Gamma)\proves_{\mathcal P}\tau A .
\]
Moreover $\tau jF\equiv F$ for every formula $F$ of $\mathcal P$, and
$\Nat$ derives $A$ from $j\tau A$ and conversely.
\end{theorem}

\begin{proof}
$\tau$ commutes with substitution and preserves the property of being a
literal equation.

\emph{From $\Nat$ to $\mathcal P$}, by induction on the derivation.
\rulename{Hyp} is the use of an assumption
(\cref{disch:lem:param-struct}). The two computation rules are the two
cases of \rulename{Compute}. \rulename{Tr} is
\cref{disch:cor:param-eq}. \rulename{QE} is
\cref{disch:lem:param-closure}(b). \rulename{DI} is
\cref{disch:lem:param-case}, and \rulename{QI} is
\cref{disch:lem:param-closure}(d). For \rulename{DE}, open the translated
major premise at $b$; by \cref{disch:lem:param-closure}(c) the opened
clause proves $q_y\tau B(b,y)$, which is $\tau(\Q{y}B(b,y))$; weaken the
translated branch and cut this formula into it; \rulename{Open} closes.

\emph{From $\mathcal P$ to $\Nat$.} Consider contexts each of whose
occurrences is an assumption clause $\operatorname{ord}(G)$ or an opened
clause $\mathcal C_B(b)$, with its origin recorded, and represent them by
the native formulas $jG$ and $\Q{y}jB(b,y)$ respectively. By induction on
a derivation in $\mathcal P$ from such a context, $\Nat$ derives $j$ of its
conclusion from the representing formulas. A \rulename{Use} of an
assumption clause is \rulename{Hyp}, and a \rulename{Use} of an opened
clause is \rulename{QE}. \rulename{Compute} is a computation rule.
\rulename{Param} is an instance of \rulename{Tr}: with the template
$u\subst{v}{s}$, $v$ fresh, if $s$ is a parameter; with the identity
template if $s$ is not a parameter and $u\equiv s$; with the constant
template $u$ otherwise. \rulename{Open} is \rulename{DE}. A
\rulename{Close} in case mode is \rulename{DI}. A \rulename{Close} in
progress mode, at the literal witness $0$ with the matrix $E$, is
\rulename{QI} for $\Q{y}E(0,y)$, followed by \rulename{QE} at $\Suc z$
under the assumption $E(0,0)$ and by \rulename{DI}.

\emph{Round trips.} $\tau j$ is the identity. For $j\tau$ the essential
equivalence in $\Nat$ is $\Q{y}B(y)\dashv\vdash\D{w}{y}B(y)$ for an unused
$w$: from left to right use \rulename{QE} under the base and \rulename{DI}
at any witness; from right to left \rulename{DE} opens exactly
$\Q{y}B(y)$. The equivalence lifts through both binders by the congruence
of binders in $\Nat$ (\cref{base:sec:structural}), and induction on $A$
gives the interderivability of $A$ and $j\tau A$. With cut in $\Nat$ this
turns the second translation into the implication from right to left of
the displayed equivalence.
\end{proof}

The proof of completeness uses the parameter case of $\psi$ and joining at
the root, never the third case, which only reverses the minor premise.

\begin{corollary}[Bridge for $\mathcal P$]
\label{disch:cor:param-ipra}
The ordinary reading of $\mathcal P$ is bi-interpretable and
equiconsistent with $\iPRA$, with whole domains, actual equality, and both
numerical composites with graph $x=y$.
\end{corollary}

\begin{proof}
Compose \cref{disch:thm:param} with \cref{base:thm:bridge}. Terms are
unchanged, so the numerical maps are identities, and $0=\Suc 0$ is fixed
by $\tau$ and $j$.
\end{proof}

\begin{remark}[Substitution, checking, Lean]
\label{disch:rem:param}
Substitution of a term for a parameter does not preserve the rules of
$\mathcal P$: it can turn the source of \rulename{Param} from a variable
into a compound term and can turn a witness of \rulename{Close} into the
literal $0$. Its admissibility follows from \cref{disch:thm:param} and
substitution in $\Nat$ (\cref{base:sec:structural}). The checker performs
structural traversals, capture-avoiding substitutions, literal comparisons
and one-step computations only, so its certificate relation is primitive
recursive; proofs and substituted terms can still grow considerably, since
the producer builds an alias around every transport in a template.

In Lean the kernel with \rulename{Tr} replaced by \rulename{Param} is
\Lean{PPrf}, with \Lean{pprf\_iff\_nprf} (\Lean{NativeParameter}); the
alias is discharged through $\D{x}{y}(\sel(y,s,x)=s)$, proved by cases with
the witness $s$ and eliminated at $y=1$ (\Lean{PPrf.aliasE}), and the
calculus is closed under injective renaming (\Lean{PPrf.ren}). The five
families of $\mathcal P$, which combine the rule with clause contexts and
the introduction with a mode, are not formalized together.
\end{remark}

\subsection{Unary congruence does not suffice}
\label{disch:sec:unary}

\rulename{VSub} and \rulename{Param} are binary: they combine two
equations. One may ask whether the equality rule can be the unary
congruence. On raw and on zero-normal terms it cannot.

\begin{proposition}[An obstruction]
\label{disch:prop:unary}
Let $\Nat^{c}$ be the calculus obtained from $\Nat$ by replacing
\rulename{Tr} with reflexivity, symmetry and the rule
\[
 \infrule{\Gamma\proves s=t}{\Gamma\proves h(s)=h(t)}
\]
for every term template $h$, all other rules being kept. Then $\Nat^{c}$
does not derive $x=y,\ y=z\proves x=z$. The same holds on zero-normal
terms, where $\It{0}{a}{u}{h}$ is identified with $a$.
\end{proposition}

\begin{proof}
Let $V=\{0,m,2\}$, ordered $0<m<2$, and call two elements \emph{near} if
they are equal or adjacent; so $0$ and $2$ are not near, and $m$ is near
to every element. Interpret $0$ as $0$, successor as the constant function
with value $2$, and
\[
 \It{n}{a}{u}{h}\ \text{ as }\
 \begin{cases}a,&n=0,\\ m,&n\neq 0,\end{cases}
\]
the body being ignored. An equation $s=t$ is \emph{true} under an
assignment of elements of $V$ to the variables if the values of $s$ and
$t$ are near. This is a truth assignment to judgments and not a structure
with actual equality. $\Q{y}B$ and $\D{x}{y}B$ are given the Boolean
readings of \cref{base:sec:reading} over $V$. A sequent is valid if every
assignment that makes its context true makes its conclusion true. The
value of a term commutes with substitution, since bodies are ignored, and
the value of a term is unchanged by zero-normalization.

\emph{The rules other than the introductions.} \rulename{Hyp},
reflexivity and symmetry are valid. The zero computation holds with equal
values. In the successor computation the left side has the value $m$,
which is near to every value. \rulename{QE} and \rulename{DE} are valid
for the Boolean readings. For the congruence rule it suffices that the
interpretation of the iterator preserves nearness in both arguments
together: if both counters are $0$ the outputs are the two seeds; if
exactly one is $0$ the outputs are a seed and $m$; otherwise both are $m$.
Zero and successor are constants.

\emph{The shape of term functions.} Every term denotes either the constant
$m$, or a function that takes the value of a \emph{payload}, which is $0$,
$2$ or one variable, when all variables of a finite \emph{gate} are $0$,
and the value $m$ otherwise. This holds for variables, $0$ and successors
with an empty gate, and for $\It{n}{a}{u}{h}$ by cases on the counter: if
$n$ is constantly $m$ or has the payload $2$, the iterator is constantly
$m$; if $n$ has the payload $0$ and the gate $G$, the iterator is $a$ with
$G$ added to its gate; if $n$ has a variable $w$ as payload, the iterator
is $a$ with $G\cup\{w\}$ added to its gate. Consequently, as a function of
one variable $v$ with the others fixed, a term takes at $(0,m,2)$ one of
the triples
\[
 (0,0,0),\quad(m,m,m),\quad(2,2,2),\quad(0,m,2),\quad(0,m,m),\quad(2,m,m).
\]

\emph{The key property.} Every formula $B(v)$ satisfies, for every
variable $v$ and all values of the other variables,
\begin{equation}
 B(0)\ \text{and}\ B(2)\ \text{imply}\ B(m).
 \label{disch:eq:convex}
\end{equation}
For an equation: at $v=m$ each side has the value $m$ or is constant in
$v$. If one side has the value $m$ there, the equation is true at $m$. If
both are constant, its truth does not depend on $v$.

For $\Q{y}E$ with $E$ an equation: the formula is false exactly if $E$ is
true at $y=0$ and $y=m$ and false at $y=2$. Falsity at $2$ means that the
two sides have the values $0$ and $2$ there, so neither side has $y$ in
its gate, all gate variables of both sides have the value $0$, and each
side is its payload. Inspecting the payloads, the formula is then false
only if the payloads are $y$ and $0$, or $y$ and a variable $w$ whose
value is $0$. Hence $\Q{y}E$ is true, or false, or equivalent to the
statement that some variable of a fixed finite set is different from $0$.
A formula of this last kind satisfies \eqref{disch:eq:convex}: if $v$
belongs to the set it is true at $v=m$, and otherwise it does not depend
on $v$. For
$\D{x}{y}E$, the reading is the existence of a value of $x$ at which
$\Q{y}E$ holds; if $x$ is in the fixed finite set the formula is true, and
otherwise it is the same statement as before. Finally, if $B$ is true,
false, or the statement that some variable of a finite set is different
from $0$, then $\Q{y}B$ is true: if $y$ is not in the set the instances
do not depend on $y$, and if it is, the instance at $y=0$ implies all
instances. So $\D{x}{y}B$ is true as well, and so is every formula with
more binders. % ed.: the classification of the term functions and of the formulas Q E is written out; the package states both without the case analysis.

\emph{The introductions.} For \rulename{DI}, take the value of the witness
$t$ for $x$. If $B(t,0)$ is true, the premise gives $B(t,\Suc z)$, that
is, $B$ at $y=2$, and \eqref{disch:eq:convex} gives $B$ at $y=m$. For
\rulename{QI}, if $E(0)$ is true, the premise with $z$ assigned $0$ gives
$E$ at $2$, and \eqref{disch:eq:convex} gives $E$ at $m$.

So every derivable sequent is valid. Under $x\mapsto 0$, $y\mapsto m$,
$z\mapsto 2$ the equations $x=y$ and $y=z$ are true and $x=z$ is false.
The contradiction $0=\Suc 0$ is false as well.
\end{proof}

Full normalization of the successor computation is outside the
obstruction: it identifies values that the interpretation only relates.
The proposition excludes this one deletion in the neighborhood of the
kernel; it does not exclude other presentations of arithmetic.

\section{Specialization without a term}
\label{disch:sec:spec}

This section draws on \pkg{deterministic\_specialization}. The calculus it
restricts is the five-rule clause calculus of \pkg{uniform\_step}.

\subsection{The calculus}

In this section terms are \emph{zero-normal}: the normalization contracts
$\It{0}{a}{u}{h}$ to $a$ in all subterms and does not contract an iterator
whose counter is a successor. Terms are the fixed points of this map, and
every substitution is followed by it, including a substitution that makes
a counter the literal $0$; in particular $\sel(0,a,b)\equiv a$. Formulas
are equations and $\D{x}{y}F$, so that every formula is a chain of binders
ending in one equation; its \emph{rank} is the number of binders. A
context is a finite list of clauses $[A;\,y.\,B(y)]$, written with the
\emph{newest} occurrence first, and the parameters of a sequent form a
list whose last member is the newest. An assumption $A$ is the clause
$[0=0;\,v.\,A]$. The order is used only to choose the target of a
specialization; clauses are reusable.

\begin{definition}[Deterministic specialization]
\label{disch:def:spec}
The rules are
\begin{align}
 &\infrule{}{\Gamma\proves t=t}\ \rulename{Refl}
 \qquad
 \infrule{\Gamma\proves\It{c}{a}{u}{h}=t\qquad
          \Gamma\proves\It{\Suc c}{a}{u}{h}=r}
         {\Gamma\proves r=h\subst{t}{u}}\ \rulename{Step}
 \label{disch:eq:step}\\[1ex]
 &\infrule{\Gamma\proves\D{x}{y}B(x,y)\qquad
           [B(b,0);\,y.\,B(b,y)],\ \Gamma\proves F}
          {\Gamma\proves F}\ \rulename{DE}
 \notag\\[1ex]
 &\infrule{[0=0;\,v.\,B(t,\eta)],\ \Gamma\proves B(t,\Suc z)}
          {\Gamma\proves\D{x}{y}B(x,y)}\ \rulename{DI},
 \qquad
 \eta=\begin{cases}
 z,&B\text{ is a literal equation and }t\equiv 0,\\
 0,&\text{otherwise}
 \end{cases}
 \notag\\[1ex]
 &\infrule{\Gamma\proves A}{\Gamma\proves B(t)}\ \rulename{Use},
 \qquad [A;\,y.\,B(y)]\text{ an occurrence in }\Gamma,\ t\text{ arbitrary}
 \notag\\[1ex]
 &\infrule{\Gamma\proves A}{\Gamma\proves B(v_i)}\ \rulename{Use}_*,
 \qquad [A;\,y.\,B(y)]=\Gamma_i,\quad
 v_i=\begin{cases}1,&i=0,\\ b_*,&i>0.\end{cases}
 \label{disch:eq:usestar}
\end{align}
In \rulename{Step} the first iterator is zero-normal, so it is $a$ when
$c\equiv 0$; repeated expressions are literally identical. In
\rulename{DE} the variable $b$ is fresh for $\Gamma$, $F$ and
$\D{x}{y}B$, and in \rulename{DI} the variable $z$ is fresh for $\Gamma$,
$t$ and $\D{x}{y}B$; each becomes the newest parameter of its premise, and
the displayed clause becomes the newest clause. In $\rulename{Use}_*$,
$\Gamma_0$ is the newest clause and $b_*$ is the newest parameter, or $0$
if there is none. $\mathcal U$ is the calculus with \rulename{Refl},
\rulename{Step}, \rulename{Use}, \rulename{DE} and \rulename{DI};
$\mathcal U_*$ has $\rulename{Use}_*$ in place of \rulename{Use}.
\end{definition}

So the newest clause is used at $1$ and every older clause at the newest
parameter. A certificate of $\rulename{Use}_*$ consists of the index of
the clause and the proof of its antecedent; it contains no term and no
mode. The witness of \rulename{DI} is still supplied, and the completeness
proof uses that facility to name the desired term.

$\mathcal U$ is the calculus of \pkg{uniform\_step} with clause contexts;
its comparison with the kernel belongs to that package, and its arithmetic
interface on the kernel is \Lean{NSPrf}, with \Lean{nsprf\_iff\_nprf}. The
facts about it that are used here are the following.

\begin{lemma}[Arithmetic from \rulename{Step}]
\label{disch:lem:step}
Using \rulename{Refl} and \rulename{Step} only, and opening no scope, both
$\mathcal U$ and $\mathcal U_*$ derive: the pivot
$s=t,\ s=r\proves r=t$; symmetry and transitivity; the congruence
$s=t\proves h(s)=h(t)$; transport in every equation template; and the
successor computation
$\It{\Suc c}{a}{u}{h}=h\subst{\It{c}{a}{u}{h}}{u}$, in particular
$\sel(\Suc z,a,b)=b$.
\end{lemma}

\begin{proof}
Let $J_s\defeq\It{1}{s}{u}{u}$. Since $\It{0}{s}{u}{u}\equiv s$,
\rulename{Step} with $c\equiv 0$ sends $s=t$ and $J_s=r$ to $r=t$. Given
$p\colon s=t$ and $q\colon s=r$, \rulename{Step} applied to $q$ and the
reflexive $J_s=J_s$ gives $J_s=r$, and \rulename{Step} applied to $p$ and
this gives $r=t$. Symmetry is the pivot with the first premise $s=s$, and
transitivity follows. For a template $h$ let $K\defeq\It{1}{s}{u}{h}$;
\rulename{Step} applied to the reflexive $s=s$ and $K=K$ gives $K=h(s)$,
and \rulename{Step} applied to $s=t$ and $K=h(s)$ gives $h(s)=h(t)$. An
equation template $p(v)=q(v)$ is transported through $p(t)=p(s)$,
$p(s)=q(s)$ and $q(s)=q(t)$. The successor computation is \rulename{Step}
applied to the reflexivities of $\It{c}{a}{u}{h}$ and
$\It{\Suc c}{a}{u}{h}$.
\end{proof}

Conversely the two computation equations and transport give
\rulename{Step}: in the computation equation for $\It{\Suc c}{a}{u}{h}$
replace $\It{c}{a}{u}{h}$ by $t$ in the body and the successor iterator by
$r$. So \rulename{Step} adds no arithmetic strength.

\begin{lemma}[The unrestricted calculus]
\label{disch:lem:spec-unrestricted}
In $\mathcal U$ cut of an assumption is admissible, and the following are
derived: case introduction for every matrix and witness; for a literal
equation $E$ and every witness $t$, the progress introduction from
$\Gamma,E(t,z)\proves E(t,\Suc z)$ to $\Gamma\proves\D{x}{y}E$; and
transport in every formula template.
\end{lemma}

\begin{proof}
Cut is an induction on the derivation: a use of the assumption clause is
replaced by the given proof. The term $c\defeq\It{1}{0}{u}{u}$ is
zero-normal, is not the literal $0$, and \rulename{Step} proves $c=0$.
Given a premise $\Gamma,E(0,0)\proves E(0,\Suc z)$ for a literal equation
$E$, start from $E(c,0)$, transport to $E(0,0)$, insert the given
derivation, and transport back to $E(c,\Suc z)$; \rulename{DI} at $c$ is
in case mode. For progress at $t$, introduce $\D{w}{y}E(t,y)$ with an
unused $w$ at the witness $0$; \rulename{DE} adds
$[E(t,0);\,y.\,E(t,y)]$; under the base $E(t,0)$ use it at $\Suc z$ and
apply the case introduction at $t$. Transport through a binder is as in
\cref{disch:lem:formula}. A progress introduction for a compound matrix is
never used; it would change the induction of the theory.
\end{proof}

\subsection{Naming a term}

For a term $t$ let
\[
 \mathrm{Nm}_t\defeq\D{x}{y}\bigl(\sel(y,0,x)=\sel(y,0,t)\bigr),\qquad
 \mathcal M_{b,t}\defeq[0=0;\,y.\,\sel(y,0,b)=\sel(y,0,t)].
\]

\begin{lemma}[Naming]
\label{disch:lem:spec-name}
In every context, $\mathcal U_*$ derives $\mathrm{Nm}_t$. Opening it by
\rulename{DE} at a fresh $b$ adds the clause $\mathcal M_{b,t}$, and in
the context of that premise, where $\mathcal M_{b,t}$ is the newest
clause, $\mathcal U_*$ derives $b=t$.
\end{lemma}

\begin{proof}
Apply \rulename{DI} at the witness $t$. The goal
$\sel(\Suc z,0,t)=\sel(\Suc z,0,t)$ has literally equal sides and is
proved by \rulename{Refl}, in either mode. The antecedent of the opened
clause is $\sel(0,0,b)=\sel(0,0,t)$, that is, $0=0$. Since the clause is
newest, $\rulename{Use}_*$ applies it at $1$, with \rulename{Refl} for its
antecedent, and gives $\sel(1,0,b)=\sel(1,0,t)$; the computations
$\sel(1,0,b)=b$ and $\sel(1,0,t)=t$ of \cref{disch:lem:step} give $b=t$.
\end{proof}

This derivation of $b=t$ is valid immediately after the opening. It does
not remain a derivation when the context is extended, because the clause
then ceases to be newest. An extension of the scope changes $b_*$, and the
insertion of a clause changes which clause is newest; so weakening is not
an operation that preserves the rule at each node of a derivation of
$\mathcal U_*$. The completeness proof replaces weakening by procedures
that produce a proof again in every larger context.

\subsection{Persistent proofs}

A \emph{situation} is a pair $(\Xi;\Gamma)$ of a parameter list and a
context over it. An \emph{extension} of $(\Xi;\Gamma)$ is a situation
$(\Xi';\Gamma')$ together with an injective map $\iota$ of the members of
$\Xi$ to members of $\Xi'$ and an injective map of the occurrences of
$\Gamma$ to occurrences of $\Gamma'$ that sends an occurrence of a clause
$C$ to an occurrence of $C\iota$. Extensions compose, and neither map need
preserve the order.

\begin{definition}[Persistent proof]
\label{disch:def:persistent}
A \emph{persistent proof} of a formula $A$ over $(\Xi;\Gamma)$ is a
function that assigns to every extension of $(\Xi;\Gamma)$ a derivation
in $\mathcal U_*$ of $\Gamma'\proves A\iota$.
\end{definition}

A persistent proof is a finite proof-producing procedure of the
metatheory; it is not an additional rule.

\begin{lemma}[Elementary persistence]
\label{disch:lem:persistent}
\begin{enumerate}[label=\textup{(\alph*)}]
\item A persistent proof over a situation restricts to a persistent proof
 over each of its extensions.
\item If $[0=0;\,v.\,A]$ occurs in $\Gamma$, then $A$ has a persistent
 proof over $(\Xi;\Gamma)$.
\item An equation $t=t$ has a persistent proof; if the two premises of a
 \rulename{Step} have persistent proofs, so has its conclusion. Hence the
 conclusions of the derivations of \cref{disch:lem:step} have persistent
 proofs when their premises have.
\end{enumerate}
\end{lemma}

\begin{proof}
(a) Compose extensions. (b) In every extension apply $\rulename{Use}_*$
to the image of the occurrence, with \rulename{Refl} for the antecedent;
the bound variable is unused, so either target gives $A\iota$. (c) Apply
the rule in each extension; \rulename{Refl} and \rulename{Step} do not
depend on the context or on the order of the parameters.
\end{proof}

\begin{lemma}[Specialization of an equation clause]
\label{disch:lem:spec-atomic}
Let $[A;\,y.\,E(y)]$ occur in $\Gamma$ with $E$ a literal equation, let
$A$ have a persistent proof over $(\Xi;\Gamma)$, and let $t$ be a term
over $\Xi$. Then $E(t)$ has a persistent proof over $(\Xi;\Gamma)$.
Consequently, if $\mathcal M_{b,t}$ occurs in $\Gamma$, at any position,
then $b=t$ has a persistent proof over $(\Xi;\Gamma)$.
\end{lemma}

\begin{proof}
Let an extension be given; write $t$, $E$, $A$ for the renamed
expressions. Derive $\mathrm{Nm}_t$ and open it by \rulename{DE} at a
fresh $b$ (\cref{disch:lem:spec-name}). In the premise, $b$ is the newest
parameter and $\mathcal M_{b,t}$ the newest clause. The persistent proof
of $A$, taken at this larger situation, gives $A$. The given clause is now
an older clause, so $\rulename{Use}_*$ applies it at $b_*\equiv b$ and
gives $E(b)$. \Cref{disch:lem:spec-name} gives $b=t$, and transport in the
equation template $E$ (\cref{disch:lem:step}), which opens no scope, gives
$E(t)$. This formula does not contain $b$, and \rulename{DE} closes. The
construction is uniform in the extension, so it is a persistent proof.

For the second statement apply the first to the clause
$\mathcal M_{b,t}$, the target $1$ and the antecedent $0=0$, which has the
persistent proof \rulename{Refl}; this gives
$\sel(1,0,b)=\sel(1,0,t)$ persistently, and
\cref{disch:lem:persistent}(c) with the selector computations gives $b=t$.
The argument is not circular: the auxiliary naming clause of the first
part is newest when it is used at $1$, and the final transport there is
between equations.
\end{proof}

\begin{lemma}[Transport and specialization at every rank]
\label{disch:lem:spec-rank}
For every $r\geq 0$ the following hold.
\begin{enumerate}[label=\textup{(\alph*)}]
\item[$(\mathrm T_r)$] Let $F(w)$ be a formula template of rank $r$ and
 let $s=t$ have a persistent proof over $(\Xi;\Gamma)$. From a derivation
 in $\mathcal U_*$ of $\Gamma\proves F(s)$ one obtains a derivation of
 $\Gamma\proves F(t)$; and from a persistent proof of $F(s)$ over
 $(\Xi;\Gamma)$ a persistent proof of $F(t)$.
\item[$(\mathrm S_r)$] Let $[A;\,y.\,F(y)]$ occur in $\Gamma$ with $F$ of
 rank $r$, let $A$ have a persistent proof over $(\Xi;\Gamma)$, and let
 $t$ be a term over $\Xi$. Then $F(t)$ has a persistent proof over
 $(\Xi;\Gamma)$.
\end{enumerate}
\end{lemma}

\begin{proof}
$(\mathrm T_0)$ is transport in an equation template
(\cref{disch:lem:step}), with the persistent proof of $s=t$ taken at the
identity extension; $(\mathrm S_0)$ is \cref{disch:lem:spec-atomic}. In
every $(\mathrm T_r)$ the second statement follows from the first, applied
in each extension to the restricted persistent proofs.

$(\mathrm T_r)$ for $r>0$, from $(\mathrm T_{r-1})$ and
$(\mathrm S_{r-1})$. Let $F(w)=\D{x}{y}G(w,x,y)$. Use the given derivation
of $F(s)$ as the major premise of \rulename{DE}, with a fresh witness $c$;
the context of the branch contains the clause
$\mathcal C\defeq[G(s,c,0);\,y.\,G(s,c,y)]$. In the branch introduce
$F(t)$ by \rulename{DI} at the witness $c$. As $c$ is a parameter it is
not the literal $0$, so this is the case mode also when $G$ is an
equation, and the premise has the assumption clause of $G(t,c,0)$ and the
goal $G(t,c,\Suc z)$. In the situation of that premise:
\begin{itemize}
\item $G(t,c,0)$ has a persistent proof by
 \cref{disch:lem:persistent}(b);
\item $t=s$ has a persistent proof by \cref{disch:lem:persistent}(a),(c),
 so $(\mathrm T_{r-1})$ for the template $G(w,c,0)$ gives a persistent
 proof of $G(s,c,0)$;
\item $(\mathrm S_{r-1})$ for the clause $\mathcal C$, this antecedent and
 the target $\Suc z$ gives a persistent proof of $G(s,c,\Suc z)$;
\item $(\mathrm T_{r-1})$ for the template $G(w,c,\Suc z)$, applied to the
 derivation that this persistent proof yields at the identity extension,
 gives a derivation of $G(t,c,\Suc z)$.
\end{itemize}
\rulename{DI} and \rulename{DE} close; $F(t)$ does not contain $c$. The
given derivation of $F(s)$ is used outside the new scopes and is not
weakened.

$(\mathrm S_r)$ for $r>0$, from $(\mathrm T_r)$ and
\cref{disch:lem:spec-atomic}. In a given extension derive $\mathrm{Nm}_t$
and open it at a fresh $b$. In the premise, the persistent proof of $A$
gives $A$, and $\rulename{Use}_*$ applies the given clause, now older, at
$b$, which gives a derivation of $F(b)$. By \cref{disch:lem:spec-atomic},
$b=t$ has a persistent proof over the situation of the premise; it is
needed in persistent form because $(\mathrm T_r)$ opens further scopes.
$(\mathrm T_r)$ gives $F(t)$, and \rulename{DE} closes. The construction
is uniform in the extension.

The recursion is well founded: $(\mathrm T_r)$ uses rank $r-1$ only, and
$(\mathrm S_r)$ uses $(\mathrm T_r)$ and rank $0$.
\end{proof}

\subsection{Equivalence}

\begin{theorem}[Deterministic specialization]
\label{disch:thm:spec}
For every parameter list, every context $\Gamma$ and every formula $A$,
\[
 \Gamma\proves_{\mathcal U}A\quad\Longleftrightarrow\quad
 \Gamma\proves_{\mathcal U_*}A .
\]
The terms, the formulas and the ordinary reading are unchanged.
\end{theorem}

\begin{proof}
Every instance of $\rulename{Use}_*$ is an instance of \rulename{Use}.
Conversely, by induction on a derivation in $\mathcal U$ of
$\Gamma\proves A$ over $\Xi$, the formula $A$ has a persistent proof over
$(\Xi;\Gamma)$.

\rulename{Refl} and \rulename{Step}: \cref{disch:lem:persistent}(c). A
\rulename{Use} of $[A';\,y.\,B]$ at $t$: the induction hypothesis gives a
persistent proof of $A'$, and $(\mathrm S_r)$ of
\cref{disch:lem:spec-rank}, with $r$ the rank of $B$, gives a persistent
proof of $B(t)$. \rulename{DE} with the major premise $\D{x}{y}B$ and the
branch over $(\Xi,b;\,\mathcal C_B(b),\Gamma)$: in a given extension
choose a fresh $b'$, take the derivation of the major premise that its
persistent proof yields there, extend the extension by $b\mapsto b'$ and
by sending the opened clause to the new opened clause, take the derivation
that the persistent proof of the branch yields at this extension, and
apply \rulename{DE}. \rulename{DI} is treated in the same way; an
injective renaming preserves the root of the matrix and the literal test
of the witness, hence the mode.

The compilation recurses on the height of the given derivation, and the
auxiliary procedures recurse on the rank of formulas; there is no
recursion on a derivation that the procedure has produced. So every
derivation in $\mathcal U$ yields a finite derivation in $\mathcal U_*$:
evaluate the persistent proof at the identity extension.
\end{proof}

Composed with the comparison of \pkg{uniform\_step} between $\mathcal U$
and the kernel and with \cref{base:thm:bridge}, the theorem gives the
bi-interpretation of the ordinary reading of $\mathcal U_*$ with $\iPRA$,
with whole domains, actual equality and identity numerical composites; the
equivalence itself changes no term, no formula and no reading.

\begin{remark}[The cost]
\label{disch:rem:spec-cost}
The argument of an arbitrary specialization is removed from the checker at
the price of larger proofs: the translation wraps every use of a clause at
a term in a naming formula, and at rank $r$ it nests the constructions of
ranks below $r$. In the test suite of the package the largest expansion is
from $699$ to $175{,}640$ bytes of compact JSON. No polynomial bound on the
size of the translation is claimed.
\end{remark}

\begin{remark}[The kernel counterpart in Lean]
\label{disch:rem:spec-lean}
On the kernel the corresponding restriction concerns the term of
\rulename{QE}. With \rulename{QE} allowed at a variable and at the numeral
$1$ only, the kernel derives the same sequents: this is \Lean{VQPrf} with
\Lean{vqprf\_iff\_nprf} (\Lean{NativeVariableQE}). An arbitrary target $t$
is named by a fresh alias $b=t$, whose discharge uses \rulename{QE} at $1$
once; \rulename{QE} is applied at $b$, and $B(b)$ is moved to $B(t)$ by
the Leibniz rule, which is proved together with \rulename{QE} at
$\Suc z$ by induction on the formula (\Lean{VQPrf.qELeib},
\Lean{VQPrf.leibniz}). The numeral cannot be dropped
(\Lean{vqprf\_var\_incomplete}): with \rulename{QE} at variables only, the
kernel does not derive $0=1$ from $\Q{y}(\sel(y,0,1)=0)$, which $\Nat$
derives by \rulename{QE} at $1$. The countermodel reads every variable as
$0$, every equation by its value in the natural numbers, and every formula
with root $\Q{}$ or $\D{}{}$ as true; each rule preserves truth, since
$B(v)$ at a variable $v$ has the truth value of $B(0)$. The form of the
package, on zero-normal terms with clause contexts and the choice of the
target by position, is not formalized; that choice is not stable under
renaming, which is why the proof above regenerates the proofs of premises.
\end{remark}

\section{Fresh aliases and backward substitution}
\label{disch:sec:alias}

This section draws on \pkg{alias\_back}
(\cref{disch:sec:alias-rules,disch:sec:alias-eq,disch:sec:alias-kernel,disch:sec:alias-newest})
and \pkg{equation\_aliases} (\cref{disch:sec:alias-atomic}).

In \cref{disch:sec:param} an alias $b=a$ is discharged only from a
conclusion that does not contain $b$ (\cref{disch:lem:alias}), and the
discharge is derived from the computations, joining at the root and the
logical rules. Here the discharge of a fresh alias is a primitive rule, and
it substitutes the defining term into the conclusion. Together with
substitution for a parameter in the backward direction it gives
reflexivity, symmetry and transport in every formula on raw terms, with no
computation, normalization or induction (\cref{disch:prop:alias-eq}). With
it the source of the equality rule and the target of every use of a clause
can both be restricted to the newest parameter
(\cref{disch:thm:alias-compiler}), and, at the price of an induction on the
rank of formulas, both new rules can be restricted to literal equations
(\cref{disch:thm:alias-atomic}). The arithmetic is unchanged: the two
computation equations, the elimination \rulename{Open} and the introduction
\rulename{Close} with a supplied witness are retained.

In this section terms are the raw terms of \cref{base:sec:terms}, and
substitution is the capture-avoiding substitution of raw syntax; it is
simultaneous and performed in one pass, so that $F\subst{t}{b}$ is defined
also when $t$ contains $b$. Nothing is normalized. A judgment
$\Pi;\Gamma\proves A$ carries its numerical scope $\Pi$ as in
\cref{wit:def:scope}, and $\nu(\Pi)$ is the newest parameter of $\Pi$, or
$0$ if $\Pi$ is empty; every rule with an eigenparameter appends it to the
scope of its premise. Where the scope plays no role it is not written. The
clause calculi of the section have the uniform clauses of
\cref{claus:def:v}, with at most one antecedent, the assumption clause
$\langle A\rangle=[\ ;\,v.\,A]$, and no reflexivity clause; $\Gamma,F$
denotes $\Gamma$ extended by $\langle F\rangle$.

\subsection{The two rules}
\label{disch:sec:alias-rules}

\begin{definition}[Alias and backward substitution]
\label{disch:def:alias}
\begin{align}
 &\infrule{\Pi,b;\ \Gamma,\ b=t\proves F}
          {\Pi;\Gamma\proves F\subst{t}{b}}\ \rulename{Alias}
 \label{disch:eq:alias}\\[1ex]
 &\infrule{\Pi;\Gamma\proves b=t\qquad\Pi;\Gamma\proves F\subst{t}{b}}
          {\Pi;\Gamma\proves F}\ \rulename{Back}
 \label{disch:eq:back}
\end{align}
In \rulename{Alias} the parameter $b$ is not in $\Pi$, hence fresh for
$\Gamma$ and for the \emph{defining term} $t$, which is over $\Pi$; the
formula $F$ is over $\Pi,b$ and may contain $b$, and every free occurrence
of $b$ in $F$ is replaced in the conclusion. In a clause calculus the added
assumption is the clause $\langle b=t\rangle$. In \rulename{Back} the
\emph{source} $b$ is a parameter of $\Pi$, the term $t$ and the formula $F$
are over $\Pi$, and $t$ may contain $b$. The conclusion $F$ is the goal: no
template is supplied, and the second premise is computed from the goal and
the first premise. The first premise is oriented; a derivation of $t=b$
does not take its place. The restricted forms are
$\rulename{Back}_{\nu}$, in which $\Pi$ is not empty and $b\equiv\nu(\Pi)$;
$\rulename{Alias}^{=}$ and $\rulename{Back}^{=}$, in which $F$ is a literal
equation; and $\rulename{Back}^{=}_{\nu}$, with both restrictions.
\pkg{equation\_aliases} calls $\rulename{Alias}^{=}$ and
$\rulename{Back}^{=}_{\nu}$ \rulename{EqAlias} and \rulename{EqBack}.
\end{definition}

The essential difference from \cref{disch:lem:alias} is that the
conclusion of \rulename{Alias} may change on discharge: from its own
assumption the rule exports $(b=t)\subst{t}{b}\equiv(t=t)$. This is the
step that is missing from the rule of parameter substitution in the reverse
direction with an ordinary elimination of a fresh name, which
\pkg{uniform\_binding} considered and discarded (\cref{ctr:sec:questions}):
once the goal no longer mentions the alias, that rule can require the goal
itself as its premise. \rulename{Alias} is \rulename{Name}
\eqref{disch:eq:name} without its paired formula: the equation $b=t$ is
assumed instead of being proved from an opened clause. \rulename{Back}
differs from \rulename{VSub} \eqref{disch:eq:vsub} in its direction and in
its scope: \rulename{VSub} substitutes forward in the left side of an
equation, and \rulename{Back} computes its premise from a goal that is an
arbitrary formula. Like \rulename{BRep} of \cref{kern:def:B} it is
goal-directed; unlike \rulename{BRep} it substitutes for a parameter and
does not replace occurrences of a compound term.

\begin{proposition}[The freshness conditions]
\label{disch:prop:alias-fresh}
Neither freshness condition of \rulename{Alias} can be dropped.
\begin{enumerate}[label=\textup{(\alph*)}]
\item If $b$ may be free in $\Gamma$, the rule derives $0=\Suc 0$ from the
 assumption $b=\Suc 0$.
\item If $b$ may occur in $t$, the rule derives $\Suc b=\Suc\Suc b$ from
 the empty context.
\end{enumerate}
Both conclusions are false in the natural numbers under every assignment
that makes the context true, while all other rules of the calculi of this
section are sound for the ordinary reading in the natural numbers.
\end{proposition}

\begin{proof}
(a) Let $\Gamma$ be the assumption $b=\Suc 0$ and $t\equiv 0$. The premise
$\Gamma,\ b=0\proves b=\Suc 0$ is a use of an assumption, and the discharge
would export $(b=\Suc 0)\subst{0}{b}\equiv(0=\Suc 0)$, although $\Gamma$ is
true at $b=1$. (b) Let $t\equiv\Suc b$. The premise
$b=\Suc b\proves b=\Suc b$ is a use of the assumption, and the discharge
would export $(b=\Suc b)\subst{\Suc b}{b}\equiv(\Suc b=\Suc\Suc b)$, which
is false at every value of $b$.
\end{proof}

The checkers of \pkg{alias\_back} validate the defining term in the outer
scope before the alias parameter is added, so that both instances above are
rejected. \rulename{Back} has no freshness condition: its replacement may
contain its source, because \rulename{Back} discharges nothing.

\begin{lemma}[Soundness in the kernel]
\label{disch:lem:alias-sound}
Both rules are admissible in $\Nat$. If $\Gamma,\ b=t\proves_{\Nat}F$ and
$b$ is not free in $\Gamma$ or in $t$, then
$\Gamma\proves_{\Nat}F\subst{t}{b}$. If $\Gamma\proves_{\Nat}b=t$ and
$\Gamma\proves_{\Nat}F\subst{t}{b}$, then $\Gamma\proves_{\Nat}F$.
\end{lemma}

\begin{proof}
For \rulename{Alias} apply \cref{base:lem:struct} with the substitution
$\theta=[t/b]$ and the context $\Gamma,\ t=t$. Since $b$ is not free in
$\Gamma$, every member of $\Gamma$ is unchanged by $\theta$, and since $b$
is not free in $t$, $(b=t)[\theta]\equiv(t=t)$; so
$\Gamma,\ t=t\proves F\subst{t}{b}$. Reflexivity
(\cref{base:lem:eq}(i)) and \cref{base:lem:cut}, with \rulename{Hyp} for
the members of $\Gamma$, give $\Gamma\proves F\subst{t}{b}$.

For \rulename{Back} let $w$ be fresh and $K\defeq F\subst{w}{b}$. Then
$K\subst{b}{w}\equiv F$ and $K\subst{t}{w}\equiv F\subst{t}{b}$, also when
$t$ contains $b$, because the substitution is simultaneous. By
\cref{base:thm:leibniz} with $\Gamma\proves b=t$, the formulas
$K\subst{b}{w}$ and $K\subst{t}{w}$ are interderivable over $\Gamma$.
\end{proof}

\subsection{Equality from three rules}
\label{disch:sec:alias-eq}

\begin{lemma}[Structural operations with unrestricted sources]
\label{disch:lem:alias-struct}
Let $\mathcal S$ have the rules \rulename{Alias} and \rulename{Back}, or
$\rulename{Alias}^{=}$ and $\rulename{Back}^{=}$, together with either the
rules of $\Nat$ other than \rulename{Tr}, or the rules \rulename{I0},
\rulename{IS}, \rulename{Use} at an arbitrary term, and \rulename{Open} and
\rulename{Close} of \cref{disch:def:param} over uniform clauses, with
$\langle B(t,\eta)\rangle$ as the assumption of \rulename{Close}. Then
weakening and injective renaming are admissible in $\mathcal S$: if
$\Pi;\Gamma\proves A$, every parameter of $\Pi$ is in $\Pi'$, in any
order, and every occurrence of $\Gamma$ is an occurrence of a context
$\Gamma'$ over $\Pi'$, then $\Pi';\Gamma'\proves A$; and if $\iota$ is an
injective map of parameters, then $\Pi\iota;\Gamma\iota\proves A\iota$.
Both transformations preserve the rule at each node and the height.
\end{lemma}

\begin{proof}
By induction on the derivation, after the eigenparameters and the alias
parameters of the derivation have been renamed away from $\Pi'$ and
$\Gamma'$, respectively from the range of $\iota$. A use of an occurrence
is a use of its image. The source of \rulename{Back} is a parameter of the
old scope, hence of the new one, and
$(F\iota)\subst{t\iota}{b\iota}\equiv(F\subst{t}{b})\iota$ for an injective
$\iota$. The defining term of \rulename{Alias} is over the old scope, and
its parameter is fresh for the new one. An injective renaming preserves the
root of a matrix and the test $t\equiv 0$, hence the mode of
\rulename{Close} and of \rulename{DI}, and it preserves literal equations.
The remaining rules are treated as in \cref{base:lem:struct} and
\cref{disch:lem:param-struct}.
\end{proof}

The lemma fails for $\rulename{Back}_{\nu}$ and for every rule whose
argument is the newest parameter: inserting a newer parameter changes
$\nu$, as in \cref{wit:ex:weakening}.

\begin{proposition}[Equality from \rulename{Hyp}, \rulename{Alias} and
\rulename{Back}]
\label{disch:prop:alias-eq}
Let $\mathcal S$ be as in \cref{disch:lem:alias-struct}. The following
rules are derived in $\mathcal S$, under every context:
\[
 \text{(a)}\ \infrule{}{t=t}
 \qquad
 \text{(b)}\ \infrule{s=t}{t=s}
 \qquad
 \text{(c)}\ \infrule{s=t\qquad H(s)}{H(t)}
 \qquad
 \text{(d)}\ \infrule{s=t\qquad h(s)=r}{r=h(t)}\ \rulename{Tr}
\]
for every term $t$, every formula template $H(v)$, which is to be a literal
equation if $\mathcal S$ has $\rulename{Alias}^{=}$ and
$\rulename{Back}^{=}$, and every term template $h(v)$. The derivations
consist of the given premises, weakened by
\cref{disch:lem:alias-struct}, uses of assumptions, and instances of
\rulename{Alias} and \rulename{Back}. The source of every instance of
\rulename{Back} that they add is the parameter of the innermost
\rulename{Alias} whose branch contains it, and it is the newest parameter
at that node. In (a), (b) and (d), and in (c) when $H$ is a literal
equation, every added instance of \rulename{Alias} and \rulename{Back} has a
literal equation as its conclusion.
\end{proposition}

\begin{proof}
(a) Choose $b$ fresh. Under the alias $b=t$ the use of its assumption
proves $b=t$, and \rulename{Alias} exports $(b=t)\subst{t}{b}$, which is
$t=t$ because $b$ does not occur in $t$.

(b) Let $\Gamma\proves s=t$, and choose $b$ fresh for $\Gamma$, $s$ and
$t$. Work under the alias $b=s$, where the given derivation is weakened.
\rulename{Back} with the first premise $b=s$ and the goal $b=t$ has the
second premise $(b=t)\subst{s}{b}\equiv(s=t)$; so it proves $b=t$. By (a),
$t=t$; the alias of that derivation is opened and discharged before the
next step, so $b$ is again the newest parameter. \rulename{Back} with the
first premise $b=t$ and the goal $t=b$ has the second premise
$(t=b)\subst{t}{b}\equiv(t=t)$; so it proves $t=b$. \rulename{Alias}
discharges $b=s$ and exports $(t=b)\subst{s}{b}\equiv(t=s)$.

(c) Let $\Gamma\proves s=t$ and $\Gamma\proves H(s)$. By (b),
$\Gamma\proves t=s$. Choose $b$ fresh for $\Gamma$, $H$, $s$, $t$ and the
names of the given derivations, and work under the alias $b=t$.
\rulename{Back} with the first premise $b=t$ and the goal $b=s$ has the
second premise $t=s$; so it proves $b=s$. \rulename{Back} with the first
premise $b=s$ and the goal $H(b)$ has the second premise
$H(b)\subst{s}{b}\equiv H(s)$, because $b$ does not occur in $H(v)$; so it
proves $H(b)$. \rulename{Alias} exports $H(b)\subst{t}{b}\equiv H(t)$. The
parameter $b$ marks exactly the distinguished positions of the template;
the occurrences of $s$ and $t$ in its passive part are not touched, the
distinguished variable may occur several times, under the binder of an
iterator or of $\D{}{}$, or not at all, and capture is excluded by the
choice of $b$ and by simultaneous substitution. If $H$ is a literal
equation, the four conclusions $b=s$, $H(b)$, $H(t)$ and those of (b) are
equations.

(d) Apply (c) to $H(v)\equiv(h(v)=r)$, which gives $h(t)=r$, and then (b).
\end{proof}

Transitivity is (c) with $H(v)\equiv(s=v)$, and joining at a common left
side is (d) with the identity template; so every instance of
\rulename{Param} \eqref{disch:eq:param} is derived.

\begin{remark}[The equality fragment]
\label{disch:rem:alias-eq}
The proposition is a completeness statement for an equality fragment alone,
on raw terms: it uses neither a computation equation nor a logical rule,
whereas the derivations of \rulename{Tr} in
\cref{disch:thm:name-transport,disch:thm:vsub-transport,disch:thm:param-congruence}
use packets, computations or joining at the root, and
\cref{disch:lem:alias} uses all three. It does not bear on the conjecture
recorded in \cref{disch:sec:further} that substitution for a variable
alone is weaker on raw terms: \rulename{Alias} is an additional rule.
In the raw calculus with substitution at a variable and no reflexivity,
reflexivity is lost (\cref{cert:thm:vs}); its invariant is not preserved by
\rulename{Alias}, whose premise $b=t\proves b=t$ is valid vacuously while
its conclusion $t=t$ may fail the head test. \pkg{alias\_back} contains
an independent checker of the fragment with \rulename{Hyp},
\rulename{Alias} and \rulename{Back} only, with a schematic unary predicate
for the template of (c); it accepts seven certificates, among them
transport below $\D{}{}$ and iterator binders, nested premise derivations
and a replacement that contains its source, and rejects twelve malformed
ones.
\end{remark}

\subsection{The kernel with aliases}
\label{disch:sec:alias-kernel}

\begin{definition}
\label{disch:def:alias-kernel}
$\Nat_{\mathrm{AB}}$ is $\Nat$ with \rulename{Tr} replaced by
\rulename{Alias} and \rulename{Back}; its rules are \rulename{Hyp},
\rulename{I0}, \rulename{IS}, \rulename{Alias}, \rulename{Back},
\rulename{QE}, \rulename{DE}, \rulename{DI} and \rulename{QI}.
$\Nat^{=}_{\mathrm{AB}}$ has $\rulename{Alias}^{=}$ and
$\rulename{Back}^{=}$ in their place. Contexts are lists of native
formulas.
\end{definition}

\begin{theorem}[The kernel with aliases]
\label{disch:thm:alias-kernel}
For every context $\Gamma$ and every native formula $A$,
\[
 \Gamma\proves_{\Nat}A
 \quad\Longleftrightarrow\quad
 \Gamma\proves_{\Nat_{\mathrm{AB}}}A
 \quad\Longleftrightarrow\quad
 \Gamma\proves_{\Nat^{=}_{\mathrm{AB}}}A,
\]
by effective transformations of derivations that are the identity on the
end sequent. Consequently the ordinary readings of $\Nat_{\mathrm{AB}}$ and
$\Nat^{=}_{\mathrm{AB}}$ are bi-interpretable and equiconsistent with
$\iPRA$, with whole domains, actual equality and both numerical composites
with graph $x=y$.
\end{theorem}

\begin{proof}
Every rule of $\Nat^{=}_{\mathrm{AB}}$ is a rule of $\Nat_{\mathrm{AB}}$.
From $\Nat_{\mathrm{AB}}$ to $\Nat$: by induction on the derivation, with
\cref{disch:lem:alias-sound} at \rulename{Alias} and \rulename{Back}; the
other rules are rules of $\Nat$. The substitution in that lemma is a
substitution in derivations of $\Nat$, where it is admissible; it is not
assumed for the new calculi. From $\Nat$ to $\Nat^{=}_{\mathrm{AB}}$: by
induction on the derivation, copying every rule other than \rulename{Tr}
and replacing every \rulename{Tr} by the derivation of
\cref{disch:prop:alias-eq}(d), whose alias parameters are chosen away from
all names of the translated premises; the premises are weakened into the
alias contexts by \cref{disch:lem:alias-struct}, applied to
$\Nat^{=}_{\mathrm{AB}}$. The translations do not change terms, formulas
or contexts, so the statement on $\iPRA$ follows from
\cref{base:thm:bridge}, and $0=\Suc 0$ is fixed.
\end{proof}

One rule of transport is replaced by two schemes. The theorem is a
different organization of equality, not a deletion of equality reasoning,
and it does not reduce the number of rule families.

\subsection{Seven rules on clause syntax and the newest parameter}
\label{disch:sec:alias-newest}

\begin{definition}[The raw template calculus]
\label{disch:def:alias-tr}
$\mathcal T_{\mathrm r}$ is the calculus over uniform clauses with the
rules \rulename{I0} and \rulename{IS}, as axioms under every context,
\rulename{Tr} \eqref{disch:eq:tr}, \rulename{Use} at an arbitrary term,
without premise for a clause without antecedent, and \rulename{Open} and
\rulename{Close} of \cref{disch:def:param}, the assumption of
\rulename{Close} being $\langle B(t,\eta)\rangle$.
\end{definition}

$\mathcal T_{\mathrm r}$ is the raw template calculus of
\pkg{parameter\_substitution} with assumptions without antecedent; both
packages of this section use it as their matched control.

\begin{lemma}[The raw template calculus]
\label{disch:lem:alias-tr}
\begin{enumerate}[label=\textup{(\alph*)}]
\item For every context $\Gamma$ of uniform clauses and every formula $A$,
 $\Gamma\proves_{\mathcal T_{\mathrm r}}A$ if and only if
 $\Gamma^{\circ}\proves_{\mathcal P}A$, where $\Gamma^{\circ}$ replaces
 every occurrence $[\ ;\,y.\,B]$ by $[H;\,y.\,B]$ as in
 \cref{claus:def:v}.
\item For every context $\Gamma$ of native formulas and every native
 formula $A$, with $\tau$ as in \cref{disch:thm:param},
 $\Gamma\proves_{\Nat}A$ if and only if
 $\langle\tau\Gamma\rangle\proves_{\mathcal T_{\mathrm r}}\tau A$.
\item In $\mathcal T_{\mathrm r}$, under every context of uniform clauses,
 reflexivity, symmetry, transport in every formula template and the case
 introduction for every matrix and every witness are derived, and
 weakening, injective renaming, cut of an assumption and substitution of
 terms for parameters ($\Gamma\proves A$ implies
 $\Gamma\sigma\proves A\sigma$) are admissible.
\end{enumerate}
\end{lemma}

\begin{proof}
(a) The occurrences of $\Gamma$ and $\Gamma^{\circ}$ correspond by
position. From $\mathcal P$ to $\mathcal T_{\mathrm r}$: \rulename{Compute}
is \rulename{I0} or \rulename{IS}; every instance of \rulename{Param} is an
instance of \rulename{Tr}, with the template $u\subst{w}{s}$ if $s$ is a
parameter, the identity template if $s$ is not a parameter and $u\equiv s$,
and the constant template $u$ otherwise (proof of \cref{disch:thm:param});
a \rulename{Use} of an occurrence $[H;\,y.\,B]$ that corresponds to
$[\ ;\,y.\,B]$ becomes the \rulename{Use} without premise, the derivation
of $H$ being discarded; every other \rulename{Use} keeps its translated
premise; \rulename{Open} adds the same clause on both sides, and
\rulename{Close} adds $\operatorname{ord}(B(t,\eta))$, which is
$\langle B(t,\eta)\rangle^{\circ}$. From $\mathcal T_{\mathrm r}$ to
$\mathcal P$: \rulename{I0} and \rulename{IS} are \rulename{Compute};
\rulename{Tr} is derived in $\mathcal P$ under every context
(\cref{disch:cor:param-eq}); a \rulename{Use} of $[\ ;\,y.\,B]$ becomes a
\rulename{Use} of $[H;\,y.\,B]$ whose premise is reflexivity
(\cref{disch:lem:root}); the other rules are copied.

(b) $\langle\tau\Gamma\rangle^{\circ}=\operatorname{ord}(\tau\Gamma)$, so
(a) and \cref{disch:thm:param} give the claim.

(c) A rule derived in $\mathcal P$ under every context, whose premises have
the context of the conclusion extended by assumption clauses, is derived
in $\mathcal T_{\mathrm r}$ by (a), since
$(\Gamma,\langle F\rangle)^{\circ}=\Gamma^{\circ},\operatorname{ord}(F)$;
the same holds for admissible weakening, renaming and cut. Reflexivity and
symmetry are \cref{disch:lem:root}, transport is
\cref{disch:lem:param-closure}(e), the case introduction is
\cref{disch:lem:param-case}, and weakening, renaming and cut are
\cref{disch:lem:param-struct}. Substitution is proved directly by induction
on the derivation, as \cref{claus:prop:subst}, with the eigenparameters made
fresh for $\sigma$: the computation axioms go to computation axioms; for
\rulename{Tr} the template variable is chosen fixed by $\sigma$ and not free
in its values, as in \cref{base:lem:struct}; \rulename{Use} and
\rulename{Open} commute with $\sigma$; a \rulename{Close} keeps its mode
unless it is in case mode with a literal equation as matrix and a witness
$t$ with $t\sigma\equiv 0$, and then the substituted premise is the premise
of the case introduction at the witness $0$, which is derived.
\end{proof}

\begin{definition}[Clause calculi with aliases]
\label{disch:def:alias-clause}
$\mathcal B$ is the calculus over uniform clauses with the rules
\rulename{I0}, \rulename{IS}, \rulename{Use} at an arbitrary term,
\rulename{Back}, \rulename{Alias}, \rulename{Open} and \rulename{Close} (of
$\mathcal T_{\mathrm r}$). $\mathcal B_{\nu}$ has the seven rules
\rulename{I0}, \rulename{IS}, $\rulename{Use}_{\nu}$,
$\rulename{Back}_{\nu}$, \rulename{Alias}, \rulename{Open} and
\rulename{Close}, where
\begin{equation}
 \infrule{\Pi;\Gamma\proves A}{\Pi;\Gamma\proves B(\nu(\Pi))}\
 \rulename{Use}_{\nu},
 \qquad[A;\,y.\,B(y)]\text{ an occurrence in }\Gamma,
 \label{disch:eq:usenu}
\end{equation}
without premise for a clause without antecedent.
\end{definition}

A certificate of $\rulename{Use}_{\nu}$ names the occurrence and contains
no term, and a certificate of $\rulename{Back}_{\nu}$ contains no source
and no template. $\mathcal B_{\nu}$ has no primitive reflexivity, symmetry,
transitivity, joining at the root, transport, template, normalization,
predecessor or logical connective. \rulename{I0} and \rulename{IS} are two
equations even when a checker groups them. The witness of \rulename{Close}
is still supplied, and its progress mode is restricted to literal
equations, so no induction for compound formulas is added. A use of an
assumption $\langle A\rangle$ by $\rulename{Use}_{\nu}$ gives $A$ whatever
$\nu(\Pi)$ is.

\begin{theorem}[The newest parameter suffices]
\label{disch:thm:alias-compiler}
Let $p$ be a derivation of $\Pi;\Gamma\proves A$ in $\mathcal B$. Let
$\theta$ be an injective map of $\Pi$ into a scope $\Pi'$, in any order
and with any further parameters, and let $\Delta$ be a context over $\Pi'$
together with an injective map of the occurrences of $\Gamma$ to
occurrences of $\Delta$ that sends an occurrence of $C$ to one of $\theta
C$. Then a derivation of $\Pi';\Delta\proves\theta A$ in $\mathcal B_{\nu}$
is constructed effectively. If every \rulename{Back} of $p$ has a literal
equation as its conclusion, so has every $\rulename{Back}_{\nu}$ of the
constructed derivation.
\end{theorem}

In the terms of the chapter, every derivation of $\mathcal B$ yields a
persistent proof (\cref{disch:def:persistent}) in $\mathcal B_{\nu}$, that
is, hereditary derivability in the sense of \cref{wit:def:hereditary},
with clause occurrences in the place of members of the context.

\begin{proof}
By induction on $p$, simultaneously for all $\Pi'$, $\theta$, $\Delta$ and
occurrence maps. The induction hypothesis for a premise of the last rule is
used with data constructed at that node, and no derivation produced by the
construction is weakened, renamed or substituted into.

\emph{\rulename{I0}, \rulename{IS}.} The image under $\theta$ is an
instance of the same axiom.

\emph{\rulename{Open}.} The major premise is translated at
$(\Pi',\theta,\Delta)$. For the branch, over $\Pi,b$ with the opened clause
$\mathcal C_B(b)$, choose $b'\notin\Pi'$, extend $\theta$ by $b\mapsto b'$
and the occurrence map by the opened clause to the new opened clause
$\mathcal C_{\theta B}(b')$, and translate the branch at
$(\Pi',b';\ \Delta,\mathcal C_{\theta B}(b'))$. \rulename{Open} concludes
$\theta A$, which is over $\Pi'$.

\emph{\rulename{Close}.} Choose $z'\notin\Pi'$, extend $\theta$ by
$z\mapsto z'$, append $\langle\theta B(\theta t,\eta')\rangle$, where
$\eta'$ is the image of $\eta$, and translate the branch there.
\rulename{Close} at the witness $\theta t$ concludes $\theta(\D{x}{y}B)$.
The mode is the same, since an injective map of parameters preserves the
root of the matrix, and $\theta t\equiv 0$ exactly if $t\equiv 0$.

\emph{\rulename{Alias}} with the defining term $t$ and the parameter $b$.
Choose $b'\notin\Pi'$, extend $\theta$ by $b\mapsto b'$ to $\theta'$,
append $\langle b'=\theta t\rangle$, and translate the branch there; it
proves $\theta'F$. \rulename{Alias} exports
$(\theta'F)\subst{\theta t}{b'}\equiv\theta(F\subst{t}{b})$.

\emph{\rulename{Back}} with the premises $v=t$ and $F\subst{t}{v}$ and the
conclusion $F$, where $v$ is a parameter of $\Pi$. Choose $b\notin\Pi'$
and work under the alias $b=\theta v$, in the scope $\Pi',b$ and the
context $\Delta,\langle b=\theta v\rangle$, where $b$ is the newest
parameter. Translate both premises there, by the induction hypothesis with
the same $\theta$ and the occurrence map into the extended context; they
give $\theta v=\theta t$ and $\theta(F\subst{t}{v})$. Then:
\begin{itemize}
\item $\rulename{Back}_{\nu}$ with the first premise $b=\theta v$, the use
 of the alias assumption, and the goal $b=\theta t$ has the second premise
 $\theta v=\theta t$; it proves $b=\theta t$.
\item Let $G\defeq(\theta F)\subst{b}{\theta v}$. Since $b$ does not occur
 in $\theta F$, $G\subst{\theta t}{b}\equiv(\theta F)\subst{\theta
 t}{\theta v}\equiv\theta(F\subst{t}{v})$, also when $t$ contains $v$.
 $\rulename{Back}_{\nu}$ with the first premise $b=\theta t$ and the goal
 $G$ proves $G$.
\item \rulename{Alias} exports $G\subst{\theta v}{b}\equiv\theta F$.
\end{itemize}
If $F$ is a literal equation, so are $b=\theta t$ and $G$.

\emph{\rulename{Use}} of an occurrence $[A';\,y.\,B]$ at a term $t$. Choose
$b\notin\Pi'$ and work under the alias $b=\theta t$. If the clause has an
antecedent, translate the derivation of $A'$ there.
$\rulename{Use}_{\nu}$ of the image of the occurrence at the newest
parameter $b$ proves $(\theta B)(b)$, and \rulename{Alias} exports
$(\theta B)(\theta t)$, which is $\theta(B(t))$. No \rulename{Back} is
used.

Every call of the induction hypothesis concerns a proper subderivation of
$p$, so the construction terminates, although contexts and terms grow.
\end{proof}

\begin{corollary}[Equivalence and bridge]
\label{disch:cor:alias-newest}
For every scope, every context $\Gamma$ of uniform clauses and every
formula $A$,
\[
 \Gamma\proves_{\mathcal B_{\nu}}A
 \quad\Longleftrightarrow\quad
 \Gamma\proves_{\mathcal B}A
 \quad\Longleftrightarrow\quad
 \Gamma\proves_{\mathcal T_{\mathrm r}}A ,
\]
and for every native context $\Gamma$ and native formula $A$,
$\Gamma\proves_{\Nat}A$ if and only if
$\langle\tau\Gamma\rangle\proves_{\mathcal B_{\nu}}\tau A$. Weakening,
renaming, substitution and cut are admissible in $\mathcal B_{\nu}$ as
operations on derivable judgments. The ordinary reading of
$\mathcal B_{\nu}$ is bi-interpretable and equiconsistent with $\iPRA$,
with whole domains, actual equality and both numerical composites with
graph $x=y$.
\end{corollary}

\begin{proof}
$\mathcal B_{\nu}\subseteq\mathcal B$: $\rulename{Use}_{\nu}$ is
\rulename{Use} at the term $\nu(\Pi)$, and $\rulename{Back}_{\nu}$ is an
instance of \rulename{Back}.

$\mathcal B\subseteq\mathcal T_{\mathrm r}$: by induction on the
derivation. \rulename{I0}, \rulename{IS}, \rulename{Use}, \rulename{Open}
and \rulename{Close} are rules of $\mathcal T_{\mathrm r}$. For
\rulename{Back}, transport along the symmetric $t=b$ in the template
$F\subst{w}{b}$, $w$ fresh, sends $F\subst{t}{b}$ to $F$
(\cref{disch:lem:alias-tr}(c)). For \rulename{Alias}, the substitution
$[t/b]$ in the translated branch gives
$\Gamma,\langle t=t\rangle\proves F\subst{t}{b}$, since $b$ is not free in
$\Gamma$ or $t$, and cut with reflexivity gives
$\Gamma\proves F\subst{t}{b}$ (\cref{disch:lem:alias-tr}(c)).
% ed.: the package proves this inclusion through the ordinary reading of the reference theory; the direct proof in T_r covers contexts of arbitrary uniform clauses, not only contexts of assumptions and opened clauses.

$\mathcal T_{\mathrm r}\subseteq\mathcal B$: by induction on the
derivation; \rulename{Tr} is \cref{disch:prop:alias-eq}(d), with the
premises weakened by \cref{disch:lem:alias-struct}, and the other rules are
copied.

$\mathcal B\subseteq\mathcal B_{\nu}$: \cref{disch:thm:alias-compiler} at
the identity map and the identity context.

The statement on $\Nat$ is \cref{disch:lem:alias-tr}(b). The structural
operations are transferred through $\mathcal T_{\mathrm r}$; they are
consequences of the equivalence, not steps of its proof. The bridge is
\cref{base:thm:bridge} composed with these translations, which leave terms
unchanged and fix $0=\Suc 0$; both formula round trips are those of
\cref{disch:thm:param}.
\end{proof}

\begin{corollary}[Backward substitution for equations only]
\label{disch:cor:alias-eqback}
\Cref{disch:cor:alias-newest} holds with $\rulename{Back}^{=}_{\nu}$ in the
place of $\rulename{Back}_{\nu}$ in $\mathcal B_{\nu}$, \rulename{Alias}
still exporting arbitrary formulas.
\end{corollary}

\begin{proof}
The translation from $\mathcal T_{\mathrm r}$ to $\mathcal B$ adds only the
instances of \rulename{Back} of \cref{disch:prop:alias-eq}(d), which have
literal equations as conclusions, and weakening preserves this. The
compiler of \cref{disch:thm:alias-compiler} then produces only instances of
$\rulename{Back}^{=}_{\nu}$. Its \rulename{Use} case exports compound
formulas by \rulename{Alias}.
\end{proof}

\begin{remark}[The kernel with the newest parameter]
\label{disch:rem:alias-kernel-newest}
The same induction applies to the kernel. Let $\Nat^{\nu}_{\mathrm{AB}}$
be $\Nat_{\mathrm{AB}}$ with $\rulename{Back}_{\nu}$ in the place of
\rulename{Back} and with \rulename{QE} restricted to the target $\nu(\Pi)$,
the witness of \rulename{DI} being arbitrary. By induction on a derivation
of $\Nat_{\mathrm{AB}}$, for all injective $\theta$ and all contexts
$\Delta\supseteq\theta\Gamma$, $\Nat^{\nu}_{\mathrm{AB}}$ derives
$\Pi';\Delta\proves\theta A$: \rulename{Hyp} is kept; \rulename{Back} is
compiled as above; \rulename{QE} at $t$ is compiled as \rulename{Use},
under the alias $b=\theta t$ with the translated major premise and base,
\rulename{QE} at $b$ and \rulename{Alias}; and \rulename{DE}, \rulename{DI}
and \rulename{QI} append a fresh eigenparameter as \rulename{Open} and
\rulename{Close} do. So $\Nat^{\nu}_{\mathrm{AB}}$ and $\Nat$ derive the
same sequents. If only \rulename{QE} is restricted and \rulename{Back} is
not, weakening in $\Nat_{\mathrm{AB}}$ (\cref{disch:lem:alias-struct})
preserves the height of a derivation, and an induction on the height
suffices: the source premises of \rulename{QE} are weakened into the alias
context in $\Nat_{\mathrm{AB}}$ and then translated, and no translated
derivation is weakened.
% ed.: the package states this lemma with ``regenerate or weaken the two premise proofs''; the weakening is admissible only for the unrestricted source derivations, before translation, and the induction is then on their height.
In the Lean development the target of \rulename{QE} can be restricted to a
variable or the numeral $1$ (\cref{disch:rem:spec-lean}), and the numeral
cannot be dropped there (\Lean{vqprf\_var\_incomplete}). With
\rulename{Alias} it can: the countermodel of
\cref{disch:rem:spec-lean} does not validate \rulename{Alias}. In the
context $\Q{y}(\sel(y,0,1)=0)$, under the alias $b=1$, \rulename{QE} at $b$
with the base $\sel(0,0,1)=0$, an instance of \rulename{I0}, proves
$\sel(b,0,1)=0$; \rulename{Alias} exports $\sel(1,0,1)=0$, and
\rulename{IS} and the derived equality rules give $0=1$, which is false in
the countermodel while the context is true there.
\end{remark}

\begin{remark}[Relation to the other restrictions of the chapter]
\label{disch:rem:alias-relations}
In $\mathcal V_0$ (\cref{disch:sec:newest}) the source of the equality rule
is the newest variable too, but completeness rests on full normal forms,
which make the rows of the packet literal, and on internally sourced
derivations (\cref{disch:def:internal}). In the derivations produced by
\cref{disch:thm:alias-compiler} every source of $\rulename{Back}_{\nu}$ is
the parameter of an \rulename{Alias}, so the alias plays the role of the
packet, raw terms suffice, and the invariant is replaced by the induction
over all target data. In $\mathcal U_*$ (\cref{disch:sec:spec}) the use of
a clause takes no term either; it uses the newest clause at $1$ and every
other clause at the newest parameter, on zero-normal terms, and needs the
naming formula $\mathrm{Nm}_t$ and an induction on rank. With an
\rulename{Alias} that exports compound formulas, $\rulename{Use}_{\nu}$
needs neither; \cref{disch:sec:alias-atomic} shows that the rank induction
returns when \rulename{Alias} exports equations only. The selection
$\nu(\Pi)$ is that of \cref{wit:def:scope}, and the families of
derivations of the proof are those of \cref{wit:def:hereditary}
(\cref{wit:rem:families}). In \cref{wit:sec:newest} the witness of the
introduction is the newest parameter and the target of \rulename{QE} is
free; here the target of specialization and the source of equality are the
newest parameter and the witness of \rulename{Close} is supplied. The
restriction of the target does not fall under
\cref{wit:thm:finite-range}, where the witness and the target are
restricted together and no rule supplies a term; here the witness is
arbitrary and \rulename{Alias} supplies an arbitrary term.
\end{remark}

\subsection{Equation-only aliases and the rank of formulas}
\label{disch:sec:alias-atomic}

\begin{definition}[The calculus $\mathcal A$]
\label{disch:def:alias-atomic}
$\mathcal A$ has the seven rules \rulename{I0}, \rulename{IS},
$\rulename{Use}_{\nu}$, $\rulename{Back}^{=}_{\nu}$,
$\rulename{Alias}^{=}$, \rulename{Open} and \rulename{Close} over uniform
clauses. The conclusion of \rulename{Open} is an arbitrary formula.
\end{definition}

So $\mathcal A$ is $\mathcal B_{\nu}$ with both new rules restricted to
literal equations. No rule of $\mathcal A$ exports a compound formula from
an alias, and no rule takes a term argument except the witness of
\rulename{Close}. The rank of a formula is the number of its binders
(\cref{disch:sec:spec}); substitution of terms does not change it.

Situations $(\Pi;\Gamma)$, their extensions and persistent proofs are
those of \cref{disch:def:persistent}, with the parameter list written $\Pi$
and $\mathcal A$ in the place of $\mathcal U_*$; a persistent proof is the
hereditary family of \cref{wit:def:hereditary}, and \pkg{equation\_aliases}
calls it a \emph{uniform proof procedure}. A \emph{local derivation} over
$(\Pi;\Gamma)$ is one derivation of $\Pi;\Gamma\proves A$ in $\mathcal A$.
Both $\rulename{Use}_{\nu}$ and $\rulename{Back}^{=}_{\nu}$ change their
meaning when a newer parameter is added, so a local derivation is not in
general a persistent proof.

\begin{lemma}[Elementary persistence]
\label{disch:lem:alias-persist}
\begin{enumerate}[label=\textup{(\alph*)}]
\item A persistent proof over a situation restricts to a persistent proof
 over each of its extensions.
\item If $\langle A\rangle$ occurs in $\Gamma$, then $A$ has a persistent
 proof over $(\Pi;\Gamma)$.
\item The instances of \rulename{I0} and \rulename{IS} over $\Pi$ have
 persistent proofs.
\end{enumerate}
\end{lemma}

\begin{proof}
(a) Extensions compose. (b) In every extension apply
$\rulename{Use}_{\nu}$ to the image of the occurrence; its body does not
contain the bound variable, so the conclusion is the image of $A$ whatever
the newest parameter is. (c) The image of an instance is an instance.
\end{proof}

\begin{lemma}[The equational base]
\label{disch:lem:alias-atomic-base}
Let $(\Pi;\Gamma)$ be a situation and $s,t$ terms over $\Pi$.
\begin{enumerate}[label=\textup{(\roman*)}]
\item $t=t$ has a persistent proof.
\item If $s=t$ has a persistent proof, so has $t=s$.
\item $(\mathrm T_0)$ If $E(v)$ is a literal equation and $s=t$ and $E(s)$
 have persistent proofs, so has $E(t)$. In particular, if $s=t$ and
 $h(s)=r$ have persistent proofs, so has $r=h(t)$.
\item $(\mathrm S_0)$ If $[A;\,y.\,E(y)]$ occurs in $\Gamma$ with $E$ a
 literal equation, and $A$ has a persistent proof (no condition for a
 clause without antecedent), then $E(t)$ has a persistent proof.
\end{enumerate}
\end{lemma}

\begin{proof}
Let an extension be given, and write $s,t,E,A$ for the renamed expressions.
In each item the derivation is that of \cref{disch:prop:alias-eq},
constructed in the extension, with every premise taken from its persistent
proof at the situation in which it is needed; this situation is an
extension of the given one, by (a) of \cref{disch:lem:alias-persist}. No
derivation is weakened. Every \rulename{Alias} and \rulename{Back} of the
constructions has a literal equation as its conclusion, and every
\rulename{Back} uses the alias opened last, which is the newest parameter
at its node; in (ii) the alias of the reflexivity is discharged before the
second \rulename{Back}. The second statement of (iii) is (iii) with
$E(v)\equiv(h(v)=r)$ followed by (ii). For (iv) choose $b$ fresh, open the
alias $b=t$, take the antecedent from its persistent proof at the extended
situation, apply $\rulename{Use}_{\nu}$ to the image of the occurrence at
the newest parameter $b$, which proves $E(b)$, and export $E(t)$ by
$\rulename{Alias}^{=}$.
\end{proof}

The antecedent $A$ in (iv) may have any rank; its persistent proof is
given, not constructed by an induction on the rank of $E$.

\begin{lemma}[Names]
\label{disch:lem:alias-name}
For a term $t$ over $\Pi$ let
\[
 \mathrm{Name}(t)\defeq\D{x}{y}\bigl(\sel(y,t,x)=t\bigr),\qquad
 \mathcal N_{b,t}\defeq[\sel(0,t,b)=t;\ y.\ \sel(y,t,b)=t],
\]
the bound variable of the selector not being free in $t$; the matrix is
that of \cref{disch:lem:alias} with $a\equiv t$.
\begin{enumerate}[label=\textup{(\alph*)}]
\item $\mathrm{Name}(t)$ has a persistent proof over every situation
 $(\Pi;\Gamma)$, and opening it by \rulename{Open} at a fresh $b$ adds the
 clause $\mathcal N_{b,t}$.
\item If $\mathcal N_{b,t}$ occurs in $\Gamma$, at any position, and $b$ is
 a parameter of $\Pi$ not occurring in $t$, then $b=t$ has a persistent
 proof over $(\Pi;\Gamma)$.
\end{enumerate}
\end{lemma}

\begin{proof}
(a) In every extension apply \rulename{Close} at the witness $t$. Its goal
$\sel(\Suc z,t,t)=t$ is the instance of \rulename{IS} with the constant
body $t$, and its assumption is not used; so the derivation is valid in
either mode, also when $t\equiv 0$ and the rule is in progress mode. The
opened clause at $b$ is $\mathcal N_{b,t}$.

(b) The antecedent $\sel(0,t,b)=t$ is an instance of \rulename{I0}, with a
persistent proof by \cref{disch:lem:alias-persist}(c). By
\cref{disch:lem:alias-atomic-base}(iv) at the term $1$,
$\sel(1,t,b)=t$ has a persistent proof, and \rulename{IS} gives
$\sel(1,t,b)=b$ persistently. Transport (iii) along $\sel(1,t,b)=t$ in the
equation template $v=b$ sends $\sel(1,t,b)=b$ to $t=b$, and (ii) gives
$b=t$. The clause $\mathcal N_{b,t}$ need not be the newest clause and $b$
need not be the newest parameter, because the specialization of (iv) opens
its own alias.
\end{proof}

In contrast with \cref{disch:lem:spec-name}, where $b=t$ is derived only
while the naming clause is newest, the equality here is persistent from the
start. This is what allows it to be used after later binders have added
newer parameters.

\begin{theorem}[Transport and specialization at every rank]
\label{disch:thm:alias-rank}
For every $r\geq 0$ and every situation $(\Pi;\Gamma)$ the following hold
in $\mathcal A$.
\begin{enumerate}[label=\textup{(\alph*)}]
\item[$(\mathrm T_r)$] Let $F(w)$ be a formula template of rank at most
 $r$. If $s=t$ and $F(s)$ have persistent proofs over $(\Pi;\Gamma)$, so
 has $F(t)$.
\item[$(\mathrm S_r)$] Let $[A;\,y.\,F(y)]$ occur in $\Gamma$ with $F$ of
 rank at most $r$, let $A$ have a persistent proof over $(\Pi;\Gamma)$
 (no condition without antecedent), and let $t$ be a term over $\Pi$. Then
 $F(t)$ has a persistent proof over $(\Pi;\Gamma)$.
\item[$(\mathrm L_r)$] For $r>0$, let $F(w)$ have rank exactly $r$. From a
 local derivation of $\Pi;\Gamma\proves F(s)$ and a persistent proof of
 $s=t$ over $(\Pi;\Gamma)$, a local derivation of
 $\Pi;\Gamma\proves F(t)$ is constructed.
\end{enumerate}
\end{theorem}

\begin{proof}
$(\mathrm T_0)$ and $(\mathrm S_0)$ are (iii) and (iv) of
\cref{disch:lem:alias-atomic-base}. No local form $(\mathrm L_0)$ is used.
Let $r>0$ and suppose $(\mathrm T_{r-1})$ and $(\mathrm S_{r-1})$.

\emph{$(\mathrm L_r)$.} Write $F(w)=\D{x}{y}B(w,x,y)$, with $B$ of rank
$r-1$. Use the given local derivation of $F(s)$ as the major premise of
\rulename{Open} at a fresh $c$; the branch has the situation
$\Theta_1=(\Pi,c;\ \Gamma,\mathcal C)$ with
$\mathcal C\defeq[B(s,c,0);\,y.\,B(s,c,y)]$. The local derivation is used
in the situation in which it was given and is never moved under $c$ or a
later parameter. In $\Theta_1$ introduce $F(t)$ by \rulename{Close} at the
witness $c$. A parameter is not the literal $0$, so the rule is in case
mode whatever $B$ is; its premise has the situation
$\Theta_2=(\Pi,c,z;\ \Gamma,\mathcal C,\langle B(t,c,0)\rangle)$ and the
goal $B(t,c,\Suc z)$. Over $\Theta_2$, which is an extension of
$(\Pi;\Gamma)$:
\begin{itemize}
\item $B(t,c,0)$ has a persistent proof by
 \cref{disch:lem:alias-persist}(b), and $t=s$ has one by restriction and
 \cref{disch:lem:alias-atomic-base}(ii); $(\mathrm T_{r-1})$ for the
 template $B(w,c,0)$ gives a persistent proof of $B(s,c,0)$;
\item $(\mathrm S_{r-1})$ for the occurrence of $\mathcal C$, with this
 antecedent, at the term $\Suc z$ gives a persistent proof of
 $B(s,c,\Suc z)$;
\item $(\mathrm T_{r-1})$ for the template $B(w,c,\Suc z)$ and the
 restricted proof of $s=t$ gives a persistent proof of $B(t,c,\Suc z)$.
\end{itemize}
Its member at the identity extension is the premise of \rulename{Close},
which proves $F(t)$ over $\Theta_1$. The formula $F(t)$ does not contain
$c$, and \rulename{Open} concludes $F(t)$ over $(\Pi;\Gamma)$. Even for
$r=1$ every appeal to transport meets the hypotheses of $(\mathrm T_0)$,
since all its premises are persistent.

\emph{$(\mathrm T_r)$.} Templates of rank below $r$ are covered by the
hypothesis. For rank $r$, given an extension, take the member of the
persistent proof of $F(s)$ there, which is a local derivation of the image
of $F(s)$, restrict the persistent proof of $s=t$ to the extension, and
apply $(\mathrm L_r)$ there.

\emph{$(\mathrm S_r)$.} Formulas of rank below $r$ are covered by the
hypothesis. For rank $r$, in a given extension, with renamed expressions,
derive $\mathrm{Name}(t)$ (\cref{disch:lem:alias-name}(a)) and open it at
a fresh $b$. In the situation $\Theta$ of the branch, $b=t$ has a
persistent proof by \cref{disch:lem:alias-name}(b). Take $A$ from its
persistent proof at $\Theta$. $\rulename{Use}_{\nu}$ of the image of the
occurrence at the newest parameter $b$ gives a local derivation of $F(b)$
over $\Theta$, and $(\mathrm L_r)$ with the persistent equation $b=t$ gives
a local derivation of $F(t)$ over $\Theta$. The formula $F(t)$ does not
contain $b$, and \rulename{Open} discharges it.

The recursion is well founded: $(\mathrm L_r)$ uses $(\mathrm T_{r-1})$
and $(\mathrm S_{r-1})$, $(\mathrm T_r)$ uses $(\mathrm L_r)$, and
$(\mathrm S_r)$ uses $(\mathrm L_r)$ and the name, which uses rank $0$
only.
\end{proof}

\begin{theorem}[Equation-only aliases]
\label{disch:thm:alias-atomic}
Every derivation of $\Pi;\Gamma\proves A$ in $\mathcal T_{\mathrm r}$ yields
a persistent proof of $A$ over $(\Pi;\Gamma)$ in $\mathcal A$, effectively;
in particular a derivation in $\mathcal A$. Consequently $\mathcal A$,
$\mathcal B_{\nu}$, $\mathcal B$ and $\mathcal T_{\mathrm r}$ derive the
same judgments over every context of uniform clauses,
$\Gamma\proves_{\Nat}A$ if and only if
$\langle\tau\Gamma\rangle\proves_{\mathcal A}\tau A$ for native $\Gamma$
and $A$, weakening, renaming, substitution and cut are admissible in
$\mathcal A$ as operations on derivable judgments, and the ordinary reading
of $\mathcal A$ is bi-interpretable and equiconsistent with $\iPRA$ with
whole domains, actual equality and both numerical composites with graph
$x=y$.
\end{theorem}

\begin{proof}
By induction on the derivation in $\mathcal T_{\mathrm r}$.
\begin{itemize}
\item \rulename{I0}, \rulename{IS}: \cref{disch:lem:alias-persist}(c).
\item \rulename{Tr}: \cref{disch:lem:alias-atomic-base}(iii), applied to
 the persistent proofs of the two premises.
\item \rulename{Use} of $[A';\,y.\,B]$ at $t$: $(\mathrm S_r)$ of
 \cref{disch:thm:alias-rank} with $r$ the rank of $B$, applied to the
 persistent proof of the antecedent given by the induction hypothesis.
 The antecedent may have a larger rank than $B$; its persistent proof comes
 from the induction on the derivation, not from the induction on rank.
\item \rulename{Open} with the major premise $\D{x}{y}B$: in a given
 extension take the member of the persistent proof of the major premise,
 choose a fresh $b'$, extend the extension by $b\mapsto b'$ and by the
 opened clause to the new opened clause, take the member of the persistent
 proof of the branch at this extension, and apply \rulename{Open}.
\item \rulename{Close}: in the same way, with the fresh $z'$ and the
 renamed assumption; the mode is preserved by injective renaming.
\end{itemize}
The calls of the induction hypothesis concern proper subderivations, and
the operations of \cref{disch:thm:alias-rank} are constructed beforehand by
the induction on rank, so the two recursions do not interact. No
derivation of $\mathcal A$ is weakened, renamed or cut. The member at the
identity extension is a derivation of $\Pi;\Gamma\proves A$ in
$\mathcal A$.

Every rule of $\mathcal A$ is an instance of a rule of $\mathcal B_{\nu}$,
so $\mathcal A\subseteq\mathcal B_{\nu}$, and the remaining statements are
\cref{disch:cor:alias-newest}.
\end{proof}

\begin{remark}[Comparison and cost]
\label{disch:rem:alias-atomic}
\Cref{disch:thm:alias-rank} has the pattern of \cref{disch:lem:spec-rank}.
There the transport at rank $0$ is \rulename{Step}, which opens no scope,
so $(\mathrm T_0)$ has a local form. Here every operation at rank $0$ opens
an alias and is available only as a persistent proof; a local derivation of
$F(b)$ obtained by $\rulename{Use}_{\nu}$ while $b$ is newest is consumed
as the major premise of \rulename{Open} before any newer parameter is
introduced, as the local derivation at the case variable is in
\cref{wit:sec:newest}. The constructions $(\mathrm L_r)$ and
$(\mathrm S_r)$ use \rulename{Open} with a compound conclusion; whether
\rulename{Open} can be restricted as well is open
(\cref{disch:sec:further}). The theorem answers the question of
\pkg{alias\_back} whether \rulename{Alias} can be restricted to equations
while every clause is used at the newest parameter only. The compiled
certificates of the two examples of \pkg{equation\_aliases} with compound
specialization at rank $1$ and $2$ have $1{,}308$ and $5{,}302$ bytes; no
bound on the size of the translation is claimed.
\end{remark}

\section{A constancy rule in place of induction}
\label{disch:sec:const}

This section draws on an appendix of \pkg{substitutive\_discharge}. It is
independent of the discharge rule: it presents an arithmetic whose one
arithmetical inference has no induction hypothesis and no step function.
The price is a primitive Boolean equality.

\subsection{The theory}

\begin{definition}[The constancy theory]
\label{disch:def:const}
The logic is intuitionistic first-order logic with equality over one sort.
Terms are
\[
 t::=x\mid 0\mid\Suc(t)\mid\delta(t,t)\mid R(n,a;\,i,u.\,h),
\]
where the recursor binds $i$ and $u$ in $h$ only. With $1\defeq\Suc 0$ and
$s\neq t\defeq(s=t\to\bot)$, the axioms are the universal closures of
\begin{align}
 R(0,a;\,i,u.\,h)&=a, \label{disch:eq:r0}\\
 R(\Suc n,a;\,i,u.\,h)&=h[n/i,\ R(n,a;\,i,u.\,h)/u], \label{disch:eq:rs}\\
 0=1&\to\bot, \label{disch:eq:sep}\\
 \bigl(s=t\land\delta(s,t)=0\bigr)&\lor
 \bigl(s\neq t\land\delta(s,t)=1\bigr), \label{disch:eq:split}
\end{align}
the substitution in \eqref{disch:eq:rs} being simultaneous, and the one
further rule is
\begin{equation}
 \infrule{\Gamma\proves f(\Suc z)=f(z)}{\Gamma\proves f(t)=f(0)}\
 \rulename{Const}
 \label{disch:eq:construle}
\end{equation}
for a term template $f(x)$, with $z$ fresh for $\Gamma$, $t$ and $f$. The
context $\Gamma$ consists of arbitrary first-order formulas. There is no
coverage axiom, no injectivity axiom for successor, no induction schema
and no assumed reflection of equality.
\end{definition}

\subsection{Recovery of quantifier-free induction}

Define
\[
 \sel(v,a,b)\defeq R(v,a;\,i,u.\,b),\qquad
 a+b\defeq R(b,a;\,i,u.\,\Suc u),\qquad
 \mathit{pd}(v)\defeq R(v,0;\,i,u.\,i).
\]
These are the counter-aware selector, addition and predecessor of
\cref{ctr:lem:decoder,ctr:thm:compiler}, for a recursor whose value at the
counter $0$ is its initializer and not the successor of it: no outer
decoder is needed, and each equation below is one computation step.

\begin{lemma}[Elementary computations]
\label{disch:lem:const-elem}
Without \rulename{Const}: $\sel(0,a,b)=a$ and $\sel(\Suc v,a,b)=b$;
$a+0=a$, $a+\Suc b=\Suc(a+b)$ and $a+1=\Suc a$; $\mathit{pd}(0)=0$ and
$\mathit{pd}(\Suc v)=v$; successor is injective; $\Suc v\neq 0$; and
\begin{equation}
 \delta(s,t)=0\lor\delta(s,t)=1,\qquad
 \delta(s,t)=0\leftrightarrow s=t .
 \label{disch:eq:deltabool}
\end{equation}
\end{lemma}

\begin{proof}
The equations are instances of \eqref{disch:eq:r0} and
\eqref{disch:eq:rs}, and injectivity follows from
$\mathit{pd}(\Suc v)=v$. If $\Suc v=0$, applying $w\mapsto\sel(w,0,1)$
gives $1=0$, against \eqref{disch:eq:sep}; no distinction of cases on a
number is used. The disjunction in \eqref{disch:eq:deltabool} is read off
\eqref{disch:eq:split}. If $s=t$, the second disjunct of
\eqref{disch:eq:split} contradicts it, so $\delta(s,t)=0$. If
$\delta(s,t)=0$, the second disjunct would give $0=1$, so $s=t$.
\end{proof}

For a quantifier-free formula $A$ define a term $q_A$, with the value $0$
meaning that $A$ holds:
\begin{align*}
 q_{s=t}&\defeq\delta(s,t), & q_\bot&\defeq 1,\\
 q_{A\land B}&\defeq\sel(q_A,q_B,1), &
 q_{A\lor B}&\defeq\sel(q_A,0,q_B),\\
 q_{A\to B}&\defeq\sel(q_A,q_B,0).
\end{align*}

\begin{lemma}[Boolean terms]
\label{disch:lem:const-bool}
For every quantifier-free $A$, without \rulename{Const},
\[
 q_A=0\lor q_A=1,\qquad q_A=0\leftrightarrow A .
\]
\end{lemma}

\begin{proof}
By induction on $A$ in the metatheory, with the two computations of
$\sel$ and \eqref{disch:eq:deltabool}. In the case of an implication
distinguish the two values of $q_A$: if it is $0$, then $A$ holds and
$A\to B$ is equivalent to $B$; if it is $1$, then $\neg A$ follows and the
implication holds. Only the decidability already proved for the
quantifier-free antecedent is used.
\end{proof}

\begin{theorem}[Quantifier-free induction from constancy]
\label{disch:thm:const-ind}
The theory of \cref{disch:def:const} derives
\[
 \infrule{\Gamma\proves A(0)\qquad\Gamma,\ A(z)\proves A(\Suc z)}
         {\Gamma\proves A(t)}
\]
for every quantifier-free $A$ and every context $\Gamma$ of first-order
formulas, $z$ fresh. In particular it derives induction for literal
equations under an arbitrary context.
\end{theorem}

\begin{proof}
Write $q(n)\defeq q_{A(n)}$ and define the terms
\[
 e(n)\defeq\sel\bigl(q(n),0,\sel(q(\Suc n),1,0)\bigr),\qquad
 a(n)\defeq R(n,0;\,i,u.\,u+e(i)),\qquad
 F(n)\defeq a(n)+q(n).
\]
These are terms of the language: $e$ and the addition are expanded before
the outer recursor is formed. The computations give $a(0)=0$ and
$a(\Suc n)=a(n)+e(n)$. By \cref{disch:lem:const-bool} the pair
$(q(n),q(\Suc n))$ has four possible values. Three of them give
$F(\Suc n)=F(n)$ with no induction hypothesis:
\[
 \begin{array}{c|c|c|c}
 q(n)&q(\Suc n)&e(n)&F(\Suc n)=F(n)\\ \hline
 0&0&0&a(n)=a(n)\\
 1&1&0&\Suc a(n)=\Suc a(n)\\
 1&0&1&\Suc a(n)=\Suc a(n)
 \end{array}
\]
The fourth, $(0,1)$, is impossible under $\Gamma$: $q(n)=0$ gives $A(n)$,
the step premise gives $A(\Suc n)$, and \cref{disch:lem:const-bool} gives
$q(\Suc n)=0$, against $0\neq 1$. By elimination of the finitely many
disjunctions, $\Gamma\proves F(\Suc z)=F(z)$. \rulename{Const} gives
$F(t)=F(0)$. The base $A(0)$ gives $q(0)=0$, so $F(0)=0$. If $q(t)=1$,
then $0=F(t)=a(t)+1=\Suc a(t)$, against $\Suc v\neq 0$. Hence $q(t)=0$,
and $A(t)$ holds.
\end{proof}

The only identities of addition used are $a+0=a$ and $a+1=\Suc a$;
commutativity and associativity are not used, so no induction is hidden in
an algebraic rearrangement.

\subsection{Both interpretations}

\begin{theorem}[The constancy theory and $\iPRA$]
\label{disch:thm:const-ipra}
The theory of \cref{disch:def:const} is bi-interpretable with $\iPRA$,
with whole domains, actual equality and both numerical composites provably
the identity. In particular the two are equiconsistent.
\end{theorem}

\begin{proof}
\emph{$\iPRA$ in the constancy theory.} \Cref{disch:thm:const-ind} applied
to the quantifier-free formula $n=0\lor n=\Suc(\mathit{pd}(n))$ gives
coverage: the case $0$ is immediate, and at a successor
$\mathit{pd}(\Suc n)=n$ gives the second disjunct without the hypothesis.
Injectivity of successor is in \cref{disch:lem:const-elem}. Every
primitive recursive definition compiles into the recursor with a counter,
as in \cref{ctr:thm:compiler} with $R$ in place of $\operatorname{Rec}$:
zero, successor, projections and composition are term operations, and
$f(0,\vec x)=g(\vec x)$,
$f(\Suc n,\vec x)=h(n,f(n,\vec x),\vec x)$ is represented by
$R\bigl(n,g(\vec x);\,i,u.\,h(i,u,\vec x)\bigr)$ after $g$ and $h$ are
compiled; the defining equations are instances of \eqref{disch:eq:r0} and
\eqref{disch:eq:rs}. With separation, injectivity and
\cref{disch:thm:const-ind} this interprets $\iPRA$.

\emph{The constancy theory in $\iPRA$.} Each fixed recursor term denotes a
primitive recursive function, and its two equations hold by construction.
\rulename{Const} is admissible: under its premise, induction for the
equation $f(n)=f(0)$ gives the conclusion. For $\delta$ an equality
characteristic is constructed as follows. Let
\[
 d(a,b)\defeq R(b,a;\,i,u.\,\mathit{pd}(u)).
\]
The computation equations and induction for equations give
\[
 d(a,0)=a,\quad d(a,\Suc b)=\mathit{pd}(d(a,b)),\quad
 d(\Suc a,\Suc b)=d(a,b),\quad d(\Suc n,n)=1 ;
\]
the third by induction on $b$ at fixed $a$, with the base
$\mathit{pd}(\Suc a)=a$ and the step obtained by applying $\mathit{pd}$ to
the preceding equation, the fourth by induction on $n$ from the third.
Inductions on the second argument of addition give $0+b=b$ and
$\Suc a+b=\Suc(a+b)$. Quantifier-free induction on $b$, with $a$ as a
parameter, proves
\begin{equation}
 d(a,b)\neq 0\ \to\ a=b+d(a,b).
 \label{disch:eq:rec}
\end{equation}
The base uses $0+a=a$. For the step put $r\defeq d(a,b)$ and suppose
$\mathit{pd}(r)=d(a,\Suc b)\neq 0$. Coverage for $r$ gives
$r=\Suc(\mathit{pd}(r))$, the case $r=0$ contradicting
$\mathit{pd}(r)\neq 0$. So $r\neq 0$, and the induction hypothesis gives
$a=b+r=b+\Suc(\mathit{pd}(r))=\Suc b+\mathit{pd}(r)$. Now define
\[
 Z(a)\defeq\sel(a,1,0),\qquad
 U(a)\defeq\sel\bigl(a,0,Z(\mathit{pd}(a))\bigr),\qquad
 E(a,b)\defeq\sel\bigl(b,Z(a),U(d(a,\mathit{pd}(b)))\bigr).
\]
Coverage and separation give that $Z(a)$ and $U(a)$ are $0$ or $1$, that
$Z(a)=1$ exactly if $a=0$, and that $U(a)=1$ exactly if $a=1$. If $b=0$,
then $E(a,b)=1$ is equivalent to $a=0$. If $b=\Suc n$, then $E(a,b)=1$ is
equivalent to $d(a,n)=1$, which implies $a=n+1=\Suc n=b$ by
\eqref{disch:eq:rec}, and conversely $a=\Suc n$ gives $d(\Suc n,n)=1$. So
$E(a,b)$ is $0$ or $1$ and equals $1$ exactly if $a=b$, and
\[
 \delta_*(a,b)\defeq Z(E(a,b))
\]
is $0$ or $1$ and equals $0$ exactly if $a=b$. It satisfies
\eqref{disch:eq:split}. This derivation uses quantifier-free inductions of
$\iPRA$; it is not part of the bootstrap of
\cref{disch:thm:const-ind}. It is a variant of the reflection argument
recorded in \cref{ctr:rem:reflection}, with the control $d$ and the test
$E$ in place of the sign and the reflection term used there.

\emph{Composites.} In the constancy theory, induction for equations
(\cref{disch:thm:const-ind}) proves the uniqueness of every fixed
primitive recursive definition under arbitrary parameters; by induction on
terms the round trip of every recursor term has the value of the original.
The same argument in $\iPRA$ treats the round trip of each primitive
recursive function symbol. For the primitive $\delta$, both $\delta$ and
the translated $\delta_*$ satisfy \eqref{disch:eq:split} in the constancy
theory, and the distinction of the two cases of the first shows
$\delta(a,b)=\delta_*(a,b)$: both are $0$ if $a=b$ and both are $1$
otherwise. Congruence and induction on formulas give both formula round
trips. The numerical isomorphisms have graph $x=y$.
\end{proof}

\begin{remark}[What the rule does and does not simplify]
\label{disch:rem:const}
The arithmetical inference concerns one term and one unconditional
recurrence with the identity as step; it has no discharged hypothesis and
no step template, and the accumulator $a(n)$ is a construction of the
proof and not of the rule. Against the kernel the presentation adds a
binary term constructor, a primitive equality split in place of coverage,
and a second bound variable of the recursor. The split
\eqref{disch:eq:split} is essential to the proof above: replacing it by
the assertion that an internally defined equality test is faithful,
without proving that assertion in the induction-free base, would be
circular. Likewise, for a tail-recursive iterator the equation
$d(a,\Suc b)=d(\mathit{pd}(a),b)$ does not by itself give
$d(a,\Suc b)=\mathit{pd}(d(a,b))$; the missing commutation is a fusion
property. No checker is supplied for this theory, and nothing is claimed
about its size.
\end{remark}

\section{Checkers and program equivalences}
\label{disch:sec:impl}

This section summarizes the programs of \pkg{witness\_discharge},
\pkg{inferred\_witness}, \pkg{deterministic\_opening},
\pkg{first\_instance}, \pkg{clause\_discharge},
\pkg{substitutive\_discharge}, \pkg{variable\_substitution},
\pkg{parameter\_substitution}, \pkg{deterministic\_specialization},
\pkg{alias\_back} and \pkg{equation\_aliases}. Sizes are in
\cref{disch:sec:sizes}.

\emph{Common contract.} Every checker is a standalone Python module that
reads one JSON value, verifies a closed theorem, prints \texttt{ACCEPT}
and exits with status zero on success, and rejects on any exception; no
decision depends on an \texttt{assert}. The checkers of
\pkg{variable\_substitution} and \pkg{deterministic\_specialization} code
syntax by tagged lists. The others use integers and bare lists: $0$ is
zero, a negative integer is a free parameter level, a positive integer is
a de Bruijn index, \texttt{[0,t]} is a successor, \texttt{[n,a,h]} an
iterator or recursor binding one variable in its last entry,
\texttt{[s,t]} an equation and \texttt{[B]} a paired formula in whose
matrix index $1$ is $y$ and index $2$ is $x$; the expected sort
distinguishes a successor from an equation. Subterms are validated before
a contraction can erase them. Inserted terms have no loose positive index,
so insertion below a binder needs no lifting. The theorems of this chapter
concern open sequents; the public interfaces check closed goals.

\emph{\pkg{witness\_discharge}.} A goal-directed checker of
$\mathcal W^{\mathcal R}$ with three packets, for \rulename{Use} (clause
index, term, premise), \rulename{Intro} (witness, branch) and
\rulename{Name} (matrix, local formula, target, major premise, equality
premise, branch). The exported formula is compared with the already
validated outer goal, which enforces the absence of the eigenparameter
also below binders and after normalization has erased an occurrence.

\emph{\pkg{inferred\_witness}.} A forward checker of
$\mathcal O^{\mathcal R}$ that infers every conclusion once. It stores the
antecedent $B(b,0)$ of the opened clause as computed and substitutes the
selected witness in the body only: if the test succeeds the antecedent is
$H$ and is unchanged by the substitution, and if it fails the substitution
is the identity, so the stored clause is $\mathcal C_B(\rho_B(b))$ in both
cases.

\emph{\pkg{deterministic\_opening}.} The checker of
\pkg{witness\_discharge} with the packet of the third rule shortened to
matrix, major premise and branch. On validated syntax the test of the row
at $0$ is that both entries are zero, and the test that the first entry of
the row at $1$ is the integer of $b$ also establishes that the row is an
equation, a paired formula having a list as its only entry.

\emph{\pkg{first\_instance}.} A checker of $\mathcal X^{\mathcal R}$ in
which every proof node is a triple; the depth of the clause context equals
the number of parameters in scope.

\emph{\pkg{clause\_discharge}.} Two goal-directed checkers on the input
\texttt{[goal, proof]}, one of $\mathcal W^{\mathcal R}$ and one of
$\mathcal X^{\mathcal R}$. The next fresh parameter is computed from the
length of the context, and one traversal performs shifting and opening.
The replacement acts on a parameter leaf exactly if it is the left side of
the row and is the fresh witness. A variant that infers conclusions was
tried and was not smaller.

\emph{\pkg{substitutive\_discharge}.} A checker of $\mathcal K$ with four
instructions and no fuel (\cref{disch:rem:sd}), its matched control with
separate rules, and the raw companion with its control.

\emph{\pkg{variable\_substitution}.} A goal-directed native checker of the
six-rule form with the newest variable, obtained from the checker of
\pkg{full\_normalization} by changing the equality packet, and a forward
checker of $\mathcal V_0$ whose equality packet consists of two premise
proofs. In the native checker the leaves for hypothesis and reflexivity
ignore an unused hint, as in its control.

\emph{\pkg{parameter\_substitution}.} A forward checker of $\mathcal P$
with five instructions. The certificate of $0=0$ joins two copies of the
computation $\It{0}{0}{u}{u}=0$. Substitution in a proper subterm acts on
negative integers only, so it cannot rewrite the zero that marks a
successor.

\emph{\pkg{deterministic\_specialization}.} A goal-directed checker of
$\mathcal U_*$ with one instruction for each of the five rules. The
instruction for a clause contains the index and the proof of the
antecedent; the target is computed from the index and the scope, and
contexts are stored newest first. A translator in both directions is
supplied and is not part of the checker.

\emph{\pkg{alias\_back}.} Two checkers of $\mathcal B_{\nu}$
on tagged trees: the input is an object with a closed goal and a proof, and
the proof has one tag for each of the computations (grouped), for
$\rulename{Use}_{\nu}$ (occurrence index and premise), for
$\rulename{Back}_{\nu}$, \rulename{Alias} (defining term and branch),
\rulename{Open} and \rulename{Close}. The public goal and every active term, formula, clause index and
scope are validated, and nothing is normalized. \rulename{Alias} validates
its defining term in the outer scope before it appends its parameter and
the clause $\langle b=t\rangle$. A use of a clause without antecedent
ignores its premise field, as in the control. The shorter checker computes
the scope from the context by \cref{disch:prop:scope}, which applies since
\rulename{Alias}, \rulename{Open} and \rulename{Close} each add one clause
and one parameter. The matched controls check $\mathcal T_{\mathrm r}$ with
the same coding and the same simplifications. Terms and formulas are coded by integers and bare
lists as in the common contract. An untrusted translator
compiles proofs of $\mathcal T_{\mathrm r}$ into $\mathcal B$ by
\cref{disch:prop:alias-eq} and then into $\mathcal B_{\nu}$ by
\cref{disch:thm:alias-compiler}; its results for $32$ inherited closed
certificates, among them compound binders, specialization, induction
hypotheses, capture avoidance and zero sources, are replayed by the
checkers, and one open example uses a parameter map that does not preserve
the order and a non-contiguous occurrence map. The checker of the equality
fragment of \cref{disch:rem:alias-eq} is not a complete checker and has no
size record.

\emph{\pkg{equation\_aliases}.} The two checkers of $\mathcal B_{\nu}$ of
\pkg{alias\_back} with the additional tests that the conclusions of
\rulename{Alias} and $\rulename{Back}_{\nu}$ are literal equations. The
readable checker keeps the scope separate from the context, for open
examples; the shorter one computes it from the context. An untrusted
compiler implements \cref{disch:thm:alias-atomic}; it compiles $34$ closed
derivations, the $32$ inherited ones and two with compound specialization
at rank $1$ and $2$, and three open examples with non-monotone parameter
maps, non-contiguous occurrence maps, duplicated clauses and an antecedent
of rank $2$ against bodies of rank $0$, $1$ and $2$. The matched controls
are the checkers of \pkg{alias\_back}, which accept the certificates with
compound \rulename{Alias} and \rulename{Back} that the new checkers
reject.

Three statements about programs, independent of the rules, are proved in
the packages.

\begin{proposition}[Scope from the context]
\label{disch:prop:scope}
In a checker whose public entry starts from the empty context and scope
$0$, and in which the branch of every scope-extending rule adds exactly
one clause and one parameter while no other call changes either, every
reachable call satisfies $\mathit{scope}=|\mathit{context}|$. The scope
argument can therefore be computed from the context.
\end{proposition}

\begin{proof}
By induction on the tree of calls: the initial call satisfies the
equation, the rules without a branch recur with both unchanged, and a
branch increments both.
\end{proof}

\pkg{deterministic\_specialization} applies this to the checker of
$\mathcal U_*$ and to its control, and \pkg{clause\_discharge} and
\pkg{first\_instance} build it into their checkers. The statement concerns
the public entry; it is not claimed for a direct call of an internal
function with a context and a scope that do not match.

\emph{Fusion of the rewrite with the continuation.}
\pkg{deterministic\_opening} reorders, in three inherited compact checkers
of an earlier guarded calculus, the rewriting of the goal by an
elimination and the construction of the new hypothesis: the hypothesis is
constructed first and the rewrite is an argument of the continuation call,
evaluated before the scope is decremented by the left-to-right evaluation
order of arguments. The two orders accept the same certificates. A rewrite
changes the goal only if the major premise is a nonempty list, and then
the construction of the hypothesis cannot perform the assignment of an
introduction; conversely an introduction changes the goal only if the
major premise is empty, and then the rewrite is the identity; so the two
operations commute. The proof of the major premise runs in its own
activation and does not affect the scope of the caller, and the rewrite
consumes no proof packet. By induction on the number of packets each new
program accepts exactly the certificates of its control and consumes the
same suffix. These programs do not implement \rulename{Open}.

\emph{One traversal for grammar and scope.} The supplement of
\pkg{deterministic\_specialization} modifies the validated compact checker
of \pkg{scalar\_uniqueness}, which checks the unchanged kernel. The
control traverses the public goal twice, once for sorts and arities and
once for scope; the new program does both in the first traversal, under a
flag that is set for the public entry only and cleared before the first
proof packet is read. By induction on the goal the fused traversal
succeeds exactly if both old traversals succeed and then returns the goal;
with the flag cleared every added guard is skipped and the remaining
computation is that of the control.

\section{Checker sizes}
\label{disch:sec:sizes}

This section lists the sizes reported by the eleven packages named at the
beginning of \cref{disch:sec:impl}. All figures follow the protocol of
\cref{base:sec:measure} and are those stated by the packages. A control is
matched if it shares the term and
formula coding, the normalizer, the context representation and the public
contract with the program and differs in the rule under comparison; where
a package compares with an unmatched program this is said.

\begin{longtable}{@{}p{0.34\textwidth}rrrp{0.40\textwidth}@{}}
\caption{Complete checkers of the packages of this chapter.}
\label{disch:tab:sizes}\\
\toprule
Program & Tokens & AST & Bytes & Control (tokens / AST / bytes)\\
\midrule
\endfirsthead
\toprule
Program & Tokens & AST & Bytes & Control (tokens / AST / bytes)\\
\midrule
\endhead
\bottomrule
\endfoot
\multicolumn{5}{@{}l}{\pkg{witness\_discharge}}\\
$\mathcal W^{\mathcal R}$ & 732 & 813 & 2,273 &
 separate elimination and transport: 769 / 841 / 2,362\\
$\mathcal W^{\mathcal R}$ with fuel & 777 & 868 & 2,486 &
 the same with fuel: 814 / 896 / 2,575\\
\midrule
\multicolumn{5}{@{}l}{\pkg{inferred\_witness}}\\
$\mathcal O^{\mathcal R}$, forward & 683 & 771 & 1,990 &
 explicit equation at opening: 676 / 765 / 1,998; explicit equation at
 output: 663 / 751 / 1,983\\
$\mathcal O^{\mathcal R}$ with fuel & 728 & 826 & 2,203 &
 unmatched: \pkg{witness\_discharge} with fuel, 777 / 868 / 2,486\\
\midrule
\multicolumn{5}{@{}l}{\pkg{deterministic\_opening}}\\
$\mathcal O^{\mathcal R}$, goal-directed & 730 & 809 & 2,253 &
 checker of \pkg{witness\_discharge}: 732 / 813 / 2,273\\
intermediate export variant & 745 & 828 & 2,296 &
 the same control\\
compact profile, fused (inherited calculus) & 267 & 302 & 508 &
 265 / 306 / 516\\
AST profile, fused (inherited calculus) & 263 & 276 & 580 &
 260 / 278 / 602\\
validated native profile, fused (inherited calculus) & 360 & 409 & 780 &
 358 / 413 / 788\\
\midrule
\multicolumn{5}{@{}l}{\pkg{first\_instance}}\\
$\mathcal X^{\mathcal R}$, triples & 704 & 804 & 2,157 &
 $\mathcal W^{\mathcal R}$, triples: 701 / 813 / 2,228\\
$\mathcal X^{\mathcal R}$ with fuel & 752 & 863 & 2,340 &
 $\mathcal W^{\mathcal R}$ with fuel: 749 / 872 / 2,411\\
\midrule
\multicolumn{5}{@{}l}{\pkg{clause\_discharge}}\\
$\mathcal W^{\mathcal R}$ (explicit discharge) & 636 & 708 & 2,010 &
 checker of \pkg{witness\_discharge}: 732 / 813 / 2,273\\
$\mathcal X^{\mathcal R}$ (deterministic discharge) & 653 & 729 & 2,058 &
 explicit discharge: 636 / 708 / 2,010\\
explicit discharge with fuel & 673 & 755 & 2,139 & \\
deterministic discharge with fuel & 690 & 776 & 2,187 &
 explicit discharge with fuel: 673 / 755 / 2,139\\
\midrule
\multicolumn{5}{@{}l}{\pkg{substitutive\_discharge}}\\
$\mathcal K$ & 713 & 798 & 2,267 &
 $\mathcal K_{\mathrm{sep}}$: 726 / 809 / 2,318\\
raw companion & 744 & 828 & 2,550 &
 separate rules: 756 / 837 / 2,606\\
\midrule
\multicolumn{5}{@{}l}{\pkg{variable\_substitution}}\\
native six rules, newest variable & 899 & 951 & 2,499 &
 checker of \pkg{full\_normalization} with \rulename{Tr}:
 908 / 965 / 2,529\\
$\mathcal V_0$ & 898 & 968 & 2,826 &
 clausal \rulename{Tr}: 907 / 982 / 2,856\\
\midrule
\multicolumn{5}{@{}l}{\pkg{parameter\_substitution}}\\
$\mathcal P$ & 769 & 892 & 2,857 &
 template rule, unused code removed: 765 / 881 / 2,907; template rule,
 shared opener: 777 / 895 / 2,937\\
\midrule
\multicolumn{5}{@{}l}{\pkg{deterministic\_specialization}}\\
$\mathcal U_*$, scope from context & 819 & 896 & 2,914 &
 $\mathcal U$, scope from context: 819 / 895 / 2,926\\
$\mathcal U_*$, explicit scope & 829 & 903 & 2,938 &
 $\mathcal U$, explicit scope (checker of \pkg{uniform\_step}, newest
 first): 829 / 902 / 2,950\\
kernel, fused grammar and scope traversal & 360 & 412 & 820 &
 checker of \pkg{scalar\_uniqueness}: 381 / 439 / 866\\
\midrule
\multicolumn{5}{@{}l}{\pkg{alias\_back}}\\
$\mathcal B_{\nu}$, readable & 825 & 966 & 3,149 &
 $\mathcal T_{\mathrm r}$, readable: 764 / 883 / 2,918\\
$\mathcal B_{\nu}$, scope from context, shorter interface & 773 & 916 &
 2,903 & $\mathcal T_{\mathrm r}$, the same simplifications:
 716 / 838 / 2,683\\
\midrule
\multicolumn{5}{@{}l}{\pkg{equation\_aliases}}\\
$\mathcal A$, readable & 843 & 989 & 3,215 &
 $\mathcal B_{\nu}$, readable (checker of \pkg{alias\_back}):
 825 / 966 / 3,149\\
$\mathcal A$, scope from context & 791 & 939 & 2,969 &
 $\mathcal B_{\nu}$, scope from context: 773 / 916 / 2,903\\
\end{longtable}

The comparisons are local. In its matched pair, \rulename{Name} against
separate rules is smaller in all three coordinates;
\rulename{Open} against an explicit equation at the opening saves $8$
bytes and costs $7$ tokens and $6$ nodes, and is $7$ bytes larger than the
control with an explicit equation at the output; \rulename{Export} saves
$71$ bytes and $9$ nodes and costs $3$ tokens in \pkg{first\_instance},
and costs $48$ bytes, $17$ tokens and $21$ nodes in
\pkg{clause\_discharge}; \rulename{SDis} is smaller in all three;
\rulename{VSub} for the newest variable saves $30$ bytes, $9$ tokens and
$14$ nodes in both profiles; \rulename{Param} saves $50$ bytes and costs
$4$ tokens and $11$ nodes; $\rulename{Use}_*$ saves $12$ bytes, leaves the
tokens unchanged and costs one node, and the computation of the scope
from the context saves a further $24$ bytes, $10$ tokens and $7$ nodes.
\rulename{Alias} with the two restrictions to the newest parameter costs,
against the template control, $61$ tokens, $83$ nodes and $231$ bytes in
the readable pair and $57$ tokens, $78$ nodes and $220$ bytes in the
shorter pair, and the restriction of both rules to equations costs a
further $18$ tokens, $23$ nodes and $66$ bytes in each pair; the interface
loses its supplied template and specialization term and gains a rule.
The programs of \pkg{alias\_back} and \pkg{equation\_aliases} read their
input as tagged trees with a fully validated closed goal.
\pkg{parameter\_substitution} also supplies, as a comparison source and
not as a new calculus, a refined checker of the two-case evaluation
calculus with $741$ tokens, $862$ nodes and $2{,}731$ bytes. Every package
states that shorter checkers of the repository with other certificate
contracts, in particular the compact checker of $516$ bytes and, where it
is named, the validated checker of $788$ bytes, remain smaller than its
programs, and none claims a minimum. For \pkg{alias\_back} and
\pkg{equation\_aliases} the compact programs named are those of the
coordinatewise envelope of $248$ tokens, $270$ nodes and $487$ bytes,
respectively $479$ bytes, attained by different programs under a promised
input alphabet.

\section{Further questions}
\label{disch:sec:further}

This section collects the questions, limits and unproved claims stated by
the eleven packages named at the beginning of \cref{disch:sec:impl}, and one
conjecture recorded in the repository.

\begin{question}[Substitution for a variable on raw terms]
On full normal forms substitution for a variable is the only equality rule
needed (\cref{disch:thm:vsub}). On raw terms the rows of a packet are not
literal, and \rulename{Param}, which also joins two equations at the root,
suffices (\cref{disch:thm:param}). Is substitution for a variable alone,
without joining at the root, strictly weaker on raw terms? The repository
records the affirmative answer as a conjecture. With the additional rule
\rulename{Alias}, substitution for a parameter in the backward direction
suffices (\cref{disch:cor:alias-newest}); \pkg{alias\_back} states that
this does not settle the conjecture, which concerns substitution for a
variable without such a rule.
\end{question}

\begin{question}[The least equality interface on raw terms]
The instances of \rulename{Param} that the completeness proof uses are
joining at the root and substitution at a fresh variable; the third case
of $\psi$ is a reversal. Is there a smaller total operation that expresses
exactly the needed instances without a larger test on syntax? An answer
has to account for arbitrary equations between closed terms as
assumptions, and not only for the computation equations.
\end{question}

\begin{question}[Unary congruence on full normal forms]
\Cref{disch:prop:unary} covers raw and zero-normal terms. Full
normalization of the successor computation identifies values that its
interpretation only relates. Whether unary congruence suffices on full
normal forms is not decided.
\end{question}

\begin{question}[The source of the selected equation]
In \rulename{Open} and \rulename{Export} the source of the selected
equation is the witness opened by the same inference. A generalization in
which the row at $1$ may name an older parameter was considered in
\pkg{clause\_discharge} and discarded; the derivation of ordinary
elimination by one \rulename{Export} fails for it. Whether a sound and
complete rule of that kind exists is open.
\end{question}

\begin{question}[The normalization that naming needs]
The packets use the two computations of $\sel$ and the two closed values
$P(1)$ and $P(2)$. Is there a smaller discipline of normalization,
coherent with substitution and with nested binders, that makes exactly
these computations literal and keeps the theorems of
\cref{disch:sec:name,disch:sec:inferred,disch:sec:vsub}? Can a syntax
expose the rows at $0$ and $1$ of a matrix without full normalization and
keep both formula round trips?
\end{question}

\begin{question}[A term template in the substitutive discharge]
Can the formula template of \rulename{SDis} be restricted to a term
template, with transport in compound formulas recovered, while the
complete checker stays smaller than that of $\mathcal K_{\mathrm{sep}}$?
The occurrences of the witness that the rule permits have to be tracked.
\end{question}

\begin{question}[An introduction that is stable under substitution]
The mode of the single introduction depends on the literal spelling of
the witness, and this is the reason why substitution does not preserve the
rules of the calculi of this chapter. Is there a test that does not depend
on the supplied witness, permits the recovery of the case introduction,
and does not cost more in the checker than it saves?
\end{question}

\begin{question}[The equality split]
Can the axiom \eqref{disch:eq:split} of \cref{disch:def:const} be derived
from a small recursion base and \rulename{Const} alone? The split must not
be assumed during such a bootstrap.
\end{question}

\begin{question}[Size of the translations]
\Cref{disch:thm:spec} produces a finite derivation with no polynomial
bound (\cref{disch:rem:spec-cost}). The translations of
\cref{disch:thm:vsub-transport,disch:thm:param-congruence} add one packet
or one alias per transport. The compiler of \cref{disch:thm:alias-atomic}
reconstructs persistent premises in every new scope and nests the
constructions of lower ranks (\cref{disch:rem:alias-atomic}). What are the
least expansions, for certificates as trees and with explicit sharing?
Shared proof graphs need a validation of their references whose cost has to
be measured as well.
\end{question}

\begin{question}[Restriction of opening]
\label{disch:q:alias-open}
In $\mathcal A$ the conclusion of \rulename{Open} is an arbitrary formula,
and the constructions $(\mathrm L_r)$ and $(\mathrm S_r)$ of
\cref{disch:thm:alias-rank} use compound conclusions essentially. Can
\rulename{Open} be restricted to literal equations as conclusions, with a
different logical encoding, while \rulename{Alias} and \rulename{Back} are
restricted to equations and every clause is used at the newest parameter?
(\pkg{equation\_aliases}.)
\end{question}

\begin{question}[The witness of \rulename{Close} in the alias calculi]
\label{disch:q:alias-witness}
In $\mathcal B_{\nu}$ and $\mathcal A$ the witness of \rulename{Close} is
an arbitrary supplied term. Can it be restricted to the newest parameter,
as in \cref{wit:sec:newest}, with a short generalized compiler?
\rulename{Alias} supplies the names, but the assumption of the branch has
to be transported through the new eigenparameter without a weakening of
restricted derivations. (\pkg{alias\_back}.)
\end{question}

\begin{question}[The equality rule of $\mathcal A$]
\label{disch:q:alias-crossed}
Can $\rulename{Back}^{=}_{\nu}$ be replaced by a crossed replacement of one
term, keeping the goal-directed form without a template, the newest source
and the restriction of \rulename{Alias} to equations, so that the checker
becomes smaller? An answer has to account for raw computations and for the
checks of scope. (\pkg{equation\_aliases}.)
\end{question}

\begin{question}[A measured benefit of the aliases]
\label{disch:q:alias-size}
The matched checkers of $\mathcal B_{\nu}$ and $\mathcal A$ are larger than
their controls in every coordinate (\cref{disch:sec:sizes}). Can
\rulename{Alias} and \rulename{Open}, both of which append a parameter and
a clause, share an implementation without hiding their different freshness
and export conditions, or can the equation tests of $\mathcal A$ share a
representation with $\rulename{Use}_{\nu}$, so that the restricted
interface lowers the size of a complete checker? (\pkg{alias\_back},
\pkg{equation\_aliases}.)
\end{question}

\begin{question}[Lower bounds]
A lower bound on the size of a checker needs a fixed public grammar, a
fixed contract for malformed input, a fixed source language and a fixed
measure. No package of this chapter fixes these, and none proves a lower
bound. Is there one complete program that improves the token, node and
byte coordinates of the repository simultaneously under one contract? Is
there a presentation that removes a genuine arithmetic or logical principle
and also has a smaller measured checker (\pkg{alias\_back})? Can
the rules of this chapter be represented in the smallest certificate
languages based on arities without losing their savings to normalization
and to the bookkeeping of scope?
\end{question}

The following are stated by the packages as not achieved or not claimed.
\begin{itemize}
\item \emph{Formalization.} No package contains a Lean or Rocq proof of
 its theorems. The repository proves the kernel forms and the clause forms
 named in
 \cref{disch:rem:name-lean,disch:rem:sd,disch:rem:param,disch:rem:spec-lean}
 and after \cref{disch:thm:inferred,disch:thm:newest}. Not formalized are:
 the four families of $\mathcal K$ on the canonical recursor syntax, the
 five families of $\mathcal P$ together, the positional calculus
 $\mathcal U_*$, the constancy theory of \cref{disch:sec:const}, and the
 calculi $\Nat_{\mathrm{AB}}$, $\mathcal B_{\nu}$ and $\mathcal A$ with the
 compilers of \cref{disch:thm:alias-compiler,disch:thm:alias-rank,disch:thm:alias-atomic}.
 \pkg{alias\_back} and \pkg{equation\_aliases} ask for one formalization of
 the compiler with arbitrary parameter embeddings and occurrence maps,
 reusable for restrictions to the newest parameter, and of the distinction
 between persistent proofs and local derivations; their finite tests
 replay compiled certificates and are not a verification of the
 metatheorems.
\item \emph{Effectiveness.} The proof translations of this chapter are
 effective; none is claimed to be primitive recursive. No primitive
 recursive bound on normalization, and no bound on the fuel needed by a
 derivable formula, is asserted (\cref{disch:rem:fuel}).
 \pkg{clause\_discharge} does not claim an arithmetized equivalence of
 consistency statements in a specified weak metatheory.
\item \emph{Executable translations.} The tests of
 \pkg{deterministic\_opening} supply, for each tested theorem, a matching
 proof in $\mathcal W^{\mathcal R}$; they do not contain a general
 executable compiler from $\mathcal O^{\mathcal R}$, nor certified
 compilers between the calculi on normal forms and the raw kernel with
 normalization traces and the repair of the introduction mode.
\item \emph{The equality fragment.} The proofs of
 \cref{disch:thm:vsub-transport,disch:thm:name-transport} use the logical
 rules; no completeness statement is made for the equality rules alone.
 The fragment of \rulename{Hyp}, \rulename{Alias} and \rulename{Back} is
 the exception: it derives reflexivity, symmetry and transport by itself
 (\cref{disch:prop:alias-eq}).
\item \emph{Size.} No package claims a smallest presentation or a smallest
 checker, and the matched comparisons of \cref{disch:sec:sizes} are not
 improvements in every coordinate. The checkers of \pkg{alias\_back} and
 \pkg{equation\_aliases} are larger than their matched controls in every
 coordinate; the packages claim a smaller equality and specialization
 interface with a proof of completeness, not a smaller checker.
\item \emph{Arithmetic.} The restrictions of this chapter change the
 equality interface and the use of clauses. They do not weaken the
 arithmetic: the construction of the predecessor and the reflection of
 equality (\cref{base:thm:pred,base:thm:reflect}) remain what a proposal
 with weaker induction would have to redo, as \pkg{parameter\_substitution}
 notes. The calculi of \cref{disch:sec:alias} keep the supplied witness of
 \rulename{Close} and its progress mode for equations, so they keep the
 induction of the theory.
\end{itemize}
